% Standalone fragment: requires amsmath, amssymb, amsthm, hyperref.
% Uses only standard theorem/proof environments. No custom mathematical macros.
\section{Complexity consequences and effective gates}
\label{sec:nn-complexity}
E and A are designated input theorems. Every use of M is conditional.
The preceding CX catalogue is background; its ETH refutations and Koiran rows
are not counted again. A proof marked below is a proof of the stated deduction
from those inputs, not an independent verification of the full input theorem.

\subsection{NN-C001 --- Quasilinear sparse systems compress on their original variables}
\begin{proposition}[NN-C001]
Fix integers $k,a,b\geq1$. Given $m\leq n\lceil\log_2 n\rceil^a$ polynomials $f_e\in\mathbb F_2[x_1,\ldots,x_n]$, each listed with at most $\lceil\log_2 n\rceil^b$ monomials of degree at most $k$, one can deterministically construct at most $6n$ linear combinations of the $f_e$ having precisely the same common zeros in $\mathbb F_2^n$. Each combination has polylogarithmically many monomials, of degree at most $k$. Construction and exact counting take
\[
 \operatorname{poly}_{k,a,b}(n+m)\,
 2^{O_{k,a,b}(n\log\log n/\sqrt{\log n})}
\]
bit operations and polynomial space. The construction introduces no auxiliary Boolean variables. This is preservation of the Boolean zero set, not equality of ideals in the unrestricted polynomial ring.
\end{proposition}
\paragraph{Exact inputs.}
E: The Exponential Time Hypothesis is False, supplied LaTeX. Section 6, Lemmas 6.1--6.2; Appendix D Lemma D.2 (lem:tensor-sparse-query).

\begin{proof}
Handle $m=0$ and bounded $n$ directly. Put $B=\max(1,\lceil m/n\rceil)$, $\mu=\lfloor\log_2m\rfloor+1$, $P=20000B^2\mu^2$ and $D_0=320000B^2\mu^3$. Apply E, Lemma 6.1, to construct $A^{(0)}$ with at most $3n$ rows and row weight at most $D_0$, such that every nonzero vector in its kernel has weight at least $m/P$. This lemma concerns arbitrary vectors of bits; it never uses the clause origin of those bits.

For $x\in\mathbb F_2^n$, set $v(x)=(f_1(x),\ldots,f_m(x))$. Write $S=\{x:v(x)=0\}$ and let $D(A)$ be the number of $x\notin S$ with $Av(x)=0$. The count returned by a common-zero query is $N(A)=|S|+D(A)$. At a stage, enumerate all $P$-tuples of pairs $(e,c)\in[m]\times\mathbb F_2$, interpreted as a row $h$ with repeated indices combined modulo two. Every false survivor has at least $m/P$ nonzero coordinates. A uniform tuple hits such a coordinate with probability at least $1-(1-1/P)^P>1/2$. Conditional on one hit, the independent coefficient of that occurrence makes $h\cdot v(x)$ a uniform bit, even if other occurrences repeat the same coordinate. Thus its rejection probability is at least $1/4$.

Consequently $\mathbb E_hD(A;h)\leq(3/4)D(A)$. Evaluate $N(A;h)$ for every tuple and append a minimizer. Since $|S|$ is the same for all candidates, this also minimizes $D$. After $3n$ rounds the false survivor count is at most $2^n(3/4)^{3n}=(27/32)^n<1$, and hence is zero. Every original zero survives every row, proving exact equality of the Boolean zero sets.

A tested system has at most $6n$ equations, degree at most $k$, and at most $\lceil\log_2n\rceil^b\max(P,D_0)$ monomials per equation. This is the interface of Lemma D.2 for one fixed constant depending only on $k,a,b$. There are at most $3n(2m)^P+1=\exp((\log n)^{O_{a,b}(1)})$ counter calls. Initial code construction has the same kind of overhead. Every fixed polylogarithm is $o(n\log\log n/\sqrt{\log n})$, so multiplying these costs gives the stated bound. Enumerate candidates sequentially, retaining only the current matrix, candidate, best candidate and counts; Lemma D.2 then preserves polynomial space.
\end{proof}
\paragraph{Status and complexity impact.}
Proved from designated input E; self-check recorded. Independent complete deduction-proof review by root: passed; imported input theorems are not claimed independently proved. Extends the linear-number-of-equations core to a quasilinear number of sparse equations while keeping the original witness parameter $n$. It supplies no polynomial-time hitting set: selection itself uses the expensive exact counter and depends on the particular map $v$. It does not preserve all zeros over extension fields and does not turn a signed circuit into a monotone circuit.

\paragraph{Geometric and representation dictionary.}
In the Boolean coordinate algebra $B_n=\mathbb F_2[x]/(x_i^2-x_i)$, evaluation gives $B_n\cong\prod_{x\in\mathbb F_2^n}\mathbb F_2$. An ideal in this product is determined by the coordinates on which it vanishes. Hence the original and compressed equations do generate the same ideal after passage to $B_n$, and define the same finite reduced, finite etale zero-dimensional scheme. The resulting isomorphism and its count are computed without listing its possibly $2^n$ components. The trace dictionary is SGA7 II XXII Corollary 3.2; the relative finite trace is SGA4 XVII 6.2.3. There is no new reductive-group action or representation multiplicity attached by this compression.

\subsection{NN-C002 --- Fixed characteristic-two fields and exact preimage tasks}
\begin{proposition}[NN-C002]
For every fixed $r,k,a,b$, NN-C001 extends to sparse systems over $\mathbb F_{2^r}$ in $n$ field variables, with $m\leq n(\log n)^a$ equations and at most $(\log n)^b$ monomials per equation. Counts, lexicographically first witnesses and exact uniform sampling from a nonempty zero set are available in $\operatorname{poly}(n)2^{O(n\log\log n/\sqrt{\log n})}$ time and polynomial space; the sampling time is expected. The same bounds apply to the fibers of a map with those coordinate-polynomial bounds.
\end{proposition}
\paragraph{Exact inputs.}
NN-C001. full proof in this register.

E: The Exponential Time Hypothesis is False, supplied LaTeX. Appendix D Lemma D.2; Theorem 1.1 for background self-reduction.

\begin{proof}
Fix once and for all an $\mathbb F_2$-basis of $K=\mathbb F_{2^r}$ and its multiplication table. Write each variable as $x_i=\sum_{j=1}^r x_{ij}b_j$. A monomial of degree at most $k$ expands in each coordinate as a sum of at most $r^k$ products of at most $k$ bits, with coefficients determined by this fixed multiplication table. Thus each equation gives $r$ equations in $N=rn$ bits, of degree at most $k$ and polylogarithmic sparsity; the number of equations remains at most $N(\log N)^{O(1)}$. This is a bijective coordinate encoding, so it is parsimonious. Apply NN-C001 and absorb the fixed factor $r$ in the asymptotic constants.

For witness search keep all $N$ declared bits. At stage $j$, add the equations fixing the already chosen prefix and test the count after setting the next bit to zero. If that count is positive choose zero; otherwise choose one. Provided the original count is positive, at least one branch is nonempty. At most $N$ further linear equations and $N+1$ counts are required, so the same bounds hold. For a fiber $F(x)=y$, subtract the fixed coordinate $y_e$ from each polynomial; this adds at most one monomial.

For exact sampling, let $A_0,A_1$ be the two child counts at a prefix. Draw an integer uniformly from $\{0,\ldots,A_0+A_1-1\}$ by rejection from $\lceil\log_2(A_0+A_1)\rceil$ random bits. Its expected number of trials is below two. Take branch zero precisely if the integer is below $A_0$. The probability of any full satisfying assignment telescopes to $1/|S|$. The largest counts have at most $N+1$ bits and only $O(N)$ queries are needed. When the count is zero the algorithm reports that sampling is unavailable. This proves the expected-time and space claims.
\end{proof}
\paragraph{Status and complexity impact.}
Proved from NN-C001; exact field/running-time parameters stated. Independent complete deduction-proof review by root: passed; imported input theorems are not claimed independently proved. The field degree $r$ is fixed. Letting $r$ grow changes both the Boolean witness count and the monomial expansion $r^k$. The proof uses characteristic two essentially: field coordinates are $\mathbb F_2$-linear. It is not a same-parameter theorem for sparse equations over arbitrary odd-characteristic fields. Search and sampling do not imply that an NP language has a polynomial-time unambiguous verifier.

\paragraph{Geometric and representation dictionary.}
The field-coordinate construction is the elementary affine coordinate description of restriction of scalars for the finite separable extension $K/\mathbb F_2$, restricted to rational points. To obtain the finite counting scheme one also imposes $x_i^{2^r}-x_i=0$, or equivalently all bit relations $x_{ij}^2-x_{ij}=0$. Its coordinate algebra is a product indexed by the actual field-valued solutions. Counting is rank/trace; conditional sampling is repeated restriction to clopen subsets. No purity beyond the trivial weight-zero Frobenius representation on these rational components is required. SGA4 XVII 6.2.3 supplies the finite trace formalism; no claim of a canonical GCT representation is made.

\subsection{NN-C003 --- Exact rational weighted partition functions with controlled bit cost}
\begin{proposition}[NN-C003]
Fix a finite alphabet size $q$, arity $w$ and density constant $C$. Given $n$ variables with $q$ states and at most $Cn$ explicitly tabulated rational factors, each depending on at most $w$ variables, their sum-product partition function is computable exactly in
\[
 \operatorname{poly}(L+n)\,
 2^{O_{q,w,C}(n(\log\log n)^2/\sqrt{\log n})}
\]
time and polynomial space, where $L$ includes every numerator and denominator bit. Signs and zero weights are allowed. The number of factors, rather than only the number of variables, is part of the hypothesis.
\end{proposition}
\paragraph{Exact inputs.}
E: The Exponential Time Hypothesis is False, supplied LaTeX. Proposition 3.4; Lemmas 4.7--4.8; high-incidence branching in Section 5; Appendix B Lemma B.1 and its bit-height method.

\begin{proof}
Encode each $q$-state variable by $t=\lceil\log_2q\rceil$ bits and add a unary-in-the-state validity factor which is one on the $q$ chosen codes and zero on the other codes. Replace each local table by its multilinear interpolation polynomial on its at most $wt$ bits, using
\[
 \Phi(x)=\sum_{a\in\{0,1\}^{I}}\Phi(a)
 \prod_{i:a_i=1}x_i\prod_{i:a_i=0}(1-x_i).
\]
Here $t,w$ are fixed, so table sizes, polynomial degrees and scopes are bounded constants. Invalid encodings contribute zero and valid encodings correspond bijectively to assignments in the original alphabet. Let $N=tn$. The number of factors and total number of factor--variable incidences are $O(N)$, and the total product degree in independent incidence variables is $O(N)$.

If a bit belongs to more than $b_0=\max(1,\lceil M/N\rceil)\lceil\sqrt{\log_2N}\rceil$ scopes, branch on it. There are at most $M/b_0\leq N/\sqrt{\log_2N}$ such bits. Each branch has $N'=N-O(N/\sqrt{\log N})$, still $\Theta(N)$, and maximum incidence bounded by a fixed power of $\log N'$. Empty scopes are scalars. Add factors equal to one to declare any unused bits. Lemmas 4.7--4.8 supply a layout for Proposition 3.4. Its incidence degree is $O(N')$ and its coefficient algebra can be $\mathbb Q$. The resulting sum-product cost is $\operatorname{poly}(N)2^{O(N'(\log\log N')^2/\sqrt{\log N'})}$. Branching multiplies by $2^{O(N/\sqrt{\log N})}$ and hence does not change the displayed exponent.

For bit complexity, let $D_{\rm in}$ be the product of all positive table-entry denominators. Then $\log_2D_{\rm in}\leq L$, and every table coefficient has denominator dividing $D_{\rm in}$. Let $f=O(N)$ be the number of factors after validity and neutral factors are included, and let $w'$ be their fixed maximum bit arity. E Section 5, Proposition 5.3, and Appendix B use the explicit common denominator $d_{\rm den}=(R-1)!Q!$ for replacement-row coefficients. Taylor substitution introduces binomial integers but no additional denominators. Each evaluated factor therefore has denominator dividing $D_{\rm in}d_{\rm den}^{w'}$, and every group-product term and final sum can use the single denominator
\[
 \Omega=(D_{\rm in}d_{\rm den}^{w'})^f.
\]
Its bit length is at most $f(L+w'\log_2d_{\rm den})$, polynomial in $N+L$. After this common denominator is cleared, a term's numerator is a product of $O(fw')$ input and substitution coefficients, each with polynomial bit length in $N+L$. Summing the enumerated terms adds only the logarithm of their number to that bit bound, even if signs cancel.

For readout, Appendix B's Hermite matrix $\mathsf V$ has dimension $M_H=QJ=\operatorname{poly}(N)$ and entries $\binom uj t^{u-j}$, of $O(M_H\log(Q+2))$ bits. Multiplying the Taylor data by $\Omega$ makes them integral. Cramer's rule gives the requested constant term as $\det\mathsf V^{(0)}/(\Omega\det\mathsf V)$. An $h\times h$ minor with entry bit bound $B$ has bit length at most $hB+\log_2(h!)+O(1)$; this bounds all fraction-free elimination intermediates and gives polynomial bit cost. Unlike the unweighted integer application, the final answer here may remain rational, which is allowed. Each branch and its interpolation therefore use polynomial-bit arithmetic; sequential processing preserves polynomial space. The final denominator divides the product of the individual factor denominators and the output numerator also has polynomial bit length. No unit-cost assumption on large rational weights is used.
\end{proof}
\paragraph{Status and complexity impact.}
Proved extension of E Proposition 3.4; uses its weaker exponent, not an unproved extension of D.1. Independent complete deduction-proof review by root: passed; imported input theorems are not claimed independently proved. Gives signed or rational weighted CSP partition functions, bounded-arity spin models and bounded-density weighted model counting. Arbitrary weights of $B$ bits do not automatically fit D.1 with precision $P_0=O(n)$: a product can need $\Theta(nB)$ bits. This proof keeps that cost in the polynomial arithmetic overhead by using the rational-algebra routine and its weaker exponent. Dense instances with superlinear numbers of factors are not covered by this statement.

\paragraph{Geometric and representation dictionary.}
For the finite split scheme $X=\operatorname{Spec}(\mathbb Q^{q^n})$, the product of local factors is an element $W\in\Gamma(X,\mathcal O_X)$, and its sum over assignments is the trace of multiplication by $W$ in that finite etale algebra. In the idempotent basis of this product algebra, multiplication by $W$ is diagonal with entries $W(x)$, so its trace is exactly the displayed sum. EGA IV4, printed p.111,18.2.1--18.2.2 is the original finite-algebra trace construction; equation18.2.1.2 states base-change compatibility. SGA4 XVII 6.2.3 provides the relative-dimension-zero trace formalism. This is a weighted endomorphism trace, not automatically a point count or a pure Frobenius trace: arbitrary rational or negative weights need not be eigenvalues of a geometric Frobenius with the desired weight. The natural basis is indexed by assignments; there is no promised short basis for the whole exponential-dimensional algebra.

\subsection{NN-C004 --- Fixed algebraic tensor alphabets and exact quantum amplitudes}
\begin{proposition}[NN-C004]
Fix a finite set $\Omega$ of algebraic numbers and a constant maximum tensor arity. A tensor network with $s$ Boolean wires and entries from $\Omega$ has an exact deterministic contraction algorithm in $\operatorname{poly}(s+L)2^{O_\Omega(s/\sqrt{\log s})}$ time and polynomial space. Consequently a $q$-qubit, $g$-gate circuit from any fixed bounded-arity algebraic gate set has each specified computational-basis amplitude computable exactly in $\operatorname{poly}(q+g)2^{O((q+g)/\sqrt{\log(q+g)})}$ time and polynomial space. A fixed field representation and complex embedding specify the exact output.
\end{proposition}
\paragraph{Exact inputs.}
E: The Exponential Time Hypothesis is False, supplied LaTeX. Appendix D Theorem D.1, coefficient ring and precision hypotheses; its equality-tree and bipartite subdivision encoding.

J.S. Milne, Fields and Galois Theory, v5.10 (2022). Theorem5.1, pp.61--62, primitive element theorem.

\begin{proof}
Enlarge the fixed alphabet by its complex conjugates if probabilities will be requested, and let $K$ be the resulting fixed number field. By Milne Theorem5.1 choose a primitive element $\beta$ for $K$. If $a_dT^d+\cdots+a_0$ is an integer polynomial annihilating $\beta$ with $a_d\ne0$, then $\alpha=a_d\beta$ is integral: multiply the relation by $a_d^{d-1}$ to obtain a monic integer equation for $\alpha$. It still generates $K$. Its degree $d$ is constant. There is a fixed positive integer $c$ such that every $\omega\in\Omega$ can be written
\[
 \omega=c^{-1}\sum_{j=0}^{d-1}a_{\omega,j}\alpha^j,
 \qquad a_{\omega,j}\in\mathbb Z.
\]
Remove isolated tensors, multiplying their entries separately at the end. The remaining number of vertices $v$ is at most $2s$. Introduce a formal variable $V$ and replace each tensor entry by its integer polynomial $\sum_j a_{\omega,j}V^j$. The contraction is an integer polynomial $H(V)$ of degree at most $(d-1)v$. If $A=\max(1,\max_\omega\sum_j|a_{\omega,j}|)$, then
\[
 \|H\|_1\leq 2^s A^v.
\]
Thus choose $\ell>(d-1)v$ and $P_0$ with $2^{P_0-1}>2^sA^v$; both are $O_\Omega(s)$. Reduction modulo $V^\ell$ loses no coefficient, and reduction modulo $2^{P_0}$ permits unique signed-integer recovery of every coefficient.

Subdivide every edge with an equality tensor to make the graph bipartite. Its wire count changes by a factor two, its maximum degree stays bounded and it has $O(s)$ nonisolated vertices. Its local entries are computable in polynomial time in the required ring $(\mathbb Z/2^{P_0}\mathbb Z)[V]/(V^\ell)$. All hypotheses of D.1 hold after increasing its fixed constant $C$. Recover $H$ coefficientwise, evaluate $H(\alpha)/c^v$, reduce to the fixed basis of $K$, and multiply the removed scalar tensors. This last arithmetic has polynomial output height and polynomial cost in the full description length.

For the quantum specialization, insert each gate as its input/output tensor and each fixed initial or final basis state as a unary tensor. A bounded-arity circuit has $O(q+g)$ wire segments: the sum of numbers of incidences at gates is $O(g)$, and each qubit contributes its initial segment. The tensor contraction is the path-sum formula for the specified amplitude. Apply the preceding result. A probability for a fixed outcome is obtained by multiplying the exact amplitude by its complex conjugate in the fixed field, with polynomial extra work.
\end{proof}
\paragraph{Status and complexity impact.}
Proved from E D.1 with explicit precision and algebraic-height bounds. Independent complete deduction-proof review by root: passed; imported input theorems are not claimed independently proved. This is a subexponential algorithm in total circuit size $q+g$, not an efficient simulation in the input-qubit parameter for arbitrary polynomial-size circuits. It gives no BQP=P consequence. The fixed finite algebraic alphabet is essential to the stated exponent: variable coefficient heights or growing number-field degree require separate accounting. The result concerns numerical contraction, not a small symbolic arithmetic circuit for all tensor entries.

\paragraph{Geometric and representation dictionary.}
The contraction is composition of evaluation/coevaluation pairings in finite-dimensional $K$-vector spaces, and can also be written as the finite-algebra trace of a diagonal weighted assignment operator as in NN-C003. The diagonal matrix calculation in NN-C003 supplies that algebraic trace dictionary directly; SGA4 XVII 6.2.3 is the verified finite-trace comparison. The temporary $V$ is a coefficient-encoding device; it is not a Frobenius coordinate, deformation of a motive, or canonical-basis parameter. Fixed algebraicity ensures exact arithmetic but does not imply purity or a representation-theoretic obstruction.

\subsection{NN-C005 --- Parity is a top Boolean coefficient, with a hard sparse quadratic boundary}
\begin{proposition}[NN-C005]
For polynomials $f_1,\ldots,f_m\in\mathbb F_2[x_1,\ldots,x_n]$, let $N$ count their common zeros in $\mathbb F_2^n$ and let $\operatorname{ML}$ denote reduction modulo $(x_i^2-x_i)$. Then
\[
 N\bmod2=[x_1\cdots x_n]\operatorname{ML}\!\left(\prod_{i=1}^m(1+f_i)\right).
\]
If $\sum_i\deg f_i<n$, this parity is zero. Nevertheless computing this coefficient is $\oplus\mathrm P$-complete already for $m=O(n)$ quadratic equations with at most three monomials per equation. Exact counting on the same class is $\#\mathrm P$-complete under parsimonious polynomial-time reductions.
\end{proposition}
\paragraph{Exact inputs.}
Elementary Boolean-coordinate-algebra calculation. full proof below.

Katz, On a theorem of Ax, AJM 93 (1971), 485--499. Theorem 1.0 and definition 0.3.

E: The Exponential Time Hypothesis is False, supplied LaTeX. Appendix D Lemma D.2, algorithm for this hard class.

\begin{proof}
On each bit assignment, $\prod_i(1+f_i(x))$ equals one exactly at a common zero. For a squarefree monomial $x^S$, its sum over all bit assignments equals $2^{n-|S|}$, reduced modulo two. This is one for $S=[n]$ and zero otherwise. Summing the multilinear representative proves the coefficient identity. Reduction to squarefree monomials never increases degree, so the top coefficient vanishes when the displayed degree sum is below $n$.

For hardness, take an arbitrary polynomial-time Boolean witness verifier and compile it, on a fixed instance, into an AND/NOT circuit. Keep its $u$ witness-input bits and introduce one bit $z_g$ for each of its $s$ gates. Encode an AND gate by $z_g+ab=0$, a NOT gate by $z_g+a+1=0$, and acceptance by $z_{\rm out}+1=0$. Constants can be substituted directly. There are $N'=u+s$ variables and $s+1=O(N')$ equations, each of degree at most two and with at most three monomials. For every assignment to the witness bits, all gate bits are uniquely determined, in topological order. The equations have one extension if that witness accepts and none otherwise. This is a parsimonious reduction from every $\#\mathrm P$ verifier; taking parity proves the $\oplus\mathrm P$ claim. Membership follows by guessing the bit assignment and checking the equations.

This also explains why the low-degree divisibility theorem does not solve the general hard instances. Over $\mathbb F_q$, Katz's Theorem 1.0 states $q^\mu\mid N$ with $\mu=\max(0,\lceil(n-\sum_i d_i)/\max_i d_i\rceil)$ for nonconstant defining polynomials of degrees $d_i$. If $\mu=0$ it imposes no divisibility. Even when $\mu>0$, the congruence does not distinguish zero from a positive multiple of $q^\mu$. In the circuit encoding the sum of degrees generally meets or exceeds the number of variables. Adding all Boolean relations to make a finite scheme also destroys the small degree-sum hypothesis. E provides the stated subexponential algorithm for the hard class, not a polynomial coefficient-extraction algorithm.
\end{proof}
\paragraph{Status and complexity impact.}
Identity and reduction proved in full; Katz divisibility used by exact classical locator. Independent complete deduction-proof review by root: passed; imported input theorems are not claimed independently proved. Identifies the precise computational gate behind a mod-two cohomological count. Low individual degree and constant monomial sparsity do not make parity easy when the number of equations grows linearly. E changes the exponential running time, while $\oplus\mathrm P$-completeness remains valid.

\paragraph{Geometric and representation dictionary.}
The Boolean quotient is the coordinate algebra of a finite split etale scheme. The coefficient functional above is its sum-of-evaluations trace over $\mathbb F_2$, expressed in the squarefree monomial basis. This basis has $2^n$ elements. Its finite scheme has weight-zero Frobenius but the trace modulo two is still $\oplus\mathrm P$-complete from a succinct equation presentation. SGA7 II XXII Corollary 3.2 gives exactly the coherent trace congruence, with properness satisfied for this finite scheme. It gives no polynomial algorithm for constructing an explicit cohomology basis.

\subsection{NN-C006 --- The effective Frobenius-data gate and the correct affine boundary}
\begin{proposition}[NN-C006]
Suppose a deterministic polynomial-time procedure, on every sparse quadratic system in NN-C005, outputs polynomially many explicitly represented $\mathbb F_2$ matrices of polynomial total dimension whose alternating trace equals the number of common bit zeros modulo two. Then $\mathrm P=\oplus\mathrm P$ and $\mathrm{NP}=\mathrm{RP}$. If it instead outputs the exact integer rank/count in polynomial time, then $\mathrm{FP}=\#\mathrm P$ and $\mathrm{PH}=\mathrm P$. The mere existence of finite coherent or etale cohomology is insufficient. For a nonproper affine scheme, use a proper compactification and its boundary, or compactly supported etale cohomology, rather than the ordinary affine coherent $H^0$.
\end{proposition}
\paragraph{Exact inputs.}
NN-C005. parsimonious sparse-quadratic reduction.

Valiant--Vazirani, NP is as easy as detecting unique solutions, TCS 47 (1986), 85--93. Theorem 1.1; Corollaries 3.1--3.4, especially Corollary 3.2 p.91.

Toda, PP is as Hard as the Polynomial-Time Hierarchy, SIAM J. Comput. 20 (1991), 865--877. main theorem PH subset P\textasciicircum{}PP; publisher abstract verified, internal numbering not claimed.

SGA 7 II, Expose XXII. Theorem 3.1 (3.1.1); Corollary 3.2 (3.2.1).

SGA 5, Expose III. equation 6.13.2; Theorem 6.13.3.

\begin{proof}
The trace of an explicit matrix is the sum of its diagonal entries. Therefore the assumed construction followed by polynomially many field operations decides the parity problem in polynomial time. By NN-C005 every language in $\oplus\mathrm P$ reduces to that problem in polynomial time, yielding $\oplus\mathrm P\subseteq\mathrm P$; the reverse containment is immediate. Valiant--Vazirani's one-sided randomized isolation reduction maps an unsatisfiable formula only to unsatisfiable formulas and maps a satisfiable formula to a uniquely satisfiable formula with inverse-polynomial probability. A parity oracle always rejects the former and accepts the latter. Polynomial amplification gives SAT in RP. Since RP is contained in NP and SAT is NP-complete, NP=RP.

For the integer version, the parsimonious reduction makes all $\#\mathrm P$ functions polynomial-time computable; all polynomial-time nonnegative polynomial-output-length functions can also be counted by choosing a binary integer below their value, so the standard functional equality follows. A PP comparison is then polynomial-time computable. Toda's main theorem puts PH in $\mathrm P^{\mathrm{PP}}$, which now equals P. This integer conclusion is stronger than the parity conclusion and must not be inferred from a single residue.

The classical trace expressions are mathematically valid without any efficiency claim. For the finite Boolean solution scheme, $H^0(\mathcal O)$ is $\mathbb F_2^N$, all higher coherent cohomology vanishes, and Frobenius acts as the identity. That matrix can have exponential dimension. An arbitrary affine scheme does not satisfy the proper coherent formula: for example $H^0(\mathbb A^1,\mathcal O)=\mathbb F_q[x]$ is infinite-dimensional. If $U\subseteq\overline U$ is an open immersion into a proper finite-type scheme and $D=(\overline U\setminus U)_{\rm red}$, then $|U(\mathbb F_q)|=|\overline U(\mathbb F_q)|-|D(\mathbb F_q)|$. Applying SGA7 XXII Corollary 3.2 to both proper schemes gives the correct difference of coherent traces modulo $p$. Alternatively SGA5 III 6.13.3 supplies the compact-support $\mathbb F_p$ etale trace formula. Neither operation supplies polynomial-size matrices or an algorithm to construct them.
\end{proof}
\paragraph{Status and complexity impact.}
Conditional gate proved; classical trace hypotheses audited; no unconditional complexity collapse claimed. Independent complete deduction-proof review by root: passed; imported input theorems are not claimed independently proved. Separates purity/existence from effective presentation. A, M and a Boolean support object do not by themselves provide the hypothesized matrix construction. The parity reduction would derandomize an NP decision only to RP through Valiant--Vazirani; claiming NP=P at this point would require an additional deterministic isolation gate.

\paragraph{Geometric and representation dictionary.}
The dictionary is literal: solution set to finite etale scheme, count to rank, parity to coherent Frobenius trace, and affine restriction to proper object minus boundary. On the representation side $H^0$ is a direct sum of trivial Frobenius representations indexed by solutions. Weight zero controls eigenvalue magnitude but not the cost of extracting the number of summands from equations. SGA7 II XXII 3.1--3.2 and SGA5 III 6.13 are the exact locators; no Grothendieck archive document is invented.

\subsection{NN-C007 --- Exact deterministic-isolation gates and why logspace isolation does not transfer automatically}
\begin{proposition}[NN-C007]
A polynomial-time computable polynomial-length list of witness restrictions, each forbidding all witnesses on no-instances and containing an odd number of witnesses for at least one list member on each yes-instance, would imply $\mathrm{NP}\subseteq\oplus\mathrm P$. If all these parity tests were in P it would imply NP=P. A single polynomial-time restriction that always has at most one witness and preserves nonemptiness would imply NP=UP. E supplies such a single isolating restriction only at its subexponential running time. Hoza's short seed generator does not give the required polynomial-size list for arbitrary NP computations.
\end{proposition}
\paragraph{Exact inputs.}
Valiant--Vazirani, NP is as easy as detecting unique solutions, TCS 47 (1986), 85--93. Theorem 1.1; Corollaries 3.1--3.4, especially Corollary 3.2 p.91.

Hoza, NL = UL, ECCC TR26-323, supplied PDF. Theorems 1.1 and 1.3; Definition 1.2; Section 1.2.

E: The Exponential Time Hypothesis is False, supplied LaTeX. Theorem 1.1 and Appendix D Lemma D.2, conditional-count self-reduction.

\begin{proof}
Let the list have $t=t(|x|)$ members with witness counts $a_1,\ldots,a_t$. In $\mathbb F_2$ the OR of their parities is
\[
 1+\prod_{i=1}^t(1+a_i).
\]
This expression is a $\#\mathrm P$ count modulo two: a nondeterministic machine has one distinguished accepting branch, and on the other branch independently chooses, for each $i$, either one dummy option or an accepting witness for restriction $i$. It accepts exactly the valid tuples. There are $1+\prod_i(1+a_i)$ accepting paths, and the total nondeterministic witness length is polynomial. On no-instances every $a_i=0$, so the parity is zero; on yes-instances some $a_i$ is odd, so the parity is one. This proves the $\oplus\mathrm P$ inclusion. If each list parity is deterministically computable in polynomial time, ordinary OR decides the NP language in P.

For the single-restriction assertion, guess its restricted witness and check it. The global promise of at most one witness makes this an unambiguous polynomial-time verifier, and preservation of nonemptiness gives the original language. Hence every NP language lies in UP; the reverse inclusion is definitional. A list with some uniquely satisfiable member is not enough for this argument: other members can accept and there is no supplied unambiguous way to select the first good member. The parity-list argument does not cure that uniqueness problem.

E counts under added unit equations. Repeatedly count the zero branch of the next prefix bit, choosing a nonempty branch, until one satisfying assignment remains. At most $n$ count queries and $n$ unit constraints suffice; an unsatisfiable input can be sent to a fixed inconsistent restriction. The construction thus preserves nonemptiness and yields at most one witness, but its time is E's subexponential bound, not polynomial.

Hoza's Theorem 1.3 gives a seed length $O(\log(wn))$ for width-$w$, length-$n$ layered graphs, hence a polynomial-size seed list when the graph has polynomial width and length. An NP verifier can have a configuration width $2^{\operatorname{poly}(N)}$ on inputs of length $N$. For that graph the seed bound is polynomial in $N$, allowing exponentially many seeds. The min-isolation test also uses the specific Reinhardt--Allender unambiguous reachability machinery cited in Hoza Section 1.2; it is not a generic polynomial-time test that an arbitrary NP witness restriction is unique. These two parameter/interface gaps prevent the proposed automatic NP transfer.
\end{proof}
\paragraph{Status and complexity impact.}
Conditional implications and E self-reduction proved; Hoza parameter gate audited. Independent complete deduction-proof review by root: passed; imported input theorems are not claimed independently proved. Makes the desired NP analogues concrete and reviewable. Randomized Valiant--Vazirani plus efficient parity yields NP=RP; deterministic polynomial-list odd isolation yields NP subset parity-P; globally unique deterministic isolation yields NP=UP. These are different conclusions with different extra hypotheses.

\paragraph{Geometric and representation dictionary.}
A restriction cuts a finite Boolean solution scheme by additional equations, equivalently selects a clopen collection of its rational components. Odd isolation asks that one selected finite scheme have odd degree; unique isolation asks for degree at most one. The finite trace of SGA4 XVII 6.2.3 records these degrees. A weight function's unique minimum is a combinatorial selection on this set, not a consequence of Frobenius purity. In the representation dictionary, direct sums encode alternative restrictions and tensor products encode simultaneous witness choices; the displayed parity OR uses their mod-two rank arithmetic.

\subsection{NN-C008 --- Subexponential computation-transcript proofs, without a natural-system size claim}
\begin{proposition}[NN-C008]
There exists a Cook--Reckhow propositional proof system in which an unsatisfiable fixed-width $k$-CNF with $n$ variables and $m$ clauses has a refutation of length $\operatorname{poly}_k(n+m)2^{O_k(n\log\log n/\sqrt{\log n})}$; at fixed clause density the sharper exponent $O(n/\sqrt{\log n})$ applies. Every tautological Boolean formula of size $L$ has a proof in this system of length $\operatorname{poly}(L)2^{O(L/\sqrt{\log L})}$. Analogous computation certificates certify emptiness of the sparse systems of NN-C001.
\end{proposition}
\paragraph{Exact inputs.}
E: The Exponential Time Hypothesis is False, supplied LaTeX. Theorem 1.1, including polynomial space and fixed-density bound.

Cook--Reckhow proof-system definition. Polynomial-time decidable proof relation, constructed explicitly below.

\begin{proof}
Fix E's deterministic exact counting machine, including a standard binary implementation of its arithmetic, and encode each complete configuration by a polynomial number of bits in the input length. A certificate for a CNF consists of its full computation transcript ending with count zero. Check that the first configuration is the prescribed start configuration, every successive pair obeys the fixed local transition relation, and the final output is zero. Verification is polynomial in the transcript length. Soundness follows from determinism and correctness of E: an accepted transcript is its actual computation and a zero count means unsatisfiability. Completeness follows by recording that computation for an unsatisfiable formula. If the run has $T$ steps and space $S$, the transcript length is $O(TS)$ up to encoding overhead. Since $S$ is polynomial, this has the asserted size.

To obtain a proof system for all tautologies, an input proof additionally contains a formula $\varphi$ and the canonical Tseitin encoding of a circuit for $\neg\varphi$, with the output forced to one. For a size-$L$ formula this has $O(L)$ variables and $O(L)$ clauses of width at most three. Local gate constraints have a unique extension for every assignment to the original variables, so the CNF is unsatisfiable exactly when $\varphi$ is a tautology. The verifier checks this deterministic encoding and an E transcript for it. Fixed-density Theorem 1.1 gives $2^{O(L/\sqrt{\log L})}$ up to polynomial factors. On invalid proof strings the Cook--Reckhow output-function version returns a fixed tautology, so its range is exactly TAUT. Replacing E's CNF solver by NN-C001 similarly gives polynomially checkable computation certificates of sparse-system emptiness.
\end{proof}
\paragraph{Status and complexity impact.}
Proved existence of an algorithm-dependent proof system; no claim about resolution, PC, Frege or IPS. Independent complete deduction-proof review by root: passed; imported input theorems are not claimed independently proved. The proofs are superpolynomial in general, so this does not imply NP=coNP. The verification procedure checks a Turing computation; it supplies no bound on width, degree, monomial count or size in a pre-existing algebraic or propositional proof system.

\paragraph{Geometric and representation dictionary.}
Unsatisfiability is emptiness of the finite Boolean solution scheme. A computation transcript certifies this from its equations; it is not automatically an effective Nullstellensatz identity or a short free resolution of its coordinate algebra. In the finite etale model, emptiness means rank zero and zero $H^0$, while SGA7 II XXII 3.2 only identifies a residue of that rank. E computes the integer rank through its own algorithm, not through a claimed small cohomological presentation.

\subsection{NN-C009 --- Small semantic quotient certificates do not supply short polynomial-calculus proofs}
\begin{proposition}[NN-C009]
Every unsatisfiable Boolean polynomial system has a small shared-circuit identity proving $1$ belongs to its generated ideal in the Boolean function algebra. However checking that identity modulo the Boolean ideal is coNP-hard, even when the defining equations are the sparse quadratics in NN-C005. E does not supply a polynomial bound for ordinary Nullstellensatz or polynomial-calculus certificates.
\end{proposition}
\paragraph{Exact inputs.}
NN-C005. Boolean quotient and parsimonious gate encoding.

E: The Exponential Time Hypothesis is False, supplied LaTeX. Theorem 1.1 / Lemma D.2 for the algorithmic distinction.

\begin{proof}
In characteristic two put $g_i=\prod_{j<i}(1+f_j)$. A telescoping product gives the formal polynomial identity
\[
 \sum_{i=1}^m f_i g_i=1+\prod_{i=1}^m(1+f_i).
\]
If there is no common bit zero, the product on the right vanishes at every Boolean assignment and is zero in $B_n=\mathbb F_2[x]/(x_i^2-x_i)$. Thus $\sum_i f_i g_i=1$ in $B_n$. All prefix products can share intermediate gates, so the collection of the $g_i$ and the displayed sum has size linear in the total supplied circuit size plus $m$.

This compact expression is not a polynomial-time checkable proof of the crucial vanishing assertion. Given any Boolean verifier circuit, construct the equations of NN-C005 and their product $P=\prod_i(1+f_i)$. The circuit for $P$ has polynomial size. It represents the zero element of $B_n$ if and only if the original verifier has no accepting input. Therefore deterministic polynomial-time verification of these semantic quotient identities would decide UNSAT in P, and hence P=NP. The construction itself does not solve that check.

In the ordinary polynomial ring, an actual Nullstellensatz certificate requires additional polynomials $h_j$ satisfying
\[
 1=\sum_i f_i g_i+\sum_{j=1}^n h_j(x)(x_j^2-x_j).
\]
The vanishing product is indeed in the Boolean ideal because $B_n$ is the algebra of all Boolean functions, but the above small circuit expression has not constructed small $h_j$ or a small derivation of membership. Expansion or multilinear reduction can involve exponentially many monomials. A transcript of E as in NN-C008 is a valid algorithmic certificate; changing its proof system to polynomial calculus, resolution or an ideal-proof system requires an additional simulation theorem with quantitative bounds. NN-C029 subsequently supplies subexponential ordinary ideal multipliers over $\mathbb Q$ for bounded-density local verifier systems by a different construction. NN-C033 supplies the fixed-characteristic circuit transfer, and NN-C034 constructs integral constant-free certificates valid after every base change. None of these gives polynomial-size multipliers or a subexponential ordinary monomial-PC simulation; NN-C021 and NN-C034 instead exhibit an exponential ratio between monomial-PC and circuit-IPS size on the same family.
\end{proof}
\paragraph{Status and complexity impact.}
Explicit identity and verification hardness proved; later NN-C029--NN-C034 supply ordinary circuit certificates while retaining the distinction from monomial polynomial calculus. Independent complete deduction-proof review by root: passed; imported input theorems are not claimed independently proved. Blocks the inference that detecting or removing supported arithmetic zeros automatically produces short algebraic refutations. It identifies the missing effective ideal-membership/derivation gate, without claiming an impossibility result for every potential simulation of E.

\paragraph{Geometric and representation dictionary.}
The quotient $B_n$ is the coordinate algebra of a finite split scheme; its idempotent decomposition makes semantic ideal membership transparent after all components are listed. The telescoping identity is an ideal-theoretic statement on that finite scheme. Producing polynomial-ring correction terms is a presentation problem, analogous to replacing an abstract zero object by an explicit exact complex. Rank or Frobenius-trace identities in SGA7 XXII do not control those presentation sizes. The translation is algebraic; no unproved categorical purity theorem is used.

\subsection{NN-C010 --- Burgisser transfer: the exact Dedekind family and the one-way conclusion}
\begin{proposition}[NN-C010]
The characteristic-zero Cook--Valiant transfer needs RH for Dedekind zeta functions of the bounded-degree, bounded-height number fields furnishing algebraic solutions to the constant-replacement systems in its proof. These fields are generally arbitrary number fields, not necessarily abelian or normal. Conditional on M, the needed RH is available, and Burgisser's transfer gives
\[
 \mathrm{VP}_{\mathbb C}=\mathrm{VNP}_{\mathbb C}
 \ \Longrightarrow\ \#\mathrm P/\mathrm{poly}=\mathrm{FP}/\mathrm{poly},
\]
with the stronger Boolean-part and nonuniform-collapse conclusions at the cited classical locators. Its contrapositive is a Boolean nonuniform counting lower bound implying an algebraic separation. A suffices for the restricted proof route when those witness fields can be chosen abelian with the same quantitative bounds; no supplied theorem establishes that restriction in general. No converse transfer is supplied.
\end{proposition}
\paragraph{Exact inputs.}
Burgisser, Completeness and Reduction in Algebraic Complexity Theory, Springer 2000, Chapter 4. Theorem 4.5(A), Corollary 4.6(1), pp.64--65; Theorem 4.12 and Remark 4.17, pp.71--75; Theorem 4.18 and Corollary 4.19, p.76; proof of 4.5(A3), pp.77--78.

A: supplied Alpoge Lean archive. DirichletLNonvanishingRightOfHalf; DirichletLZerosOnLine; LFunction\_ne\_zero\_of\_half\_lt\_re; LFunction\_eq\_zero\_re; Remark p.4 step20.

M: Masullo claimed general Hecke GRH. Thm12.5.1.1 pp.519--520; Thm13.3.1.1 p.522; all uses conditional on M.

Prior CX register. CX-K01--CX-K03: analogous Koiran witness-field audit, used as background only.

\begin{proof}
The classical proof to be imported is Burgisser, Chapter 4, proof of Theorem 4.5(A3), with Theorem 4.12 and the root-modulo-primes estimate of Corollary 4.19. It replaces arbitrary circuit constants by variables and imposes agreement on all Boolean inputs. A bounded algebraic solution is represented in one number field $K=\mathbb Q(\alpha)$. The prime-ideal theorem used in Section 4.6 assumes RH for $\zeta_K$; degree-one prime ideals yield roots modulo rational primes. The proof does not pass to a Galois closure. The constants of one sufficiently large finite-field specialization become polynomial-length nonuniform advice. Theorem 4.5(A) and Corollary 4.6(1) state the resulting transfer and collapse. These are exact classical proof locators, rather than an asserted new proof of that theorem.

For the new premise audit, if $K/\mathbb Q$ is finite abelian, its Dedekind zeta function is a product of the corresponding primitive Dirichlet $L$-functions. A therefore gives its RH. If $K$ is arbitrary, the trivial Hecke character of $K$ has $L$-function $\zeta_K$, so M supplies precisely the missing input, conditional on M. Neither statement requires holomorphy of a nontrivial Artin $L$-function. The existence of algebraic solutions alone does not make their fields abelian: for example $X^3-2=0$ has roots generating non-Galois cubic fields; none is contained in an abelian extension of $\mathbb Q$, since every subextension of an abelian Galois extension is Galois. This example illustrates the field-class gap, not a proof that any logically possible consequence of A cannot cover it.

Finally take the contrapositive of the displayed implication. E's supplied lower bound concerns near-linear circuits for a language in $\mathrm E^{\mathrm{NP}}$. It is not a superpolynomial lower bound for the Boolean functions computing $\#\mathrm P$ values, and hence does not instantiate the antecedent. Uniform $\mathrm P\ne\mathrm{NP}$ alone also does not state a nonuniform counting lower bound. Conversely, an algebraic separation does not reverse an implication whose antecedent is algebraic equality. Thus the combination closes the analytic premise only under M (or a separately proved abelian-witness promise), leaving the Boolean lower-bound gate intact.
\end{proof}
\paragraph{Status and complexity impact.}
Classical transfer by exact locator; field-family audit proved; general characteristic-zero use conditional on M. Independent complete deduction-proof review by root: passed; imported input theorems are not claimed independently proved. Use function notation $\#\mathrm P\not\subseteq\mathrm{FP}/\mathrm{poly}$, or explicitly encode output bits as Boolean functions, rather than silently writing a counting class as a decision class. Under the same classical hypothesis, Corollary 4.6(1) also gives $\mathrm{NC}^3/\mathrm{poly}=\mathrm P/\mathrm{poly}=\mathrm{NP}/\mathrm{poly}=\mathrm{PH}/\mathrm{poly}$ and a uniform second-level collapse of PH if VP=VNP. These do not assert VP=VNP, and E does not choose an alternative.

\paragraph{Geometric and representation dictionary.}
The constant-replacement equations define an affine scheme over $\mathbb Z$. A closed point of its characteristic-zero fiber supplies $\operatorname{Spec}K$; reductions with residue degree one give rational points over finite fields. Its finite permutation representation has Dedekind zeta function $\zeta_K$. The relevant trace counts prime-ideal splitting, not satisfying assignments. The elementary matrix trace of multiplication supplies the finite-algebra language, and Burgisser Section 4.6 supplies the exact arithmetic estimate. A covers abelian character decompositions; the general field need not admit a decomposition into Dirichlet characters. This is the same type of analytic field gate as Koiran, but the output here is nonuniform circuit simulation.

\subsection{NN-C011 --- Inversion and collision attacks for sparse low-degree maps}
\begin{proposition}[NN-C011]
Fix $r,k,a,b$. Let $F_n:\mathbb F_{2^r}^n\to\mathbb F_{2^r}^{m(n)}$ have $m(n)\leq n(\log n)^a$ explicitly given coordinate polynomials of degree at most $k$ and sparsity at most $(\log n)^b$. There is a deterministic $2^{O(n\log\log n/\sqrt{\log n})}\operatorname{poly}(n)$ algorithm that inverts every image point. The same bound finds a collision whenever one exists, and counts ordered unequal colliding pairs. Thus these maps cannot have $2^{\Omega(n)}$ inversion or collision security against all algorithms in a security parameter proportional to $n$. This does not rule out ordinary polynomial-time one-wayness.
\end{proposition}
\paragraph{Exact inputs.}
NN-C001--NN-C002. quasilinear sparse-system counting and search.

E: The Exponential Time Hypothesis is False, supplied LaTeX. Appendix C Corollary C.2 is the previously known constant-locality special case, not newly claimed here.

\begin{proof}
For inversion solve the equations $F_{n,e}(x)-y_e=0$. The parameter bounds and exact preimage search are those of NN-C002. The algorithm succeeds for every $y$ in the image, so in particular its success probability is one when $y$ is sampled as $F_n(x)$ under any distribution on $x$.

For collisions introduce two input vectors $x,x'$, giving $2n$ field variables, and impose $F_{n,e}(x)-F_{n,e}(x')=0$. The coordinate degree remains at most $k$ and sparsity at most twice the original bound. After conversion to the fixed binary basis there are $2rn$ bits. Partition the unequal pairs according to the first bit position $j$ where the binary encodings differ. In branch $j$, add linear equations $x_i+x'_i=0$ for $i<j$ and $x_j+x'_j+1=0$. These at most $rn$ equations preserve the quasilinear equation bound. Every unequal pair belongs to exactly one branch. Count each branch with NN-C001 and sum to obtain the ordered unequal-pair count; search any nonempty branch for a collision. There are only $rn$ branches and replacing $n$ by $2rn$ changes the exponent by a fixed factor. If the field-output length is smaller than the field-input length, a collision exists by the pigeonhole principle; otherwise the algorithm can correctly report none.

For every fixed $c>0$, the algorithm's exponent divided by $n$ tends to zero, so its running time is eventually below $2^{cn}$. This contradicts a claimed lower bound against every algorithm at that exponential security scale, including an average-case claim requiring negligible attack success, because the inversion attack succeeds on every image. A security parameter with an unrelated scaling must be substituted explicitly. A subexponential algorithm is still superpolynomial, so the standard definition of a one-way function, which excludes only polynomial-time adversaries, is not contradicted.
\end{proof}
\paragraph{Status and complexity impact.}
Proved attack consequence for the stated representation and parameter; no blanket claim about deployed cryptosystems. Independent complete deduction-proof review by root: passed; imported input theorems are not claimed independently proved. Extends E's constant-locality inversion corollary to polylogarithmic monomial sparsity and a quasilinear output length, and adds a direct collision-counting construction. Dense multivariate quadratics with $\Theta(n^2)$ monomials per coordinate are outside this theorem. A generic polynomial-size Boolean circuit encoding has $n+s$ variables and can lose subexponentiality in the intended security parameter; the input representation is part of the result.

\paragraph{Geometric and representation dictionary.}
An inversion fiber is a rational-point fiber of the polynomial map $F_n$. Collisions are rational points of its fiber product with itself, with the diagonal removed; the first-differing-bit decomposition is a disjoint clopen stratification of this finite Boolean scheme. Counts are degrees of those finite etale pieces, with trace formalism SGA4 XVII 6.2.3. This construction supplies an algorithm for the succinct fibers, not a theorem that GRH or purity makes arbitrary fibers easy.

\subsection{NN-C012 --- PIT transfer fails at both the identity interface and the time quantifiers}
\begin{proposition}[NN-C012]
E's Boolean sparse-system algorithm does not, through the stated reductions, satisfy the Kabanets--Impagliazzo formal-PIT hypothesis. First, vanishing on the Boolean cube differs from polynomial identity. Second, $2^{O(N\log\log N/\sqrt{\log N})}$ is a mildly subexponential bound, not the stronger $\mathrm{NSUBEXP}=\bigcap_{\epsilon>0}\mathrm{NTIME}(2^{N^\epsilon})$ bound used by the cited PIT transfer. Third, general circuit encodings charge their gates in the witness variable count. A current independent unconditional arithmetic/Boolean lower-bound disjunction remains a disjunction and does not close these gates.
\end{proposition}
\paragraph{Exact inputs.}
E: The Exponential Time Hypothesis is False, supplied LaTeX. Appendix D Lemma D.2; Corollary 1.4 and Appendix C Corollary C.1.

Kabanets--Impagliazzo, Derandomizing Polynomial Identity Tests Means Proving Circuit Lower Bounds, Computational Complexity 13 (2004), 1--46. Main theorem as stated on author research page; refined statement at CGGM26 Theorem 1.1.

Chukhin--Gajulapalli--Golovnev--Mihajlin, A Little Derandomization Goes a Long Way, ECCC TR26-313, 5 Nov 2026. Theorem 1 p.5; Theorem 1.1 and definition of NSUBEXP p.3; Theorem 8.8 and proof of Theorem 1 pp.46--47.

\begin{proof}
The nonzero polynomial $x^2-x$ is zero on both Boolean inputs. Thus any procedure whose only assertion is whether a circuit or polynomial vanishes on the Boolean cube cannot serve as formal PIT. Over a finite field one can instead evaluate in a sufficiently large extension field, but then a reduction must encode that extension and every arithmetic gate; the constant-degree sparse-polynomial interface is not an automatic encoding of an arbitrary formal circuit.

The time issue is independent. Set $g(N)=N\log\log N/\sqrt{\log N}$. For every fixed $0<\epsilon<1$,
\[
 \frac{g(N)}{N^\epsilon}
 =N^{1-\epsilon}\frac{\log\log N}{\sqrt{\log N}}\longrightarrow\infty.
\]
Therefore the upper bound supplied by E cannot be replaced by $2^{N^\epsilon}$ for every $\epsilon>0$. This is a statement about what the known bound guarantees, not a lower bound against possible future faster algorithms. Even a linear reduction into E's interface would not supply the required NSUBEXP guarantee.

For circuit encodings, an $n$-input circuit with $s$ gates has the parsimonious sparse quadratic encoding in $n+s$ variables from NN-C005. Taking $s=n^d$ gives a bound exponential in approximately $n^d/\sqrt{\log n}$, not the all-polynomial-circuit-size $2^n/\operatorname{poly}(n)$ type saving required by standard stronger SAT-to-lower-bound transfers. The already supplied E circuit corollary is valid at its own exact near-linear scale; it cannot be promoted by dropping $s$.

For current-source completeness, CGGM26 Theorem 1 states the unconditional alternative $\mathrm{NE}\not\subseteq\mathrm{SIZE}[2^n/(2n)]$ or $\mathrm{VP}_{\mathbb Z}^{0}\ne\mathrm{VNP}_{\mathbb Z}^{0}$. This theorem is independent of E and A. It does not select the algebraic alternative, and its constant-free formal-degree-bounded classes must not be silently replaced by ordinary $\mathrm{VP}_{\mathbb C}$ and $\mathrm{VNP}_{\mathbb C}$. Its precise quantitative PIT gates are useful further research targets, but E has not met them here.
\end{proof}
\paragraph{Status and complexity impact.}
Explicit counterexample and asymptotic/interface audit proved; no formal independence claim. Independent complete deduction-proof review by root: passed; imported input theorems are not claimed independently proved. Prevents conflating ETH-style $2^{o(n)}$ savings with the stronger SUBEXP quantifier, and numerical Boolean evaluation with symbolic arithmetic identity. No PRG, hitting set or BPP=P conclusion follows from the supplied E circuit lower bound alone. CGGM26 supplies an up-to-date background disjunction, not a new consequence counted as ours.

\paragraph{Geometric and representation dictionary.}
Formal identity asks whether a regular function is zero in an affine coordinate ring. Boolean testing asks whether its restriction to the finite reduced subscheme cut out by $x_i^2-x_i$ is zero. The restriction map has a nonzero kernel, as the example shows. A trace on the finite subscheme forgets still more data. Finite-algebra matrix traces and the SGA7 XXII trace statement therefore do not identify formal PIT with Boolean zero testing. On the representation side, one must retain the polynomial family and its arithmetic circuit model; the finite evaluation representation is only a quotient.

\subsection{NN-C013 --- Full energy histograms and sharper Hamming-weight enumerators}
\begin{proposition}[NN-C013]
Fix $k,w,C,d$. Let $F$ be any $k$-CNF on $n$ declared variables. Suppose an energy $E(x)$ is the sum of at most $Cn$ explicitly tabulated nonnegative integer local costs, each on at most $w$ bits, with total possible range at most $M\leq n^d$. The entire polynomial
\[
 Z_F(T)=\sum_{x:F(x)=1}T^{E(x)}
\]
is computable coefficientwise in $\operatorname{poly}(n+m+M)2^{O(n(\log\log n)^2/\sqrt{\log n})}$ time and polynomial space. Hence minimum/maximum feasible energy and the numbers of optimizers are computable within that bound. For Hamming weight $E(x)=\sum_i x_i$, a separate linear-size addition-tree encoding improves the exponent to $O(n\log\log n/\sqrt{\log n})$. In particular, the full independence polynomial of an arbitrary graph has that sharper bound in its number of vertices.
\end{proposition}
\paragraph{Exact inputs.}
E: The Exponential Time Hypothesis is False, supplied LaTeX. Proposition 3.4; Lemmas 4.7--4.8, 5.1--5.2, 6.1--6.3; Appendix B; Theorem 1.1.

NN-C001. compression preserves the full solution set, not merely its cardinality.

\begin{proof}
Apply E's disjoint sparsification lemma with a rational parameter of reciprocal polylogarithmic size, for example $\sigma=1/\lceil\sqrt{\log_2n}\rceil$. For fixed $k$, the leaves have $n(\log n)^{O_k(1)}$ clauses on the original variables, their satisfying sets partition that of $F$, and their number is at most a polynomial times $2^{\sigma n}$. Compress each leaf as in NN-C001. The final equations $P_i=0$ have count $a=O(n)$, fixed degree and polylogarithmic sparsity, and their Boolean zero set is exactly that leaf's satisfying set.

Set $\ell=n+1$ and work in the finite-dimensional commutative $\mathbb Q$-algebra
\[
 \mathcal A=\mathbb Q[V,T]/(V^\ell,T^{M+1}).
\]
For each $P_i$ use the parity-lift factor from E Lemmas 5.1--5.2,
\[
 B_i(x,V)=\sum_{j<\ell}e_j(L_i(x))\,[L_i(x)(L_i(x)-1)]^jV^j,
\]
where $L_i$ is its integer monomial lift. For a local cost table $c_h$ insert the multilinear interpolation factor $W_h(x,T)$ whose value on a bit assignment is exactly $T^{c_h(x)}$. Its input-variable degree is at most $w$, regardless of the exponent of $T$. The product of parity factors has incidence degree at most $k(a+2(\ell-1))=O(n)$ after the global $V$ truncation, and the cost factors add at most $wCn=O(n)$ incidence degree. If $M_{\rm inc}$ is the total scope size, then $M_{\rm inc}=n(\log n)^{O(1)}$. Set
\[
 b_0=\max\{1,\lceil M_{\rm inc}/n\rceil\}\lceil\sqrt{\log_2n}\rceil.
\]
Branch on all variables whose incidence exceeds $b_0$. There are at most $M_{\rm inc}/b_0\leq n/\sqrt{\log_2 n}$ such variables, and the remaining incidence is polylogarithmic. Lemmas 4.7--4.8 and Proposition 3.4 then evaluate the sum of the full product in the asserted weaker-exponent time.

After evaluating on Boolean assignments, apply the coefficient-sum map in $V$ only. This map is used after multiplication and is not treated as a ring homomorphism. Every dropped term of $V$-order at least $\ell$ is divisible by $2^\ell$ at integer arguments. Consequently the resulting integer polynomial in $T$ is coefficientwise congruent modulo $2^\ell$ to the leaf energy enumerator. Each genuine coefficient is between zero and $2^n$, so the least nonnegative residues recover it exactly. The $T$ truncation loses nothing because the sum of the local maximum energies is at most $M$. The algebra dimension is $(n+1)(M+1)$, polynomial in the stated parameters. E's rational height argument applies coordinatewise: only polynomially many bits are needed, and algebra operations have polynomial overhead. Process leaves and branches sequentially, adding their recovered coefficient arrays. Disjointness proves the final polynomial is exactly $Z_F$ and keeps every final coefficient at most $2^n$.

For the sharper Hamming specialization, form a balanced tree adding the $n$ one-bit inputs. At level $j$ there are $O(n/2^j)$ additions of $O(j)$-bit integers, using $O(j)$ Boolean gates per addition. The total gate count is $O(n\sum_{j\geq1}j/2^j)=O(n)$. Equality of the final binary sum to a chosen $h\in\{0,\ldots,n\}$ uses $O(\log n)$ further gates. Its Tseitin clauses have a unique extension for every original assignment and introduce $O(n)$ variables. Conjoin these clauses with $F$, obtaining fixed width $\max(k,3)$ and $O(n)$ variables, even when $F$ is dense. E Theorem 1.1 computes the coefficient of $T^h$ with the sharper general-density exponent. Repeat for $n+1$ values of $h$. For a graph use $F=\bigwedge_{uv\in E}(\neg x_u\vee\neg x_v)$; the coefficient counts independent sets of each cardinality. The polynomial input prefactor includes the edge list.
\end{proof}
\paragraph{Status and complexity impact.}
Proved coefficient-algebra extension; Hamming specialization separately sharpened by a parsimonious linear-size circuit. Independent complete deduction-proof review by root: passed; imported input theorems are not claimed independently proved. Gives complete polynomial-size distributions and exact optimization, not only a single unweighted count. Polynomially bounded total energy and $O(n)$ bounded-scope cost tables are essential. Arbitrarily many cost factors or exponentially many output coefficients fall outside the bound. The independence polynomial can be evaluated at any supplied rational activity with polynomial extra bit cost; this uses its $n+1$ coefficients and does not assert a general arbitrary-weight shortcut for dense CSPs.

\paragraph{Geometric and representation dictionary.}
On the finite Boolean solution scheme, the energy partitions the rational components into clopen subsets $X_j$. The enumerator is $\sum_j\operatorname{rank}H^0(X_j,\mathcal O)T^j$, equivalently the character of the graded vector space with one basis vector per solution in its energy degree. The finite-algebra trace is immediate in the component idempotent basis; SGA4 XVII 6.2.3 supplies the verified finite-trace formalism. The variable $T$ records this chosen grading; it does not turn $T^j$ into a Tate twist or a geometric Frobenius eigenvalue without a separate realization. Thus the result produces an effective graded rank computation but no plethysm multiplicity formula or GCT obstruction.

\subsection{NN-C014 --- Two unconditional limits on universal hitting-set interpretations of E}
\begin{proposition}[NN-C014]
A set $H\subseteq\mathbb F_2^n$ meeting the nonempty zero set of every system of $n$ sparse linear Boolean equations must equal the entire cube. Also, a fixed $a\times m$ binary matrix with $a<m$ cannot preserve the Boolean zero set of every $m$-tuple of sparse linear polynomials by replacing $f$ with $Af$. Therefore E's small collection of parity checks is necessarily instance dependent; its theorem is not a universal subexponential assignment hitting set or a universal rank-deficient linear sketch of violations.
\end{proposition}
\paragraph{Exact inputs.}
Elementary singleton and kernel arguments. complete proof below.

E: The Exponential Time Hypothesis is False, supplied LaTeX. Section 6, Lemmas 6.1--6.2, instance-dependent greedy completion.

\begin{proof}
For each $b=(b_1,\ldots,b_n)\in\mathbb F_2^n$, take the $n$ equations $x_i+b_i=0$. Their common zero set is exactly $\{b\}$, and each equation has degree one and at most two monomials. If $H$ meets every such nonempty set then $b\in H$ for every $b$, so $H=\mathbb F_2^n$. This is an information-theoretic obstruction to the specified point-query hitting-set model; it says nothing against algorithms that manipulate the supplied equations.

For the second assertion, let $A$ have fewer rows than columns. Choose a nonzero vector $z\in\ker A$. For $n\geq1$ set $f_i(x)=z_i x_1$. Every $f_i$ has degree at most one and at most one monomial. The original system vanishes exactly on $x_1=0$, since some $z_i=1$. Yet $Af(x)=(Az)x_1=0$ identically, so the compressed system accepts the whole cube. Thus no such fixed matrix preserves zero sets for all instances. Equivalently a universal linear sketch separating every possible nonzero violation vector must have trivial kernel and at least $m$ rows.

E's Lemma 6.1 is consistent with this obstruction: its universal initial checks exclude only vectors of small support. Lemma 6.2 then selects additional checks using the particular image $\{f(x):x\in\mathbb F_2^n\}$ and exact survivor counts. The completed matrix only needs its kernel to avoid the nonzero vectors in that image. NN-C001 retains precisely this dependence and computational cost.
\end{proof}
\paragraph{Status and complexity impact.}
Proved structural limitation. Independent complete deduction-proof review by root: passed; imported input theorems are not claimed independently proved. Separates E's conditional-expectation construction from a black-box SAT hitting set and from formal PIT hitting sets. A uniform efficient constructor depending on the instance is a different question; no P versus NP lower bound is inferred from the two elementary impossibility statements.

\paragraph{Geometric and representation dictionary.}
Each singleton system defines one rational reduced point, a proper finite etale scheme of degree one. A universal set of evaluation points must meet every such closed subscheme of the Boolean scheme, forcing every component to occur. In its coordinate algebra $\mathbb F_2^{2^n}$, these are the primitive coordinate idempotents. The sketch counterexample changes the ideal and doubles the selected clopen locus. Finite trace and degree are as in SGA4 XVII 6.2.3; weight-zero purity does not remove the need to retain each possible component in a universal evaluation set.

\subsection{NN-C015 --- The constant-free model removes the analytic constant-elimination gate}
\begin{proposition}[NN-C015]
For constant-free integer arithmetic circuits with polynomial size and polynomial formal degree, the implication
\[
 \mathrm{VP}_{\mathbb Z}^{0}=\mathrm{VNP}_{\mathbb Z}^{0}
 \quad\Longrightarrow\quad \#\mathrm P\subseteq\mathrm{FP}/\mathrm{poly}
\]
holds without A or M. Thus the arbitrary-constant characteristic-zero transfer in NN-C010 and the constant-free alternative of CGGM26 have different analytic requirements. A separation for these constant-free classes does not itself establish $\mathrm{VP}_{\mathbb C}\ne\mathrm{VNP}_{\mathbb C}$. This is a classical model refinement, with a direct proof supplied here, not a claimed new theorem from E.
\end{proposition}
\paragraph{Exact inputs.}
Definitions of VP\textasciicircum{}0\_Z and VNP\textasciicircum{}0\_Z. CGGM26 p.5: constant-free circuits of polynomial size and formal degree.

Chukhin--Gajulapalli--Golovnev--Mihajlin, A Little Derandomization Goes a Long Way, ECCC TR26-313, 5 Nov 2026. Proposition 8.1 p.40, coefficient-norm bound; Theorem 1 p.5.

Burgisser, Completeness and Reduction in Algebraic Complexity Theory, Springer 2000, Chapter 4. Theorem 4.5(A), Corollary 4.6(1), pp.64--65; Theorem 4.12 and Remark 4.17, pp.71--75; Theorem 4.18 and Corollary 4.19, p.76; proof of 4.5(A3), pp.77--78.

\begin{proof}
Let $R(x,w)$ be any polynomial-time witness verifier defining a $\#\mathrm P$ function. For each input length, build a polynomial-size AND/NOT verifier circuit and introduce a Boolean variable for each gate. Over $\mathbb Z$, the gate-validity indicator for $z=ab$ is
\[
 zab+(1-z)(1-ab),
\]
and that for $z=1-a$ is $z(1-a)+(1-z)a$. Each is a constant-free bounded-degree polynomial, equal to one exactly when the corresponding Boolean gate is correct. Multiply all these indicators and the output bit. The resulting polynomial $Q(x,w,z)$ has a constant-free circuit of polynomial size and formal degree. For each $x,w$ there is exactly one consistent assignment to the gate variables, and its product is one precisely when $R(x,w)$ accepts. Hence $\sum_{w,z\in\{0,1\}^{\operatorname{poly}(n)}}Q(x,w,z)$ belongs to $\mathrm{VNP}_{\mathbb Z}^{0}$ and equals the desired count on Boolean $x$.

Assume equality of the two constant-free classes. Its counting polynomial has a constant-free arithmetic circuit with $s=\operatorname{poly}(n)$ gates and formal degree $d=\operatorname{poly}(n)$. Here constants have formal degree one, multiplication adds formal degrees, and addition or subtraction takes their maximum. Write $\|P\|_1$ for the sum of absolute coefficient values. A direct induction in topological gate order gives $\|P_g\|_1\leq2^{t d_g}$ at gate number $t$: addition increases the larger norm by at most a factor two, while multiplication multiplies the two bounds and adds formal degrees. The leaves $0,1,-1$ and variables satisfy the base case. Thus every used gate has coefficient norm at most $2^{sd}$. At Boolean inputs its integer value has at most $sd+O(1)$ bits. Ordinary integer addition and schoolbook multiplication evaluate the whole circuit in polynomial time. The circuit description is polynomial-length advice, proving the functional containment.

The formal-degree restriction matters for this characteristic-zero height argument: repeatedly squaring the constant two produces $2^{2^s}$ from $O(s)$ gates despite ordinary polynomial degree zero. Its formal degree grows exponentially and is excluded from that argument. NN-C020 gives a different modular simulation of bounded integer outputs which does not require this restriction. Finally, constant-free polynomial-formal-degree circuits form a more restrictive model than arithmetic circuits with arbitrary complex constants. Hardness against the former does not imply hardness against the latter. The proof above uses explicit integer arithmetic and never chooses a prime or a residue-field specialization, so it has no hidden GRH assumption.
\end{proof}
\paragraph{Status and complexity impact.}
Classical model refinement proved directly. Independent complete deduction-proof review by root: passed; imported input theorems are not claimed independently proved. Explains the restricted transfer underlying the current unconditional win--win. It does not convert E's near-linear $\mathrm E^{\mathrm{NP}}$ lower bound into the required nonuniform counting lower bound. It also does not supply a converse from ordinary algebraic separation to Boolean separation.

\paragraph{Geometric and representation dictionary.}
The verifier product gives a finite Boolean scheme together with its acceptance idempotent; summing the idempotent is the finite-algebra trace. Over $\mathbb Z$, the supplied integer presentation and formal-degree bound make all intermediate coefficients effectively bounded. The trace dictionary follows directly by using the component idempotent basis, as in NN-C003; compare SGA4 XVII 6.2.3. By contrast NN-C010 starts with arbitrary complex constants, whose replacement by finite-field data uses number-field points and Dedekind zeta estimates. These are different effective presentations of related rank/count constructions, not different purity statements.

\subsection{NN-C016 --- A genuine symbolic transfer for bounded-arity partition polynomials}
\begin{proposition}[NN-C016]
Fix arity $w$ and factor density $C$. For a factor hypergraph with $n$ Boolean variables and $m\leq Cn$ factors of scope at most $w$, attach an independent formal weight $y_{f,a}$ to every local assignment $a$ of factor $f$. Its full partition polynomial
\[
 Z_H(y)=\sum_{x\in\{0,1\}^n}\prod_f y_{f,x|I_f}
\]
has a deterministically constructible arithmetic circuit over $\mathbb Q$ of size
\[
 \operatorname{poly}(n)\,2^{O_{w,C}(n(\log\log n)^2/\sqrt{\log n})}.
\]
The rational constants have polynomial bit length. In particular the symbolic perfect-matching polynomial of a graph of bounded maximum degree has subexponential arithmetic circuit size in its vertex/edge parameter. This is a restricted symbolic consequence of E's full common routine; it is stronger than what follows merely by regarding E as a numerical SAT counter.
\end{proposition}
\paragraph{Exact inputs.}
E: The Exponential Time Hypothesis is False, supplied LaTeX. Proposition 3.4, especially algebra-operation count and rational Hermite interpolation; Lemmas 4.7--4.8; Section 3.1 product interface and measurement-separation proof.

\begin{proof}
Use the coefficient algebra $\mathcal A=\mathbb Q[y_{f,a}]$. For each factor take its Boolean multilinear interpolation
\[
 \Phi_f(x_{I_f})=\sum_{a\in\{0,1\}^{I_f}}y_{f,a}
       \prod_{i:a_i=1}x_i\prod_{i:a_i=0}(1-x_i).
\]
It is linear in the formal weights and has degree at most $w$ in its Boolean inputs. The sum-product is exactly $Z_H(y)$, and its independent-incidence input degree is at most $wm=O(n)$. The number of incidences is also $O(n)$. Branch on variables of incidence above $\max(1,\lceil M_{\rm inc}/n\rceil)\lceil\sqrt{\log_2n}\rceil$, as in NN-C003, and retain the original scopes even if a substitution makes some coefficient zero. There are at most $n/\sqrt{\log_2n}$ branched bits and every resulting layout satisfies Lemmas 4.7--4.8.

Proposition 3.4 explicitly allows every commutative $\mathbb Q$-algebra, counts polynomially many operations in that algebra per enumerated vector, and permits final interpolation by the inverse of a rational Hermite matrix. Implement an algebra element as a pointer to a gate in a shared arithmetic directed acyclic graph. Each addition, subtraction or multiplication in $\mathcal A$ creates one arithmetic gate; truncated Taylor arrays have polynomial length and require polynomially many such gates per operation. Factor evaluation takes $O(2^w)$ times a polynomial number of algebra operations because $w$ is fixed. For interpolation compute the inverse matrix over $\mathbb Q$ once and use its entries as rational constants, followed by linear combinations of the symbolic Taylor data. Do not run fraction-free elimination with symbolic pivot entries: no division by an element involving $y$ is needed.

All control choices are independent of symbolic coefficients. The layout construction uses only the scopes. The measurement-separation scans select representatives by their integer measurements and fixed input bits, not by comparing values in $\mathcal A$. Branching fixes declared Boolean variables and does not ask whether a symbolic factor is zero. Therefore the complete computation can be unrolled into a circuit without a symbolic identity oracle. Proposition 3.4 and Lemma 4.8 bound the number of operations by the displayed exponential times a polynomial, including all high-incidence branches. Layout and Hermite constants have polynomial bit length by E's row-coefficient and Appendix B bounds. Circuit construction time is polynomial in its output size; storing algebra elements as gate references avoids expansion into monomials. The output polynomial has degree $m$ in the weights, even though it may have exponentially many monomials.

For a bounded-degree graph, put one bit $x_e$ on each edge. At every vertex insert the fixed-arity indicator that exactly one incident edge is selected. On each edge insert the unary factor $(1-x_e)+y_e x_e$. The product is zero off perfect matchings and equals $\prod_{e\in M}y_e$ on a perfect matching. With maximum degree fixed, the number of bits, factors and incidences is $O(v+e)$; isolated vertices immediately give the zero polynomial when they must be matched. Thus this is a specialization of the partition construction, preserving circuit size. In the bipartite case it is the permanent polynomial with the graph's forbidden entries set to zero.
\end{proof}
\paragraph{Status and complexity impact.}
Proved symbolic extension from E Proposition 3.4. Independent complete deduction-proof review by root: passed; imported input theorems are not claimed independently proved. This corrects an overly broad counting-only objection in this restricted setting. It does not give a polynomial-size circuit or VP=VNP. A dense $r\times r$ permanent encoded by edge bits has $\Theta(r^2)$ Boolean variables, so substituting that size into the result gives no useful subexponential bound in $r$ (and is worse than the usual $2^{O(r)}$ upper bound). Compression through parity lifts does not automatically preserve this characteristic-zero symbolic argument: coefficientwise reduction modulo a power of two is not an arithmetic-circuit operation over $\mathbb Q$. This original construction uses rational constants. NN-C033 subsequently proves a modified construction over every prime field, and NN-C034 gives integral constant-free circuits of polynomial formal degree over Z itself. These are strengthened constructions; the original rational circuit is not simply reduced modulo a denominator prime.

\paragraph{Geometric and representation dictionary.}
The partition polynomial is the trace of multiplication by $\prod_f\Phi_f$ on the finite free algebra $\mathbb Q[y][x]/(x_i^2-x_i)$ over the parameter scheme $\operatorname{Spec}\mathbb Q[y]$. In the orthogonal idempotent basis indexed by bit assignments, this multiplication matrix is diagonal with the corresponding factor products on the diagonal. Its trace is their sum. Applying a coefficient-ring homomorphism entrywise commutes with summing the diagonal, proving the claimed base-change compatibility directly. The original EGA IV4, printed p.111,18.2.1--18.2.2 records this trace, with base change in equation18.2.1.2; compare the relative-dimension-zero formalism in SGA4 XVII 6.2.3. Specializing formal weights commutes with the trace and with the constructed arithmetic circuit. This is an actual family-level symbolic computation, unlike a single numerical rank. It still does not provide orbit-closure multiplicity obstructions, border lower bounds, or a finite Weil realization of nonstandard canonical bases.

\subsection{NN-C017 --- Sparse equations over every fixed finite field by prime-adic total-order truncation}
\begin{proposition}[NN-C017]
Fix a finite field $K=\mathbb F_{p^r}$ and constants $k,a,b$. Given $n$ variables over $K$ and at most $n(\log n)^a$ polynomial equations of degree at most $k$, with at most $(\log n)^b$ listed monomials per equation, their exact number of common zeros is computable deterministically in
\[
 \operatorname{poly}(L+n)\,
 2^{O_{K,k,a,b}(n(\log\log n)^2/\sqrt{\log n})}
\]
bit operations and polynomial space. Here $L$ is the explicit input length. Witness search and exact uniform sampling from a nonempty zero set have the same bound, with expected time for sampling. This extends NN-C002 to odd characteristic at the weaker Proposition 3.4 exponent; it does not assert the sharper Appendix D bound there.
\end{proposition}
\paragraph{Exact inputs.}
E: The Exponential Time Hypothesis is False, supplied LaTeX. Proposition 3.4; Lemmas 4.7--4.8; Lemmas 5.1--5.2; Proposition 5.3 and Appendix B; complete proofs of Lemmas 6.1--6.2.

NN-C003. explicit common-denominator and Hermite-minor bit-cost method.

NN-C002. prefix self-reduction and exact rejection sampling.

\begin{proof}
We first count common zeros of $m\leq CN$ sparse bounded-degree polynomials in $N$ Boolean variables with coefficients in the fixed prime field $\mathbb F_p$. Choose integer representatives of their coefficients and let $L_i(x)$ be the resulting integer polynomials. On Boolean inputs they reduce to the prescribed polynomials modulo $p$. Set $Y_i=L_i^{p-1}$. Fermat's elementary identity gives $Y_i(x)\equiv0$ or $1\pmod p$, with the first case exactly when $L_i(x)\equiv0\pmod p$. For completeness, for $u\not\equiv0\pmod p$, multiplication by $u$ permutes the $p-1$ nonzero residues; comparing their products and cancelling $(p-1)!$ gives $u^{p-1}=1$ in $\mathbb F_p$.

Set $\ell=N+1$ and use exactly E Lemma 5.1's integer polynomial
\[
 p_\ell(Y)=(1-Y)^\ell\sum_{j=0}^{\ell-1}\binom{\ell+j-1}{j}Y^j
           =\sum_{j=0}^{\ell-1}e_j(Y)h(Y)^j,
 \qquad h(Y)=Y(Y-1),\quad\deg e_j\leq1.
\]
The proof of that lemma establishes the polynomial congruences $p_\ell\equiv1\pmod{Y^\ell}$ and $p_\ell\equiv0\pmod{(1-Y)^\ell}$, independently of characteristic. Therefore $p_\ell(Y_i(x))$ is the zero indicator modulo $p^\ell$, and $p\mid h(Y_i(x))$. In $\mathcal A=\mathbb Q[V]/(V^\ell)$ use
\[
 B_i(x,V)=\sum_{j<\ell}e_j(Y_i(x))h(Y_i(x))^jV^j.
\]
After multiplying all factors, sum the retained $V$ coefficients. Each discarded term has total $h$-order at least $\ell$ and is divisible by $p^\ell$ at every Boolean input. Thus the sum-product recovers the number of simultaneous zeros modulo $p^\ell$. Since $p^\ell>2^N$, its least nonnegative residue is the exact count. This readout is performed after multiplication, not as a ring homomorphism.

If the original degree is at most $d$, then $\deg Y_i\leq d(p-1)$. Even with independent input copies for every factor, every retained coefficient has degree at most
\[
 d(p-1)\bigl(m+2(\ell-1)\bigr)=O_{p,d,C}(N).
\]
Scopes remain polylogarithmic and total incidence is $N(\log N)^{O(1)}$. Add neutral unary factors for unused declared bits. If total incidence is $M$, branch on bits appearing in more than
$\max(1,\lceil M/N\rceil)\lceil\sqrt{\log_2N}\rceil$ scopes. At most $N/\sqrt{\log_2N}$ bits are branched; every branch retains $\Theta(N)$ declared bits and has polylogarithmic maximum incidence. E Lemmas 4.7--4.8 and Proposition 3.4 now apply. Factor evaluation computes the sparse $L_i$, its fixed power $p-1$, and successively the $h$ powers in the truncated algebra; it uses polynomially many rational operations and never enumerates a factor's entire scope.

The bit bound is explicit. With E's row denominator $d_{\rm den}=(R-1)!Q!$, a substituted $Y_i$ has denominator dividing $\Delta=d_{\rm den}^{d(p-1)}$. Each $B_i$ has denominator dividing $\Delta^{2\ell-1}$ and their product has denominator dividing $\Omega=\Delta^{m(2\ell-1)}$. Coefficient bits are polynomial by E Lemma 5.1 and Appendix B; the additional fixed power $p-1$ changes only constants. Clearing $\Omega$ bounds every Taylor datum by a polynomial-bit integer: a product has polynomially many polynomial-bit factors and summing the enumerated terms adds only the logarithm of their number. The integer Hermite matrix and its fraction-free minor bound are exactly as in NN-C003. Final integer reduction is modulo $p^\ell$, which has $O_p(N)$ bits. The time is the displayed weaker exponent and sequential evaluation gives polynomial space.

Next extend the equation count from $CN$ to $N(\log N)^{O(1)}$. E Lemma 6.1's graph, its shifts, its real-valued Walsh mixing estimate, and its shortening do not depend on the coefficient field. Keep its parameters
\[
 B=\max(1,\lceil m/N\rceil),\quad \mu=\lfloor\log_2m\rfloor+1,
 \quad P=20000B^2\mu^2,\quad D_0=320000B^2\mu^3.
\]
Use the same $d_G,q_G,c_G,\delta$ and replace only the binary local matrix by a uniformly random $c_G\times d_G$ matrix over $\mathbb F_p$. A fixed nonzero vector belongs to its kernel with probability $p^{-c_G}$. If $h=\lceil\delta d_G\rceil-1\geq1$, the number of nonzero vectors of weight at most $h$ is at most
\[
 \sum_{j=1}^h\binom{d_G}{j}(p-1)^j
 \leq(h+1)(p d_G)^h
 \leq(p d_G)^{2h}
 <p^{d_G/(25B)}.
\]
Indeed $d_G\leq m<2^\mu$, so $\log_p(p d_G)\leq1+\mu\leq2\mu$, while $h<d_G/(100B\mu)$. Since $c_G\geq d_G/(8B)$, the union bound is below one. If $h=0$ there is nothing to exclude. Exhaustive enumeration of local matrices and the low-weight vectors costs $\exp(\operatorname{poly}(B,\log(m+2)))$ for fixed $p$. The rest of Lemma 6.1 uses only supports and graph mixing, so it produces at most $3N$ rows, each of support at most $D_0$, whose nonzero kernel vectors have support at least $m/P$.

For the remaining greedy stage, use lists in $([m]\times\mathbb F_p)^P$. A list hits the support of a false survivor with probability at least $1-e^{-1}>1/2$. Conditional on a hit, one uniform coefficient makes the tested sum uniform in $\mathbb F_p$, even with repeated indices. Its rejection probability is consequently at least $(1-e^{-1})(1-1/p)>1/4$. A candidate minimizing the exact count leaves at most $3/4$ of the false survivors, exactly as in E Lemma 6.2. There are at most $2^N$ Boolean assignments, so $3N$ stages remove all false survivors. Every candidate has at most $6N$ equations and polylogarithmic monomial sparsity. The number of core calls is at most $3N(pm)^P=\exp((\log N)^{O(1)})$, absorbed by the claimed running time. All enumeration is sequential and uses polynomial space. This compression preserves the common zero set on the Boolean cube.

Finally encode the original $\mathbb F_{p^r}$ variables. Fix an $\mathbb F_p$ basis; each original variable has $r$ coordinates and each equation gives $r$ coordinate equations. A degree-$k$ monomial expands into at most a fixed multiple of $r^k$ coordinate monomials. Encode each prime-field coordinate by $t=\lceil\log_2p\rceil$ bits, interpreted as the integer $\sum_{j<t}2^jx_j$. Add one polynomial equation per coordinate saying that this bit code is below $p$: explicitly, take one minus the sum of the bit-pattern indicator monomials for $0,\ldots,p-1$. On bits it is zero exactly for valid codes. It has constant degree and at most $2^t$ expanded pattern terms times a fixed constant. There are $N=rtn=\Theta_K(n)$ bits, $rm+rn=N(\log N)^{O(1)}$ equations, fixed degree and polylogarithmic monomial sparsity. Valid codes are in bijection with the original field assignments. The Boolean prime-field counter therefore gives exactly the desired number. Prefix fixing adds only $O(N)$ linear equations, so NN-C002's search and exact sampling proofs apply unchanged.
\end{proof}
\paragraph{Status and complexity impact.}
Complete new deduction proved; complete self proof-check passed. Independent complete deduction-proof review by root: passed; imported input theorems are not claimed independently proved. This removes an actual algorithmic restriction: sparse systems over fixed odd-characteristic fields now have a deterministic mildly subexponential exact counter with the original variable parameter, even when their scopes are polylogarithmic. The modulus and field extension degree must remain fixed; their powers and coordinate expansion enter the constants and cannot be hidden if they grow. The stronger characteristic-two Appendix D exponent has not been transferred. Neither polynomial-time counting nor a uniform identity test follows.

\paragraph{Geometric and representation dictionary.}
After adjoining $x_i^{|K|}-x_i$, the original zero set is a finite reduced split scheme over $K$ whose rank is the count. The basis expansion and valid-code equations give a bijection of rational components with a finite Boolean scheme over $\mathbb F_p$; they are an effective coordinate presentation, not an identification of unrestricted affine zero schemes. The $p$-adic order in $h(Y)$ is a divisibility filtration on integer-valued indicators. The formal nilpotent $V$ records this filtration and is not a Frobenius coordinate or a Tate twist. On the resulting finite scheme, the component-idempotent basis makes counting the trace of the identity; SGA7 II XXII Corollary 3.2 supplies its mod-$p$ coherent trace expression and SGA4 XVII 6.2.3 the finite-trace formalism. The algorithm computes this succinctly presented rank without constructing the exponential basis and supplies no new representation-theoretic positivity statement.

\subsection{NN-C018 --- Sparse polynomial equations over every fixed finite commutative ring}
\begin{proposition}[NN-C018]
Let $R$ be any fixed finite commutative ring. Fix $k,a,b$. Exact counting, witness search, and exact uniform sampling for systems of at most $n(\log n)^a$ degree-at-most-$k$ polynomials over $R$, with at most $(\log n)^b$ monomials per equation, have the NN-C017 time bound
\[
 \operatorname{poly}(L+n)\,
 2^{O_{R,k,a,b}(n(\log\log n)^2/\sqrt{\log n})}
\]
and polynomial space. Sampling has expected running time. In a fixed $p$-primary ring, the same bound computes a weighted sum over the solutions when there are $O(n)$ bounded-arity factors whose rational table entries lie in a fixed finite alphabet. For a general $R$, this weighted assertion holds when the factors separate as products over its primary components. Arbitrary input-height rational weights and weights coupling distinct primary components are not included in this weighted assertion.
\end{proposition}
\paragraph{Exact inputs.}
NN-C017. complete prime-field Boolean-system counter and p-ary equation compression.

J.S. Milne, Group Theory, v4.01 (2025). Theorems1.54 and1.57(c), pp.25--27, finite abelian groups as products of prime-power cyclic groups.

E: The Exponential Time Hypothesis is False, supplied LaTeX. Proposition3.4; Lemmas4.7--4.8 and5.1--5.2; AppendixB.

\begin{proof}
The zero ring is trivial, so assume $R$ has at least two elements. Write its characteristic as $M=\prod_p p^{a_p}$. The elementary integer Chinese remainder construction gives integers $e_p$ which are one modulo $p^{a_p}$ and zero modulo every other prime-power divisor of $M$. Their images in $R$ are orthogonal central idempotents with sum one. Hence
\[
 R\longrightarrow\prod_p R_p,\qquad R_p=e_pR,\qquad
 x\longmapsto(e_px)_p
\]
is a ring isomorphism: its inverse adds the components, and the identities just stated check both compositions and multiplication. Every equation therefore holds exactly when all its component equations hold. Variables in distinct components are independent. Thus the solution set is the Cartesian product of the component solution sets, and the exact counts multiply. There are only a fixed number of components. It suffices to handle a fixed $p$-primary ring.

By the stated finite-abelian-group theorem, fix additive generators $b_1,\ldots,b_d$ of this ring with unique coordinates
\[
 x=\sum_{j=1}^d z_jb_j,
 \qquad0\leq z_j<p^{a_j},\qquad
 (R,+)\cong\bigoplus_{j=1}^d\mathbb Z/p^{a_j}\mathbb Z.
\]
The multiplication table $b_ib_j=\sum_t c_{ijt}b_t$, with fixed integer representatives for the $c_{ijt}$, is part of this fixed presentation. Expanding one degree-$k$ monomial in these coordinates introduces only a constant depending on $R,k$ many integer monomials, all of degree at most $k$. Constant coefficients in $R$ are handled by the same fixed multiplication table. Consequently a ring equation is equivalent to $d$ congruences
\[
 L_t(z)\equiv0\pmod{p^{a_t}},
\]
where $L_t$ is an integer polynomial of fixed degree and polylogarithmic monomial sparsity. Intermediate reductions modulo the coordinate moduli are unnecessary: the additive relations $p^{a_t}b_t=0$ ensure that the final coordinate congruence gives the correct element.

Encode each coordinate by $\lceil\log_2(p^{a_j})\rceil$ bits and impose its fixed-size validity table, as in NN-C017. This gives $N=\Theta_R(n)$ Boolean variables. Every coordinate polynomial after substitution still has fixed degree and polylogarithmic sparsity. We next replace divisibility by $p^{a_t}$ by a fixed number of genuinely sparse $\mathbb F_p$ polynomial equations, taking care not to invert a factorial divisible by $p$.

For every integer $u$ and $j\geq0$, define the integer-valued polynomial
\[
 D_j(u)=\binom{u}{p^j}.
\]
For fixed $a$ and $j<a$, this is periodic modulo $p$ with period $p^a$. Indeed the formal power series identity
$(1+T)^{u+p^a}=(1+T)^u(1+T)^{p^a}$ holds also for negative integer $u$, with integral binomial coefficients. Modulo $p$, the last factor is $1+T^{p^a}$, by repeated Frobenius. Thus coefficients below degree $p^a$ do not change. For $0\leq u<p^a$, write $u=\sum_{i<a}u_ip^i$ with $0\leq u_i<p$. Then
\[
 (1+T)^u\equiv\prod_{i<a}(1+T^{p^i})^{u_i}\pmod p.
\]
The coefficient of $T^{p^j}$ is $u_j$: the maximum total exponent available from the lower factors is $p^j-1$, and any higher factor has exponent greater than $p^j$. By periodicity, for every integer $u$ the residues $D_0(u),\ldots,D_{a-1}(u)$ are exactly the base-$p$ digits of $u\bmod p^a$. In particular,
\[
 p^a\mid u\quad\Longleftrightarrow\quad
 D_j(u)\equiv0\pmod p\quad(0\leq j<a).
\]

Apply this to each $L_t(x)$. As a rational polynomial,
$D_j(L_t)=\prod_{h=0}^{p^j-1}(L_t-h)/(p^j)!$ has fixed degree and polylogarithmic monomial sparsity: $p^j$ is a constant depending only on $R$. Expand it over $\mathbb Q$, replace every positive bit power $x_i^h$ by $x_i$, and collect like monomials. The resulting multilinear polynomial has \emph{integer} coefficients. To prove this, for any rational polynomial $Q$ taking integer values on the Boolean cube, its unique multilinear representative has coefficient
\[
 [x^S]\operatorname{ML}(Q)
   =\sum_{T\subseteq S}(-1)^{|S|-|T|}\,Q(\mathbf1_T),
\]
where coordinates outside $T$ are zero. This is the inversion of
$Q(\mathbf1_T)=\sum_{S\subseteq T}[x^S]\operatorname{ML}(Q)$ and is an integer. Here $Q=D_j(L_t)$ is integer-valued because $L_t$ has integer Boolean values. Computationally one simply expands and collects the sparse polynomial; the coefficient formula proves integrality and does not require enumerating all assignments on the polylogarithmic scope. All rational denominators before collection divide the fixed factorial, and coefficient bit lengths are polynomial in the input length. Reduce the resulting integer coefficients modulo $p$ only after this collection.

The final Boolean system has $N(\log N)^{O(1)}$ equations, fixed degree, and polylogarithmic monomial sparsity. The number of coordinate and digit equations per original equation is fixed. It includes the constant-size validity equations. NN-C017 applies and is parsimonious with respect to valid coordinate encodings. The primary-component counts multiply as above; their outputs have $O_R(n)$ bits. Self-reduction and sampling work componentwise, and independent uniform component samples give the uniform distribution on their Cartesian product. This proves the unweighted claims.

For the weighted $p$-primary assertion, first compress the constraint system by the \emph{unweighted} NN-C017 counter; the selected equations preserve its exact Boolean zero set. Let the fixed rational weight alphabet have a common positive denominator $D$, and let $A\geq1$ bound the absolute values of all the integer weights obtained by multiplying by $D$. If there are $f\leq Cn$ weight factors, the numerator of the weighted sum after this scaling has absolute value at most $2^N A^f$. Choose $\ell=O_{R,C,A}(N)$ so that
\[
 p^\ell>2^{N+1}A^f.
\]
Express each scaled table by its integer multilinear interpolation polynomial on its fixed number of coordinate bits. Fill invalid codes by zero; the validity constraints already exclude them. Multiply these weight factors by the constraint factors in the same algebra $\mathbb Q[V]/(V^\ell)$ used in NN-C017. At Boolean points every omitted order-$\ell$ term remains divisible by $p^\ell$ after multiplication by the integer weight. The total incidence degree is still $O(N)$: the compressed constraint contribution is $O(m+\ell)$ and the $O(n)$ weight factors each have bounded degree and scope. The added scopes preserve polylogarithmic total incidence. The same layout, coefficient-height, interpolation, and space proof applies. The residue has a unique representative in $(-p^\ell/2,p^\ell/2)$ within the stated magnitude bound; this is the desired signed numerator. Divide by $D^f$ to recover the rational weighted sum.

If the ring has several primary components and each weight factor is a product of factors in the individual components, the full weighted summand and the solution domain separate, so the total partition function is the product of the computed component partition functions. Without that hypothesis, the product-of-counts argument does not apply. Likewise a variable $B$-bit weight alphabet can require $\ell=\Theta(n+nB)$, so this proof cannot simply hide arbitrary weight height in a polynomial overhead while retaining the same exponent.
\end{proof}
\paragraph{Status and complexity impact.}
Complete deduction proved; complete self proof-check passed. Independent complete deduction-proof review by root: passed; imported input theorems are not claimed independently proved. Nilpotents and fixed composite characteristic do not by themselves obstruct the sparse-equation algorithm: the digit reduction keeps the number of variables, equation degree and monomial sparsity within the required fixed-parameter interface. This is a theorem for each fixed ring, including fixed prime-power residue rings, not an algorithm uniform in a growing modulus such as an RSA modulus. The weighted restrictions are substantive, and the full rational bounded-arity theorem NN-C003 remains a different interface.

\paragraph{Geometric and representation dictionary.}
A polynomial solution over a finite nonreduced ring is an $R$-point of the affine equation scheme, not in general a geometric point of a reduced scheme. The primary decomposition is a decomposition of $\operatorname{Spec}R$ into clopen components; its functor of points turns the equation solution set into the displayed Cartesian product. Additive coordinates and binomial digits then present this finite set by a finite reduced Boolean scheme over $\mathbb F_p$. That presentation preserves the set and its rank, not the nilpotent structure, tangent spaces or deformation theory of the original equation scheme. Counting is the trace of the identity on the product algebra indexed by these points, and weighted counting is the diagonal multiplication trace; these assertions follow directly in the idempotent basis, with SGA4 XVII6.2.3 as the finite-trace dictionary. SGA7 II XXII Corollary3.2 applies to the resulting finite scheme for its mod-$p$ trace, not automatically to the original nonreduced ring-valued functor. No canonical-basis or orbit-closure positivity is inferred.

\subsection{NN-C019 --- E constructs one polynomial-height isolating weight vector for each instance}
\begin{proposition}[NN-C019]
Let $\mathcal F\subseteq\{0,1\}^n$ be a witness set whose prefix-restricted counts admit one of the exact algorithms NN-C001, NN-C002 or NN-C017. If $\mathcal F$ is nonempty, a deterministic algorithm in the same time bound up to a polynomial factor constructs a vector
\[
 w\in\{1,n+1\}^n\subseteq\{1,\ldots,2n\}^n
\]
and a unique minimum-weight witness $x^*\in\mathcal F$, with minimum weight at most $n$. It also identifies that witness. For fixed-width CNF, E Theorem1.1 gives $\operatorname{poly}(n+m)2^{O_k(n\log\log n/\sqrt{\log n})}$ time and polynomial space. For a fixed-density CNF, the construction retains the sharper $2^{O_{k,C}(n/\sqrt{\log n})}$ exponent.
\end{proposition}
\paragraph{Exact inputs.}
NN-C002. complete exact prefix-count self-reduction.

NN-C017. fixed-field extension and prefix constraints.

E: The Exponential Time Hypothesis is False, supplied LaTeX. Theorem1.1, including fixed-density statement.

Mulmuley--Vazirani--Vazirani, Matching is as Easy as Matrix Inversion, STOC1987. Section3, Lemma1, printed p.347; compare Lemma2, p.348 for the unique-lowest-term valuation argument.

\begin{proof}
Use exact prefix counts, with the bit order $0<1$, to find the lexicographically first witness $x^*$. Let $S=\{i:x_i^*=1\}$. No proper subset of $S$ is the support of a witness: its bit string would agree until its first removed one, where it has zero and $x^*$ has one, contradicting lexicographic minimality. Define $w_i=1$ for $i\in S$ and $w_i=n+1$ for $i\notin S$.

The weight of $x^*$ is $|S|\leq n$. If another witness has support $T\subseteq S$, then $T\ne S$ would contradict the preceding inclusion minimality; equality gives the same bit string. Every other witness therefore has an element outside $S$ and has weight at least $n+1>|S|$. Thus $x^*$ is the unique minimum. The argument also covers $S=\varnothing$. If the initial count is zero, the algorithm reports emptiness instead of claiming an isolated witness.

There are at most $n+1$ counting calls. In the sparse-system interfaces each prefix adds at most $n$ linear equations, with fixed degree and one or two monomials, so the corresponding bounds are preserved. For fixed-width CNF, a prefix consists of at most $n$ unit clauses; padding a unit clause by repeated literals keeps it within the width bound, or one substitutes its bit directly. The number of variables never increases, and an initial bound $m\leq Cn$ becomes at most $(C+1)n$ clauses if all variables are retained. Thus E's fixed-density exponent $2^{O_{k,C}(n/\sqrt{\log n})}$ is preserved as well. All counts, the current prefix, and the final weights use polynomial space.

MVV Lemma1 states that independent uniform weights in $[2n]$ isolate any nonempty family with probability at least $1/2$. The construction here selects a weight in that same range deterministically by first solving the instance at E's running time. It is not a polynomial-time replacement for the MVV random choice, and it is not one small weight family that works uniformly for all succinctly represented witness sets. It also does not eliminate signed cancellation at an unweighted specialization: isolating a term in a signed expression produces a nonzero leading term of a new weight-generating polynomial, whereas the original specialization can still vanish.
\end{proof}
\paragraph{Status and complexity impact.}
Complete deduction proved; complete self proof-check passed. Independent complete deduction-proof review by root: passed; imported input theorems are not claimed independently proved. This makes the available deterministic NP-style isolation precise: instance-dependent isolation and witness recovery are mildly subexponential for the stated sparse interfaces, with polynomially bounded weights. A polynomial-size output does not imply polynomial construction time. It supplies neither $\mathrm{NP}=\mathrm{UP}$ nor $\mathrm{NP}\subseteq\oplus\mathrm P$ with polynomial uniformity. Hoza's logarithmic-seed layered-graph construction remains a different polynomial-time uniform theorem; NN-C007 states the missing NP transfer gates.

\paragraph{Geometric and representation dictionary.}
The finite witness scheme has a diagonal grading operator with one basis vector of degree $w\cdot x$ for every witness. Its character is $\sum_{x\in\mathcal F}T^{w\cdot x}$, and the least-degree graded piece has rank one. Equivalently, $w$ exposes a single vertex of the convex hull of the witness bit vectors and selects one monomial in the initial form of the positive generating polynomial. The count/character interpretation follows directly from the component-idempotent basis and the finite-trace formalism of SGA4 XVII6.2.3. These are finite-set and torus-grading statements; they do not construct a GCT canonical basis or make a signed structure-constant formula positive. A numerical signed specialization forgets the grading and can reintroduce cancellation.

\subsection{NN-C020 --- Unconditional Boolean simulation from a small explicit coefficient algebra}
\begin{proposition}[NN-C020]
Let $C_n$ be a division-free arithmetic circuit of polynomial size over $B_n\otimes_{\mathbb Z}\mathbb Q$, where $B_n$ is a commutative finite free $\mathbb Z$-algebra supplied by a basis $e_0=1,e_1,\ldots,e_{r(n)-1}$ of polynomial rank and polynomial-bit integer multiplication constants. Suppose every circuit constant is supplied by polynomial-bit rational coordinates in this basis, and the total circuit and algebra description has polynomial size. Assume that on every Boolean input the output is $g_n(x)e_0$, where $g_n(x)$ is an integer in $[0,2^{w(n)})$ for a polynomial $w$. Then $g\in\mathrm{FP}/\mathrm{poly}$, without GRH and without a formal-degree bound. If these descriptions and the output bound are constructible in polynomial time, then $g\in\mathrm{FP}$.

In particular, the conclusion applies when constants lie in a number field of polynomial degree supplied by an integral primitive element with polynomial-bit monic minimal polynomial and polynomial-bit rational power-basis coordinates. A collapse giving such circuits for all $\#\mathrm P$ counting families implies $\#\mathrm P\subseteq\mathrm{FP}/\mathrm{poly}$ unconditionally.
\end{proposition}
\paragraph{Exact inputs.}
Elementary quotient-algebra simulation. complete proof below.

NN-C015. constant-free counting-family construction; coefficient-height proof refined here.

Burgisser, Completeness and Reduction in Algebraic Complexity Theory, Springer 2000, Chapter 4. Theorem 4.5(A), Corollary 4.6(1), pp.64--65; Theorem 4.12 and Remark 4.17, pp.71--75; Theorem 4.18 and Corollary 4.19, p.76; proof of 4.5(A3), pp.77--78.

EGA IV, fourth part, Publications Mathematiques IHES32 (1967). Printed p.111,18.2.1--18.2.2; base-change equation18.2.1.2; primary PDF page110 (zero-based109).

\begin{proof}
Write each constant as $d_g^{-1}\sum_j a_{g,j}e_j$, with $d_g$ a positive integer. Put $D=\prod_g d_g$. This product has polynomial bit length, because the total constant description does. Every circuit value lies in the localization $B_n[D^{-1}]$. Choose the integer modulus
\[
 M=D\,2^{w(n)+1}+1.
\]
Its bit length is polynomial, it exceeds $2^{w(n)}$, and it is coprime to every $d_g$. In fact, an explicit inverse of $d_g$ modulo $M$ is
\[
 -\frac{D}{d_g}\,2^{w(n)+1}\pmod M,
\]
because multiplying this expression by $d_g$ gives $1-M$.

Therefore the canonical map $B_n\to B_n\otimes_{\mathbb Z}\mathbb Z/M\mathbb Z$ extends to a ring homomorphism
\[
 B_n[D^{-1}]\longrightarrow
 B_n\otimes_{\mathbb Z}\mathbb Z/M\mathbb Z.
\]
Evaluate the full circuit in this quotient algebra. Each stored value has $r(n)$ coordinates modulo $M$. Addition is coordinatewise and multiplication uses the supplied table; a naive multiplication needs only $O(r(n)^3)$ modular scalar operations. All integers have $O(\log M)$ bits. Consequently the whole evaluation uses polynomially many bit operations and polynomial space, independent of the magnitude of unreduced intermediate characteristic-zero values. Exponentially large formal degrees cause no difficulty: repeated squaring remains a sequence of the originally supplied polynomially many modular gates.

By the homomorphism property, the final vector is $(g_n(x)\bmod M,0,\ldots,0)$. The algebra remains free with the supplied basis after base change, so this coordinate statement is literal, not merely equality in an unspecified residue field. Since $0\leq g_n(x)<2^{w(n)}<M$, the least nonnegative representative of the first coordinate is the exact desired integer. This proves the simulation. Polynomial descriptions serve as advice in the nonuniform case; in the uniform case they are constructed before evaluation in polynomial time.

For the number-field specialization, let $f_n\in\mathbb Z[T]$ be the supplied monic minimal polynomial of the integral primitive element $\alpha_n$. Then $B_n=\mathbb Z[\alpha_n]\cong\mathbb Z[T]/(f_n)$ is free in the power basis of size $\deg f_n$. Reducing products of two basis powers by the monic relation constructs an integral multiplication table of polynomial bit length: at most $O(\deg f_n)$ successive reductions are needed, and each can increase the coefficient bit bound by at most a polynomial in the degree and the bit height of $f_n$. The given rational constant coordinates lie in $B_n[D^{-1}]$. The general argument applies even if $f_n$ becomes reducible or inseparable modulo $M$; neither a field quotient nor a root modulo $M$ is required.

Applying the simulation to the VNP counting polynomials constructed in NN-C015 proves its $\#\mathrm P$ consequence directly for this larger effective coefficient model. In particular, constant-free circuits use $B_n=\mathbb Z$ and $D=1$. Their formal-degree restriction is part of the definition of the usual constant-free Valiant classes and was needed for NN-C015's characteristic-zero height proof, but it is not needed for this modular Boolean simulation of promised bounded integer outputs. The proof also applies to explicitly supplied rational constants of polynomial total bit height.

This does not replace Burgisser's general complex-constant theorem. An arbitrary polynomial-size circuit over $\mathbb C$ need not come with a polynomial-rank finite algebra, polynomial-bit constants, or an effectively small algebraic replacement. Burgisser's constant-replacement system constructs bounded algebraic witness fields of possibly much larger degree, and its small residue-degree-one prime argument addresses that different presentation problem. The present proof uses the entire quotient algebra in polynomial rank and thus has no corresponding prime-splitting requirement.
\end{proof}
\paragraph{Status and complexity impact.}
Complete model refinement proved; complete self proof-check passed. Independent complete deduction-proof review by root: passed; imported input theorems are not claimed independently proved. The analytic transfer gate is now more precisely located: polynomial rank and an explicit polynomial-bit coefficient presentation suffice, with composite-modulus arithmetic and no RH input. The contrapositive transfers Boolean counting lower bounds to this restricted but broad arithmetic model unconditionally. It does not transfer ordinary $\mathrm P\ne\mathrm{NP}$ to $\mathrm{VP}_{\mathbb C}\ne\mathrm{VNP}_{\mathbb C}$, and does not give a converse to Burgisser. Arbitrary complex constants, large algebraic degree, absent presentations, and unrestricted division gates remain outside the theorem.

\paragraph{Geometric and representation dictionary.}
The coefficient object is a finite locally free affine scheme $\operatorname{Spec}B_n\to\operatorname{Spec}\mathbb Z$ with an explicit polynomial-rank trivialization of its pushforward algebra. The algorithm performs base change to $\operatorname{Spec}(\mathbb Z/M\mathbb Z)$ and evaluates in the full fiber algebra, including all of its components and nilpotents. It does not seek a rational point of that fiber. The scalar-output condition makes the resulting section descend to the base ring, where the magnitude bound identifies its unique integer lift. EGA IV4, printed p.111,18.2.1--18.2.2 gives the finite-algebra trace dictionary and base-change formula; the computational proof itself is the explicit multiplication table above. A weight or purity assertion without this small presentation does not supply the simulation, and no GCT orbit-closure obstruction is constructed.

\subsection{NN-C021 --- E computation transcripts have no polynomial-size translation to general resolution or polynomial calculus}
\begin{proposition}[NN-C021]
For the computation-transcript proof system of NN-C008, there is no fixed polynomial upper bound converting the size of every transcript refutation into the size of a resolution refutation of the same CNF. The analogous statement holds for ordinary polynomial calculus over $\mathbb F_2$, when proof size counts explicitly written nonzero monomials. More concretely, there are fixed-density 3-CNF contradictions on $N$ variables whose E transcripts have size
\[
 2^{O(N/\sqrt{\log N})},
\]
while every refutation in the relevant natural system has size $2^{\Omega(N)}$. The resolution and polynomial-calculus lower-bound families may be different. These natural-system lower bounds are classical inputs; the stated comparison with E's new transcript system is the deduction.
\end{proposition}
\paragraph{Exact inputs.}
E: The Exponential Time Hypothesis is False, supplied LaTeX. Theorem1.1, bounded-clause-density exponent.

NN-C008. complete computation-transcript proof-system construction.

Buss--Grigoriev--Impagliazzo--Pitassi, Linear Gaps between Degrees for the Polynomial Calculus Modulo Distinct Primes, JCSS62(2001),267--289. Author-hosted journal-final manuscript: Section2 pp.5--6; Lemma3 pp.7--8; Theorem15 p.17; Corollaries21--22 and proof p.19.

Ben-Sasson--Wigderson, Short Proofs Are Narrow: Resolution Made Simple, JACM 48 (2001), 149--169. Lemma4.2 p.155; Theorem4.4 and Corollary4.5 p.156; size-width Corollary3.6 p.155; proof of Theorem4.4 Section6.1 p.161.

Impagliazzo--Pudlak--Sgall, Lower Bounds for the Polynomial Calculus and the Groebner Basis Algorithm, Computational Complexity8(1999),127--144. Author-hosted accepted manuscript: Section6, Theorem6.2 and complete proof p.13; Corollary6.3 pp.13--14.

\begin{proof}
For resolution, take 3-regular connected expanding graphs with an odd total vertex charge and the associated Tseitin CNFs. BSW Lemma4.2 gives at most four clauses per graph vertex and one variable per edge, hence $O(N)$ clauses of width three. Their Corollary4.5 gives resolution size $2^{\Omega(N)}$; equivalently, combine their linear width lower bound Theorem4.4 with Corollary3.6. E Theorem1.1 nevertheless counts each such fixed-density 3-CNF in $2^{O(N/\sqrt{\log N})}$ time and polynomial space. Its count is zero. Recording all configurations as in NN-C008 multiplies the time by only a polynomial, absorbed by the same exponent. This gives the asserted transcript upper bound on this family.

For polynomial calculus over $\mathbb F_2$, use BGIP's Boolean mod-$p$ Tseitin family with $p=3$, over a fixed-degree expanding graph. Their Corollary21 states a linear degree lower bound for this Boolean encoding over characteristic two, and Corollary22's proof explicitly states that its 3-CNF translation has $O(N)$ clauses and variables and exponential monomial-size lower bounds. There is, however, a directional omission in the printed proof of Corollary21: its appeal to Lemma2 supplies a reduction from the Boolean system to the root-of-unity system, whereas the definition and Lemma1 require the reverse reduction to transfer the cited hardness. We supply that reverse reduction here instead of using that inference as printed.

Work in the fixed field $F=\mathbb F_4$, choose a primitive cube root $\omega$, and use BGIP's root-of-unity system $\operatorname{TS}_n(3)$ in edge variables $y_e$. Define the Boolean indicator substitutions
\[
 x_{e,i}=\delta_i(y_e),\qquad
 \delta_i(T)=\frac13\sum_{j=0}^{2}\omega^{-ij}T^j
 \quad(0\leq i<3).
\]
The scalar $3$ is invertible in characteristic two. On $T=\omega^a$ this polynomial is one for $a=i$ and zero otherwise. Every substituted Boolean axiom and clause of $\operatorname{BTS}_{n,3}$ involves only a fixed number $s$ of edge variables and has fixed degree. Include the domain equations $y_e^3-1$ for these variables and the corresponding local root-of-unity vertex or reverse-edge constraints $g_1,\ldots,g_t$. These local equations force exactly the local Boolean constraint in question.

Here is a complete constant-degree derivation argument. The algebra
\[
 F[y_1,\ldots,y_s]/(y_1^3-1,\ldots,y_s^3-1)
 \cong\prod_{a\in\{1,\omega,\omega^2\}^{s}}F
\]
has distinct roots in every coordinate. Let $q$ be the substituted Boolean constraint polynomial. At a tuple where all $g_j$ vanish, $q$ vanishes. At every other tuple, choose the first nonzero $g_j$ and set the value of $h_j$ to $q/g_j$, setting the other $h$ values to zero. Lagrange interpolation gives representatives $h_j$ of individual degree at most two, hence total degree at most $2s$, and
\[
 q=\sum_j h_jg_j+\sum_{e=1}^{s}u_e(y)(y_e^3-1).
\]
The last equality is an equality of ordinary polynomials: reduce the difference successively by the monic equations $y_e^3-1$. The remainder has individual degrees below three and vanishes at every root tuple, so it is zero. Since $s$, all input degrees and all interpolation degrees are fixed, the degrees of every displayed product are bounded by a constant $C$. Expand each multiplier into its finitely many monomials; each monomial multiple of an axiom is derivable by successive single-variable multiplication, and their linear combination gives $q$ in PC degree at most $C$. This proves a $(2,C)$ reduction from $\operatorname{TS}_n(3)$ to the Boolean system in the exact sense of BGIP's definition and Lemma1. Substitution of a degree-$d$ Boolean PC refutation therefore gives a root-of-unity refutation of degree at most $\max(C,2d)$.

BGIP Theorem15 bounds that latter degree below by $\Omega(n)$. A refutation over $\mathbb F_2$ is also one over $F$, so the Boolean system requires degree $\Omega(n)$ over $\mathbb F_2$. Its numbers of variables and clauses are $O(n)$ and its clause width is fixed. BGIP Lemma3 and its proof turn this into a 3-CNF with $O(n)$ variables and clauses while preserving the degree lower bound up to a constant factor. Impagliazzo--Pudlak--Sgall Theorem6.2 says that a proof with $M$ nonzero monomials gives a proof of degree at most $\max(d_0,2\lceil\sqrt{N\log M}\rceil+1)$, where $d_0$ is the initial degree. For constant $d_0$ and minimum refutation degree $d=\Omega(N)$, this implies $M=2^{\Omega(d^2/N)}=2^{\Omega(N)}$. This is also the exact size-degree chain invoked in BGIP Corollary22's proof. Apply E Theorem1.1 and NN-C008 to this 3-CNF family exactly as before.

Now suppose a fixed polynomial $P$ bounded the length of a natural-system refutation in terms of the length $L$ of every E transcript for the same formula. On one of the corresponding hard families,
\[
 P(L)\leq2^{O(N/\sqrt{\log N})}=2^{o(N)},
\]
contradicting the classical $2^{\Omega(N)}$ lower bound. This excludes even a universal polynomial size bound, and therefore in particular excludes a polynomial-time proof translation of that kind. The argument compares proofs of the same formula under its stated encoding; it does not forbid changing the target proof system by extension axioms or new inference rules.

For the familiar parity Tseitin equations over $\mathbb F_2$ themselves, summing all vertex equations already yields the contradiction in low degree. Thus the polynomial-calculus example deliberately uses mod three constraints over characteristic two, as in BGIP; the resolution example alone cannot justify a characteristic-two polynomial-calculus lower bound. This field distinction is essential.

No lower bound on unrestricted satisfiability algorithms follows from these proof lower bounds. BGIP Corollary22's algorithmic conclusion concerns implementations that explicitly represent polynomials as sums of monomials and generate the corresponding polynomial-calculus derivation. E's interpolation and global count computation need not produce such a derivation. NN-C009's compact semantic identity likewise does not constitute a short polynomial-calculus refutation.
\end{proof}
\paragraph{Status and complexity impact.}
Complete deduction and reverse low-degree reduction proved; primary source clauses and printed directional omission checked. Independent complete deduction-proof review by root: passed; imported input theorems are not claimed independently proved. This gives a concrete obstruction at the proof-system interface left open by NN-C009. E removes the exponential-time hypothesis but does not remove exponential resolution or characteristic-two polynomial-calculus proof-size lower bounds. The transcript system is still not polynomially bounded on all tautologies, so this comparison does not imply $\mathrm{NP}=\mathrm{coNP}$ or separate those classes.

\paragraph{Geometric and representation dictionary.}
For every hard contradiction, the Boolean-coordinate equation scheme is empty; its finite algebra has rank and Frobenius trace zero. E computes this vanishing efficiently relative to the former exhaustive exponent. Resolution width and polynomial-calculus degree/monomial size concern the cost of deriving the unit ideal from the chosen generators in specified syntactic systems. Thus the same zero scheme can have a short computation transcript and a long ideal-derivation certificate. The finite trace dictionary is SGA7 II XXII Corollary3.2, but that identity places no bound on the chosen presentation's derivation complexity. This is an effective distinction between a geometric zero object and an explicit proof of emptiness in a fixed calculus, not a new failure of purity or a representation-theoretic multiplicity statement.

\subsection{NN-C022 --- Local-map inversion and image distinguishing retain the original seed parameter}
\begin{proposition}[NN-C022]
Fix $k$. Let $F:\{0,1\}^n\to\{0,1\}^m$ be specified by a list of output scopes of size at most $k$ and their Boolean truth tables. For every target $y$, the number $|F^{-1}(y)|$, image membership, an inverse when one exists, and exact uniform sampling from a nonempty fiber are computable in
\[
 \operatorname{poly}_k(n+m)\,
 2^{O_k(n\log\log n/\sqrt{\log n})}
\]
time and polynomial space; sampling has this expected-time bound. If $m\leq Cn$, replace the exponent by $O_{k,C}(n/\sqrt{\log n})$. The number of output coordinates may be polynomially larger than $n$ in the first bound.

When $m>n$, these algorithms give a deterministic distinguisher for $F(U_n)$ against $U_m$, with advantage at least $1-2^{n-m}$ and the same time bound. Thus fixed-locality maps with polynomial output length cannot have $2^{\Omega(n)}$ inversion security or pseudorandom-generator security. This conclusion concerns the original seed length and does not assume expansion, random scopes, injectivity, or regular fibers.
\end{proposition}
\paragraph{Exact inputs.}
E: The Exponential Time Hypothesis is False, supplied LaTeX. Theorem1.1, arbitrary-density and bounded-density counting bounds.

NN-C002. prefix-count self-reduction and exact conditional sampling proof.

NN-C001. quasilinear fixed-degree polylog-sparse systems, for the logarithmic-locality extension.

Goldreich, Candidate One-Way Functions Based on Expander Graphs, 2011 revision of ECCC TR00-090. Section2, equation(1), printedp.78; Section2.2, equation(2), printedp.79; fixed versus logarithmic locality explicitly distinguished in Section2.

\begin{proof}
For each output coordinate $i$, enumerate the at most $2^k$ assignments $a$ to its scope $I_i$ for which $F_i(a)\ne y_i$. Exclude each one by the single clause
\[
 C_{i,a}=\bigvee_{j\in I_i}
 \begin{cases}x_j,&a_j=0,\\ \neg x_j,&a_j=1.\end{cases}
\]
This clause is false precisely on the forbidden local assignment. The conjunction $\Phi_y$ has exactly $n$ original variables, at most $2^k m$ clauses, and width at most $k$. Constant false coordinates give an immediate empty fiber; constant true constraints may be omitted. E applies with fixed width $\max(2,k)$. The satisfying assignments of $\Phi_y$ are exactly $F^{-1}(y)$, with no auxiliary gate variables or multiplicities. This proves the displayed counting bounds.

To find an inverse, keep the same $n$ declared variables and repeatedly count after fixing the next bit to zero, choosing that branch when its count is positive. There are at most $n$ extra unit clauses and $n+1$ counting calls. The bounded-density hypothesis remains bounded density when $m=O(n)$. Exact sampling uses the two child counts as probabilities; binary rejection sampling of a uniform integer below their sum has expected fewer than two trials. The proof in NN-C002 applies verbatim, including polynomial space. A nonempty fiber is inverted with certainty, for every target in the image; the argument is not merely an average-case attack.

Set $D_F(y)=1$ exactly when the fiber count is positive. Then
\[
 \Pr[D_F(F(U_n))=1]=1,\qquad
 \Pr[D_F(U_m)=1]=\frac{|\operatorname{im}F|}{2^m}
 \leq2^{n-m}.
\]
Subtracting gives the claimed advantage. In particular one bit of stretch already gives advantage at least one half, and larger stretch improves it. For polynomially bounded $m$, the polynomial description/read-time factor is $2^{o(n)}$, so both algorithms run in $2^{o(n)}$. If the map family is uniformly constructible, the attack is uniform; otherwise its given description is part of the input, or polynomial advice for that family. This refutes an exponential-in-$n$ security claim for this model, while allowing the possibility of hardness against all polynomial-time attacks.

Goldreich's equation(1) has exactly this form when the local predicate has a fixed number of arguments, regardless of the chosen expanding graph. His alternative equation(2) publishes the predicate table as part of the output; conditioning on that published table likewise leaves only the $n$ seed bits unknown. The paper also permits locality growing like $\log n$. An arbitrary truth table of that arity produces growing-width clauses, for which E's constants are not uniform in the width; the fixed-locality result alone does not cover that regime.

There is a useful algebraic extension of that boundary. Suppose the output predicates, although they may depend on $O(\log n)$ coordinates, have multilinear algebraic normal form over $\mathbb F_2$ of a fixed degree $d$, and suppose $m\leq n(\log n)^a$. Each such form has at most $\sum_{j=0}^{d}\binom{O(\log n)}j=(\log n)^{O_d(1)}$ monomials. Subtracting $y_i$ preserves this bound, so NN-C001 applies on the same $n$ original bits. The same inverse and image distinguisher therefore have $2^{O(n\log\log n/\sqrt{\log n})}$ time. Computing the algebraic normal forms from the explicit $2^{O(\log n)}$ truth tables takes polynomial time by Boolean finite differences. No assertion is made for a logarithmic-locality predicate of growing algebraic degree.

More generally, suppose some output coordinates are fixed-locality tables while all remaining coordinates are supplied by a Boolean circuit containing $s$ gates. Give a variable only to each of these $s$ gates and impose its constant-width Tseitin constraints, together with the requested output bits. The extension is unique for every seed. All local coordinates are still encoded without auxiliary variables. The resulting width-$\max(k,3)$ formula has $N=n+s$ variables and $O_k(m+s)$ clauses, hence costs
\[
 \operatorname{poly}_k(n+m+s)\,
 2^{O_k(N\log\log N/\sqrt{\log N})},
\]
with the sharper exponent $O(N/\sqrt{\log N})$ when $m+s=O(N)$. Thus it is the size of the genuinely nonlocal circuit part that enlarges the witness parameter, not the total number of local output constraints. All search and distinguishing conclusions persist with this explicit parameter substitution.
\end{proof}
\paragraph{Status and complexity impact.}
Complete deduction proved from E and NN-C001/002; self-check recorded. Independent complete deduction-proof review by root: passed; imported input theorems are not claimed independently proved. This narrows the cryptographic gate left by a generic Cook--Tseitin reduction: arbitrarily many polynomially listed local observations do not enlarge the unknown seed space. In particular, expansion can obstruct specific local or sequential inversion heuristics while failing to establish exponential security against E's global counting method. The conclusion does not refute the existence of ordinary one-way functions or pseudorandom generators, does not estimate E's constants at a practical key size, and does not make a claim for arbitrary dense multivariate-quadratic maps or logarithmic-locality random predicates of growing degree. The image test uses the original nonnegative integer fiber count; parity alone cannot replace it when a fiber has positive even cardinality.

\paragraph{Geometric and representation dictionary.}
Let $X_n=\operatorname{Spec}\mathbb Z[x_1,\ldots,x_n]/(x_j^2-x_j)$, the disjoint union of $2^n$ copies of $\operatorname{Spec}\mathbb Z$. Each Boolean output table has a unique integer multilinear interpolation which is an idempotent in this algebra, so $F$ defines a morphism $X_n\to X_m$. Over a Boolean target section $y$, its fiber is the finite split scheme with coordinate algebra $\mathbb Z^{|F^{-1}(y)|}$. The local clauses give a concise presentation of this same fiber. Its finite-locally-free rank, or trace of the identity, is the inverse count; reduction to $\mathbb F_2$ gives only the parity of that rank. E computes the rank without expanding the exponential product algebra. This is the explicit multiplication-matrix trace of EGA IV4,18.2.1--18.2.2, and the finite trace formalism of SGA4 XVII6.2.3. The image distinguisher compares support in the finite target with its uniform counting measure; the image has at most the source's number of components. No zeta-function estimate, cohomological purity theorem, or orbit-closure representation is needed, and A or M does not alter this parameter calculation.

\subsection{NN-C023 --- Parity hardness persists when the coherent Frobenius matrix has only one entry}
\begin{proposition}[NN-C023]
There is a deterministic polynomial-time construction taking a 3-CNF $\Phi$ to a nonzero homogeneous polynomial $F$ of degree $r$ in exactly $r\geq3$ variables over $\mathbb F_2$, given by an arithmetic formula of size $O(n+m)$, such that for the proper hypersurface
\[
 X_\Phi=V(F)\subset\mathbb P^{r-1}_{\mathbb F_2}
\]
one has
\[
 H^0(X_\Phi,\mathcal O)=\mathbb F_2,\qquad
 H^{r-2}(X_\Phi,\mathcal O)=\mathbb F_2,
\]
with all other coherent cohomology zero, and absolute Frobenius acts on the latter one-dimensional space by the scalar
\[
 \#\operatorname{SAT}(\Phi)\pmod2.
\]
Here $r=O(n+m)$, and the basis of the one-dimensional space has an explicit universal description. Consequently computing this single Frobenius entry is $\oplus\mathrm P$-complete on the specified factored class. A polynomial-time algorithm for that entry would imply $\mathrm P=\oplus\mathrm P$ and $\mathrm{NP}=\mathrm{RP}$. E computes it in subexponential time on this class. The hypersurfaces are allowed to be singular and reducible; no smoothness or purity assertion is part of the theorem.
\end{proposition}
\paragraph{Exact inputs.}
NN-C005. parsimonious Boolean gate encoding and parity-counting completeness.

NN-C006. effective Frobenius-data gate being sharpened.

Katz, Algebraic Solutions of Differential Equations (p-Curvature and the Hodge Filtration), Inventiones Mathematicae18(1972),1--118. Equations2.3.7.12--2.3.7.13 and Algorithm2.3.7.14, p.35; rank-one Special Case2.3.7.17, p.36.

SGA 7 II, Expose XXII. Theorem 3.1 (3.1.1); Corollary 3.2 (3.2.1).

E: The Exponential Time Hypothesis is False, supplied LaTeX. Theorem1.1.

Valiant--Vazirani, NP is as easy as detecting unique solutions, TCS 47 (1986), 85--93. Theorem 1.1; Corollaries 3.1--3.4, especially Corollary 3.2 p.91.

\begin{proof}
First suppose $\Phi$ has no empty clause and at least one variable. Regard each clause as a Boolean relation on its distinct variable scope. Add a tautological unary relation to every unused variable; this does not change the satisfying assignments and ensures every variable occurs. Let $\mathcal I$ be the set of factor--variable incidences and put $M=|\mathcal I|$. Then $n\leq M\leq3m+n$ after these additions. For each incidence $e=(f,i)$ introduce two distinct variables $z_{e,0},z_{e,1}$. Define integer polynomials
\[
 P_f(z)=\sum_{a\in\{0,1\}^{I_f}:\,a\text{ satisfies }f}
              \prod_{i\in I_f}z_{(f,i),a_i},
 \qquad
 Q_i(z)=\prod_{f:\,i\in I_f}z_{(f,i),0}
             +\prod_{f:\,i\in I_f}z_{(f,i),1}.
\]
Every $P_f$ is homogeneous of degree $|I_f|$ with at most eight distinct monomials. Every $Q_i$ is homogeneous of its positive incidence degree and has two distinct monomials. Put
\[
 P=\prod_f P_f\prod_i Q_i,\qquad
 Z=\prod_{e\in\mathcal I}z_{e,0}z_{e,1}.
\]
Both $P$ and $Z$ have degree $2M$. In a choice of one term from every factor contributing to the coefficient $[Z]P$, the choice of the second or first term of $Q_i$ forces every adjacent $P_f$ to choose the complementary rail. Thus all occurrences encode a single Boolean value of variable $i$. The chosen local assignment must satisfy each clause. Conversely every satisfying assignment makes exactly one consistent choice of all $P_f$ and $Q_i$ terms. All coefficients are positive integers equal to one before expansion, so this proves the exact identity
\[
 [Z]P=\#\operatorname{SAT}(\Phi)
\]
over $\mathbb Z$, including its reduction modulo two.

Adjoin $z_0$, set $r=2M+1$ and $F=z_0P$, and reduce coefficients modulo two. No $P_f$ or $Q_i$ is zero: their monomials are distinct and every clause has a satisfying local assignment. A polynomial ring over $\mathbb F_2$ is an integral domain, so $F\ne0$. It is homogeneous of degree $r$ in $r$ variables. Each rail occurs in exactly two factors, and the total number of factor incidences and multiplication gates needed to describe $F$ is $O(M)$; the eight-term local sums cost only a constant per clause. Moreover
\[
 [z_0Z]F=\#\operatorname{SAT}(\Phi)\pmod2.
\]
If an empty clause was present, use instead $F=z_1^3$ in three variables, whose relevant coefficient is zero. For the finitely many zero-variable cases use $z_1^3$ or $z_1z_2z_3$ according as the count is zero or one. These are again nonzero homogeneous degree-three equations.

We give the coherent calculation, which does not require smoothness. For every nonzero homogeneous degree-$r$ form $F$ in $r$ variables, its multiplication map on projective space is injective and gives
\[
 0\longrightarrow\mathcal O_{\mathbb P^{r-1}}(-r)
 \xrightarrow{\,F\,}\mathcal O_{\mathbb P^{r-1}}
 \longrightarrow\mathcal O_X\longrightarrow0.
\]
On the standard coordinate affine cover, the top Cech cohomology of $\mathcal O(-r)$ is represented by homogeneous Laurent monomials of total degree $-r$ having every exponent negative; monomials with a nonnegative exponent lie in a lower intersection and are zero in the top quotient. There is exactly one remaining monomial,
\[
 b=\frac1{z_1\cdots z_r}.
\]
The same Cech complex, split by Laurent monomial, is acyclic in intermediate degrees: for each exponent pattern other than the all-negative top pattern or the all-nonnegative global-section pattern, its nonzero components form the augmented simplex complex on the coordinates that may be inverted. Pairing subsets by adjoining a fixed admissible coordinate gives a contracting homotopy. Consequently $H^{r-1}(\mathbb P^{r-1},\mathcal O(-r))$ is the one-dimensional span of $b$, all its other cohomology groups vanish, and $\mathcal O$ has only $H^0=\mathbb F_2$. The long exact sequence gives the displayed cohomology of $X$ and the canonical boundary identification of its top group with the span of $b$.

The Frobenius square map on $\mathcal O_X$ lifts to the square map on $\mathcal O_{\mathbb P^{r-1}}$ and the map
\[
 a\longmapsto F a^2
\]
on $\mathcal O(-r)$, because $F(Fa^2)=(Fa)^2$. The latter map has the correct degree $-r$ when applied to a homogeneous section of degree $-r$. Naturality of the boundary map therefore identifies Frobenius on the top group with
\[
 b\longmapsto\frac{F}{(z_1\cdots z_r)^2}.
\]
In the top Cech quotient a numerator monomial survives only when every exponent is at most one. Since the numerator has total degree exactly $r$, that monomial must be $z_1\cdots z_r$. Thus the map is multiplication by $[z_1\cdots z_r]F$. This is precisely the rank-one version of Katz's coefficient formula, now verified for this possibly singular hypersurface. Relabeling $z_1,\ldots,z_r$ as $z_0$ and the rails proves the Frobenius assertion. SGA7 II XXII Corollary3.2 also yields
\[
 \#X_\Phi(\mathbb F_2)\equiv1+\#\operatorname{SAT}(\Phi)\pmod2.
\]

For hardness, the construction is a polynomial-time reduction from parity-3-SAT; that problem is $\oplus\mathrm P$-complete by the parsimonious circuit encoding in NN-C005. For membership on the specified factored class, nondeterministically choose one monomial from each displayed factor, and accept exactly when their exponent vectors sum to the all-one vector. There are constantly many choices at each factor, and exponent checking takes polynomial time. The number of accepting choices is the coefficient, so its parity is in $\oplus\mathrm P$. A polynomial-time Frobenius-entry algorithm therefore decides every $\oplus\mathrm P$ language, giving $\mathrm P=\oplus\mathrm P$; Valiant--Vazirani then gives $\mathrm{NP}=\mathrm{RP}$ as in NN-C006.

For E's upper bound one may simply count the original $\Phi$. There is also a direct bound from the factored input: encode the constant number of monomial choices of each factor by a constant number of bits, add validity constraints, and impose for every rail that exactly one of its two adjacent factors selects it. These are constant-arity Boolean constraints on $O(n+m)$ choice bits. There are $O(M)$ constraints, and $M=O(n+m)$, so their truth-table CNF has fixed width and bounded density. E Theorem1.1 gives time $2^{O((n+m)/\sqrt{\log(n+m)})}$ up to polynomial factors. Thus the single entry has E's subexponential algorithm, while its polynomial-time computation remains the precise parity-counting gate.
\end{proof}
\paragraph{Status and complexity impact.}
Complete construction and singular coherent calculation proved; full self-check and finite coefficient/point-count checks recorded. Independent complete deduction-proof review by root: passed; imported input theorems are not claimed independently proved. This strictly sharpens the effective-presentation issue in NN-C006. Exponential cohomology rank is not the only obstruction: here the relevant basis and rank are explicitly known and the entire missing matrix is one bit, yet computing that bit is parity-counting complete. A compact equation, a one-dimensional coherent cohomology group, and the existence of a trace formula do not constitute an efficient algorithm for Frobenius. The result concerns a succinct factored input; expansion into all monomials may already cost exponential time. It gives neither a lower bound for unrestricted arithmetic circuits nor a claim that the smooth hypersurface case is hard. A and M do not supply the missing entry.

\paragraph{Geometric and representation dictionary.}
The construction is already a scheme-theoretic theorem: parity-3-SAT is encoded in the absolute-Frobenius endomorphism of an explicitly trivialized rank-one coherent cohomology space of a projective hypersurface. Its sheaf-theoretic exact sequence and Cech basis prove the dictionary directly. Katz's original pp.35--36 locate the corresponding smooth Hasse--Witt formula, and SGA7 II XXII3.1--3.2 locate the zeta congruence and trace-point-count identity for arbitrary proper schemes. These reducible hypersurfaces are not asserted to furnish pure l-adic realizations of canonical bases, and their one-dimensional coherent Frobenius space is not a GCT representation multiplicity. The theorem identifies an effective arithmetic operation inside the geometric formalism, rather than importing Boolean hardness merely by analogy.

\subsection{NN-C024 --- The parity gate survives smooth hypersurfaces and genuine finite Weil objects}
\begin{proposition}[NN-C024]
Let $F$ be the degree-$r$ form constructed in NN-C023, so $r$ is odd and $c=[z_1\cdots z_r]F=\#\operatorname{SAT}(\Phi)\bmod2$. Put
\[
 H=\sum_{i=1}^r z_i^r,\qquad F_u=H+uF.
\]
The hypersurface $V(F_u)$ over $\mathbb F_2(u)$ is smooth and projective, its top coherent cohomology is a known one-dimensional space, and its absolute Frobenius is semilinear with matrix coefficient $uc$. Thus deciding whether the linearized Frobenius on this line is zero is parity-counting complete, even for this guaranteed generically smooth one-parameter family.

There is also a randomized polynomial-time reduction to smooth projective hypersurfaces over explicitly constructed finite fields $\mathbb F_q$, with $\log_2q=O(r\log r)$. Its output is smooth with probability greater than three quarters; on every smooth output the $\mathbb F_q$-linear $q$-power Frobenius scalar on the top coherent line is exactly $c$. Therefore a polynomial-time procedure correct on smooth inputs for this scalar, or for the point-count parity, would imply $\oplus\mathrm P\subseteq\mathrm{BPP}$ and then $\mathrm{NP}=\mathrm{RP}$. This randomized promise reduction does not assert $\mathrm P=\oplus\mathrm P$ for the finite-field problem.
\end{proposition}
\paragraph{Exact inputs.}
NN-C023. homogeneous coefficient encoding and complete coherent Frobenius calculation.

SGA1, ExposeII, Section4. Theorem4.10(iii) and Corollary4.12, Jacobian criterion for a hypersurface; text edition labels Theorem10 and Corollary12 within Section4.

Katz, Algebraic Solutions of Differential Equations, Inventiones Mathematicae18(1972),1--118. Algorithm2.3.7.14 pp.35--36; rank-one Special Case2.3.7.17 p.36.

SGA 7 II, Expose XXII. Theorem 3.1 (3.1.1); Corollary 3.2 (3.2.1).

Deligne, La conjecture de Weil I, Publications Mathematiques IHES43(1974),273--307. Theorem1.6, printedp.276.

Valiant--Vazirani, NP is as easy as detecting unique solutions, TCS 47 (1986), 85--93. Theorem 1.1; Corollaries 3.1--3.4, especially Corollary 3.2 p.91.

E: The Exponential Time Hypothesis is False, supplied LaTeX. Theorem1.1, for the inherited evaluation bound.

\begin{proof}
Because $r$ is odd, the partial derivatives of $H$ in characteristic two are $z_i^{r-1}$. We give an explicit degree bound for bad parameters without computing a discriminant. Set
\[
 L=r(r-2)+1,\qquad
 T=\binom{L+r-1}{r-1}=\binom{r(r-1)}{r-1}.
\]
Every degree-$L$ monomial $m$ is divisible by some $z_i^{r-1}$: otherwise every exponent would be at most $r-2$ and their sum at most $L-1$. Choose the least such $i=i(m)$. For each of the $T$ degree-$L$ monomials form
\[
 \frac{m}{z_{i(m)}^{r-1}}\frac{\partial F_u}{\partial z_{i(m)}}
 =m+u\frac{m}{z_{i(m)}^{r-1}}\frac{\partial F}{\partial z_{i(m)}}.
\]
In the degree-$L$ monomial basis their coefficient matrix is $B(u)=I+uA$ for a fixed matrix $A$ over $\mathbb F_2$. Consequently
\[
 \Delta(u)=\det B(u),\qquad \Delta(0)=1,\qquad \deg\Delta\leq T.
\]
Whenever $\Delta(a)\ne0$, every degree-$L$ monomial is in the ideal generated by the partial derivatives of $F_a$. They cannot have a common nonzero zero over an algebraic closure, since a nonzero coordinate would give a nonvanishing monomial $z_i^L$ in that ideal. The hypersurface is therefore smooth by SGA1 II4.10(iii), or its hypersurface Corollary4.12. In a projective coordinate chart this is the usual nonvanishing-differential criterion; Euler's identity $\sum_i z_i\partial_iF_a=rF_a$ ensures that at a point of $V(F_a)$ some derivative in a chart coordinate is nonzero whenever the homogeneous gradient is nonzero.

This proves smoothness over $\mathbb F_2(u)$, since $\Delta$ is a nonzero polynomial. It also proves that at most $T$ field elements can be excluded by this sufficient smoothness test. The exponentially sized matrix $B$ is used only in the proof of this bound. The reduction never forms $B$, computes $\Delta$, expands $F$, or tests smoothness. The integer $T$ has $O(r\log r)$ bits and can be computed by the product formula for the binomial coefficient in polynomial time. The form $H+uF$ has a polynomial-size circuit. It is nonzero for every $u$: the constructed $F$ is divisible by its extra coordinate, whereas $H$ is not; the explicit cubic edge-case replacements in NN-C023 are also not scalar multiples of $H$.

For completeness the required finite field can be built deterministically with polynomial bit cost. Choose $b$ to be a power of two, large enough that $2^b\geq8(T+1)$ and minimal with that property. Then $b=O(r\log r)$. Build fields of degrees $1,2,4,\ldots,b$ by successive quadratic extensions. In a constructed field $K$ of size $2^d$, the absolute trace
\[
 \operatorname{Tr}(a)=\sum_{j=0}^{d-1}a^{2^j}
\]
is a nonzero $\mathbb F_2$-linear functional: it takes values in $\mathbb F_2$, and its defining nonzero polynomial has degree $2^{d-1}<2^d$, so cannot vanish on every element. Evaluate it on a supplied vector-space basis to find an $a$ of trace one. Then $Y^2+Y+a$ has no root in $K$, since $\operatorname{Tr}(y^2+y)=0$, and hence is irreducible. Adjoin a root and use the evident doubled basis. Repeated squaring, testing the basis vectors, and the tower's arithmetic all take polynomial time in $b$.

If a single defining irreducible polynomial is desired, the last adjoined element $\theta$ has degree $b$ over $\mathbb F_2$: its degree divides the power of two $b$, and every proper subfield of $\mathbb F_{2^b}$ lies in the preceding field $\mathbb F_{2^{b/2}}$, whereas $\theta$ does not. Equivalently, any smaller degree divides $b/2$ and makes $\theta$ a root of $X^{2^{b/2}}-X$, whose complete root set is that preceding field. Linear algebra on $1,\theta,\ldots,\theta^b$ gives its monic minimal polynomial and the change to the power basis in polynomial time. Thus the construction uses an explicit polynomial-length finite-field presentation, not an oracle for unspecified algebraic constants.

Retain a basis beginning with $1$. Choose the parameter $a$ uniformly from the affine half-space of field elements whose coefficient on that basis vector is one. Every such $a$ is nonzero, there are $q/2$ choices, and exactly $b-1$ fair random bits suffice. At most $T$ choices are roots of $\Delta$, so
\[
 \Pr[V(F_a)\text{ is not smooth}]\leq\frac{2T}{q}<\frac14.
\]
The complete reduction has worst-case polynomial running time; it does not use an unbounded rejection loop.

Now apply the coherent calculation of NN-C023, or Katz2.3.7.17 on a smooth fiber. The coefficient of $z_1\cdots z_r$ in $H$ is zero because $r\geq3$. Absolute $2$-power Frobenius on the top coherent line therefore sends its displayed basis vector $e$ to $ac\,e$ and is semilinear on coefficients. Over the rational function field the coefficient is $uc$, so its linearization is zero exactly when $c=0$ and is an isomorphism otherwise; one must not call this a linear endomorphism over $\mathbb F_2(u)$. Over $\mathbb F_q$, iterate this semilinear map $b$ times. Its $\mathbb F_q$-linear scalar is
\[
 (ac)^{1+2+\cdots+2^{b-1}}=(ac)^{q-1}=c,
\]
because $a\ne0$ and $c\in\{0,1\}$. The distinction between the semilinear matrix coefficient $ac$ and the linear $q$-Frobenius scalar $c$ is essential.

All other coherent groups are as in NN-C023: $H^0=\mathbb F_q$, top degree $r-2$ has rank one, and the remaining groups vanish. SGA7 II XXII3.1 and3.2 consequently give, on a smooth sampled fiber (indeed on every nonzero fiber),
\[
 \#V(F_a)(\mathbb F_q)\equiv1+c\pmod2,
 \qquad
 Z(V(F_a)/\mathbb F_q,t)\equiv\frac{1-ct}{1-t}\pmod2.
\]
Here $r-2$ is odd. In particular the full zeta series modulo two is $1$ when $c=1$ and $(1-t)^{-1}$ when $c=0$.

Suppose a polynomial-time procedure computes the top $q$-Frobenius scalar on every smooth input in this class. Clock it at its polynomial bound and treat malformed outputs or timeouts arbitrarily. Its behavior on nonsmooth inputs is unrestricted. On a sampled input, correctness nevertheless has probability greater than three quarters because the smooth event alone has that probability. Independent repetitions and majority vote therefore put parity-3-SAT, and hence all of $\oplus\mathrm P$, in BPP. The same argument applies to a smooth-input point-count parity procedure, after adding one. It is a randomized promise reduction; it gives no deterministic parity algorithm.

To justify the stated NP consequence, use Valiant--Vazirani isolation followed by the BPP parity procedure, with its error amplified below the inverse polynomial isolation success probability. This gives a BPP procedure for SAT. For a satisfiable formula, self-reduce one variable at a time with decision error at most $1/(10n)$ per query; with probability at least nine tenths all needed answers are correct and a satisfying assignment is found. Check the resulting assignment directly. On an unsatisfiable input no string passes that final check. Thus SAT is in RP, and NP=RP. This verified-witness step is required to obtain one-sided error.

The generic-family problem has a stronger deterministic implication: a deterministic polynomial-time test of whether its linearized top Frobenius is zero directly decides parity-3-SAT, so gives P=$\oplus$P. E instead computes $c$ through the original formula or the monomial-choice encoding of NN-C023. On the bounded-density family this takes $2^{O(r/\sqrt{\log r})}$ time up to polynomial factors, including the polynomial field-construction overhead.
\end{proof}
\paragraph{Status and complexity impact.}
Complete smooth-pencil, field-construction and randomized-promise proofs; complete deduction-level self-reading passed. Independent complete deduction-proof review by root: passed; imported input theorems are not claimed independently proved. This removes singularity from the geometric obstruction in NN-C023, with a precisely stated randomized finite-field gate. The good fibers are genuine smooth projective varieties with the pure l-adic Frobenius representations guaranteed by Deligne's theorem. Producing such objects and knowing purity still leaves parity information computationally encoded in a rank-one coherent Frobenius map. This proof itself gives a randomized specialization; NN-C030 subsequently closes its deterministic finite-field specialization gate by using the actual encoding's pairwise homogeneity. A polynomial bound on all l-adic Betti numbers or an efficient full expansion of the succinct equation is still not asserted. Neither result separates VP from VNP. A and M do not calculate this coefficient or supply the missing polynomial-time procedure.

\paragraph{Geometric and representation dictionary.}
Let $U=\operatorname{Spec}\mathbb F_2[u,\Delta(u)^{-1}]$. The pencil defines a smooth projective hypersurface family over $U$, with a canonically based coherent line $R^{r-2}\pi_*\mathcal O$; its linearized relative Frobenius is multiplication by $uc$. The open smooth family is established by an explicit homogeneous Jacobian minor, which is the local differential criterion of SGA1 II4.10--4.12. The finite-field trace and full zeta congruence are exactly SGA7 II XXII3.1--3.2. Every good closed fiber has pure l-adic cohomology by Deligne, Weil I1.6; purity constrains the complex absolute values of Frobenius eigenvalues and does not compute this coherent characteristic-two invariant. The construction supplies genuine finite Weil objects for these counting gadgets, not canonical realizations of GCT's nonstandard basis or its structure constants. It sharpens an effectiveness gate while leaving that separate functorial/canonical realization problem intact.

\subsection{NN-C025 --- Geometrically integral smooth quadratic affine schemes retain exact counts above parity}
\begin{proposition}[NN-C025]
A $k$-CNF $\Phi$ with $n\geq1$ variables and $m$ clauses has a deterministic polynomial-time construction of a geometrically integral smooth affine scheme $Y_\Phi/\mathbb F_2$ of dimension exactly $n$, together with a finite etale degree-two map
\[
             Y_\Phi\longrightarrow\mathbb A^n_{\mathbb F_2},
\]
such that
\[
        \#Y_\Phi(\mathbb F_2)=2\#\operatorname{SAT}(\Phi).
\]
It has a presentation by $s+1$ equations of degree at most two in $n+s+1$ coordinates, with $s=O(n+km+\log(m+2))$, and an explicitly displayed Jacobian minor identically equal to one. Thus rational-point existence is NP-complete on this promised class, exact point counting is $\#\mathrm P$-hard by one division by two, and point counting modulo four encodes parity-SAT. Point counting modulo two is identically zero on the entire class. E evaluates these counts through the original $n$ witnesses, with its exact Theorem1.1 bound; geometric auxiliary coordinates do not become independent Boolean witnesses.
\end{proposition}
\paragraph{Exact inputs.}
NN-C006 and prior CX-E07. finite Boolean solution schemes and the correct affine trace boundary; used only for contrast, not repeated.

E: The Exponential Time Hypothesis is False, supplied LaTeX. Theorem1.1, original-variable exact k-CNF counter.

SGA1, ExposeII, Section4. Theorem4.10(iii) and Corollary4.12, Jacobian smoothness criterion.

Deligne, La conjecture de Weil I, PMIHES43(1974),273--307. Section1.3--1.5, especially compact-support trace equation1.5.1, printedpp274--275; contrast projective smooth Theorem1.6 p276.

Valiant--Vazirani, NP is as easy as detecting unique solutions, TCS 47 (1986), 85--93. Theorem 1.1; Corollaries 3.1--3.4, especially Corollary 3.2 p.91.

\begin{proof}
For a clause $C$, let $v_C(x)$ be the product of its false-literal indicators: a positive literal contributes $1+x_i$, and a negative literal contributes $x_i$. Put
\[
 S_\Phi(x)=\prod_C(1+v_C(x)),\qquad P_\Phi(x)=1+S_\Phi(x).
\]
On $\mathbb F_2^n$, $P_\Phi(x)=0$ exactly for satisfying assignments. The polynomial has degree at most $km$ and has a fan-in-two arithmetic circuit of $O(n+km)$ gates over $\mathbb F_2$. Empty products and empty clauses have their usual values, so tautological and contradictory inputs are included.

Choose the odd integer $d=2km+3$ and form
\[
        f_\Phi(x)=P_\Phi(x)+x_1^d+x_1.
\]
The added two terms vanish on every bit assignment. They ensure that the total degree is exactly the odd integer $d$, with leading monomial $x_1^d$. Binary powering takes $O(\log(d+1))$ additional gates. Define
\[
 Y_\Phi=\operatorname{Spec}
 \mathbb F_2[x_1,\ldots,x_n,t]/(t^2+t+f_\Phi(x)).
\]
The coordinate algebra is free of rank two over $\mathbb F_2[x]$, with basis $1,t$, because the equation is monic in $t$. Its relative derivative in $t$ is one. The hypersurface Jacobian criterion therefore makes its projection smooth of relative dimension zero, hence etale; alternatively it is finite flat and its relative differentials vanish. Composition with the smooth affine-space projection makes $Y_\Phi$ smooth of dimension $n$.

Geometric integrality is also explicit. Over any algebraic closure $\overline k$ of $\mathbb F_2$, a factorization of $T^2+T+f_\Phi$ over $\overline k(x)$ would give a rational function $h$ with $h^2+h=f_\Phi$. Write $h=A/B$ in coprime polynomials. The relation $A^2+AB=f_\Phi B^2$ forces every irreducible factor of $B$ to divide $A$, so $B$ is constant. Thus $h$ is a polynomial. If nonconstant, $h^2+h$ has even total degree $2\deg h$; if constant, it has degree zero or is zero. Neither possibility equals the odd-degree polynomial $f_\Phi$. The quadratic is therefore irreducible over the fraction field. Its monic quotient embeds in that field extension, so is a domain. This proves geometric integrality, not merely connectedness or reducedness.

For an $\mathbb F_2$ point, $t^2+t=0$ for both possible values of $t$. Each satisfying input $x$ therefore has exactly two points above it, and every nonsatisfying input has none. This proves the exact factor-two identity. No equations $x_i^2-x_i=0$ were imposed on the scheme: its geometric dimension is $n$, and its projection has two geometric points over every geometric base point. The vanishing or splitting of the rational fiber over a bit assignment is the arithmetic distinction being used.

To obtain the quadratic complete-intersection presentation, order the gates of the circuit computing $f_\Phi$ topologically. Introduce one coordinate $y_j$ for each gate, with equation
\[
 y_j-a-b=0,\quad y_j-ab=0,\quad\hbox{or}\quad y_j-a-1=0
\]
as appropriate, where $a,b$ are inputs or earlier gates. In characteristic two subtraction equals addition. Add the final equation $t^2+t+y_{\rm out}=0$. In the columns $(y_1,\ldots,y_s,t)$ and this order of equations, the Jacobian submatrix is lower triangular with diagonal entries all one: each gate equation has derivative one in its new gate, and the final equation has derivative one in $t$. Its determinant is identically one over the whole coordinate ring. Eliminating the gates, in order, recovers precisely the displayed Artin--Schreier coordinate algebra. The gate equations are a regular sequence by successive elimination, and the final monic polynomial is a nonzerodivisor. Thus this is a complete intersection, with constant-degree sparse equations and a proved global smoothness certificate. All coordinates and equations are written in polynomial time.

The exact-count complexity claim concerns this presentation. Guessing its polynomially many bit coordinates and checking its equations is a $\#\mathrm P$ computation. From its integer count one recovers $\#\operatorname{SAT}(\Phi)$ by division by two, so hardness is under a polynomial-time numerical reduction, not a falsely claimed parsimonious reduction for an always-even output. The divided count is itself a $\#\mathrm P$ function: guess just $x$ and check $P_\Phi(x)=0$. Its parity is exactly half of the point count modulo four. Consequently a deterministic polynomial-time mod-four algorithm gives P=$\oplus$P and, by Valiant--Vazirani, NP=RP. A polynomial-time existence algorithm would directly give P=NP. All these are conditional gates, not conclusions from smoothness.

For the upper bound, the constructor retains the source $k$-CNF presentation. Each assignment to its $n$ input bits determines every gate coordinate, and the extra coordinate $t$ contributes the known multiplicity two. E Theorem1.1 therefore gives
\[
 \operatorname{poly}_{k}(n+m)\,
 2^{O_k(n\log\log n/\sqrt{\log n})}
\]
bit operations and polynomial space, with exponent $O_{k,D}(n/\sqrt{\log n})$ when $m\leq Dn$. Applying a sparse-equation counter blindly to all $n+s+1$ coordinates would instead use that larger variable parameter. The smaller bound follows from the explicit functional graph and the original clause structure; it is not claimed for arbitrary succinct arithmetic circuits. Inputs with no variables have their constant count evaluated separately.
\end{proof}
\paragraph{Status and complexity impact.}
Complete affine construction, geometric-integrality and parameter proofs; complete deduction-level self-reading passed. Independent complete deduction-proof review by root: passed; imported input theorems are not claimed independently proved. This strengthens the affine boundary of NN-C006: even geometric integrality, smoothness, an everywhere invertible quadratic Jacobian minor, and a finite etale cover of affine space do not remove exact counting hardness. It is different from the earlier zero-dimensional Boolean solution scheme, whose number of components already equals the count. Here the geometric object is integral of the original witness dimension, and every rational count is even. The precise information lost by reduction to characteristic two is one binary digit: modulo four retains it. The construction does not supply a smooth projective compactification in polynomial-size ordinary coordinates, nor permit the pure smooth-projective conclusion of Deligne1.6. It does not resolve P versus NP or VP versus VNP; E changes the available running time while the reductions remain valid.

\paragraph{Geometric and representation dictionary.}
For every fixed-width CNF $\Phi$ on $n\geq1$ Boolean variables there is an explicitly constructed polynomial $f_\Phi\in\mathbb F_2[x_1,\ldots,x_n]$, represented by a polynomial-size arithmetic circuit and having odd degree, such that
\[
 U_\Phi=\operatorname{Spec}\mathbb F_2[x,t]/(t^2+t-f_\Phi(x))
\]
is smooth and geometrically integral of dimension $n$, finite etale of degree two over $\mathbb A^n$, and satisfies $\#U_\Phi(\mathbb F_2)=2\#\operatorname{SAT}(\Phi)$. Its graph presentation is a polynomial-size complete intersection of quadratic equations with a displayed triangular Jacobian minor equal to one. In particular its count is always zero modulo two, whereas its count modulo four is twice the Boolean parity. The proof above constructs $f_\Phi=1+S_\Phi+x_1^d+x_1$ with $d$ odd and larger than the verifier degree; the odd pole at infinity proves that $f_\Phi$ is not $h^2+h$ over the algebraic closure. Thus this geometric integrality claim has an explicit certificate argument.

The map is an Artin--Schreier torsor under the constant group $\mathbb Z/2\mathbb Z$. Over a satisfying rational input its fiber is $\operatorname{Spec}(\mathbb F_2\times\mathbb F_2)$; over a nonsatisfying one it is $\operatorname{Spec}\mathbb F_4$. These fibers have the same geometric degree, but respectively two and zero rational points. The multiplication trace of the unit in either rank-two algebra is zero in characteristic two; its algebraic trace locator is EGA IV4,18.2.1--18.2.2, printedp111. SGA4 XVII6.2.3 is the corresponding sheaf-theoretic finite trace, a distinct formalism. The whole affine point count is the alternating compact-support etale trace, by Deligne Weil I(1.5.1); it is not the trace on infinite-dimensional ordinary affine coherent $H^0$. Its integer residue modulo four is a separate operation, not a map from an odd-adic coefficient field to $\mathbb Z/4\mathbb Z$. Exact point counting on these circuit-presented smooth objects is counting-complete by division by two, while E evaluates the constructed objects through the original $n$-variable formula. No canonical GCT multiplicity or positive polytope formula is identified.

\subsection{NN-C026 --- Deterministic pure projective curves encode counting in the next binary digit}
\begin{proposition}[NN-C026]
From a $k$-CNF $\Phi$ with $n\geq1$ variables and $m$ clauses one can deterministically construct, in polynomial time, an explicitly presented field $\mathbb F_q$ with $\log_2q=O(n)$ and a smooth geometrically integral projective curve $C_\Phi/\mathbb F_q$ such that
\[
               \#C_\Phi(\mathbb F_q)=1+2\#\operatorname{SAT}(\Phi).
\]
The curve is given by a polynomial-size arithmetic circuit for a monic odd-degree Artin--Schreier equation, together with its explicit weighted-projective completion. Its degree and genus are exponentially large and are encoded in binary. There is no promise of a polynomial-degree dense equation or a polynomial-genus curve.

The integer count, and the integer trace of the pure weight-one l-adic Frobenius representation, are therefore $\#\mathrm P$-hard from these succinct descriptions. Reduction of the count modulo four is parity-counting complete after subtracting one and dividing by two. Nevertheless every curve in this class has
\[
 Z(C_\Phi/\mathbb F_q,t)\equiv(1-t)^{-1}\pmod2,
\]
and its entire coherent Frobenius matrix is determined by a simple entry formula from its genus and field, independently of the satisfying-assignment count. This entry formula does not require listing an exponentially large matrix. Thus a genuine pure finite Weil object can be constructed deterministically here, while the next integer counting digit remains a separate effective gate.
\end{proposition}
\paragraph{Exact inputs.}
NN-C024. deterministic quadratic-tower finite-field construction.

NN-C025. Boolean indicator and odd-degree Artin--Schreier irreducibility arguments.

E: The Exponential Time Hypothesis is False, supplied LaTeX. Theorem1.1, evaluation through the original formula.

SGA1, ExposeII, Section4. Theorem4.10(iii) and Corollary4.12, used on the two explicit affine curve charts.

The Stacks Project, Cohomology of Schemes. Lemma30.2.2, tag01XB (affine acyclicity); Lemma30.2.6, tag01XD (affine-cover Cech comparison).

SGA 7 II, Expose XXII. Theorem 3.1 (3.1.1); Corollary 3.2 (3.2.1).

Deligne, La conjecture de Weil I, PMIHES43(1974),273--307. Trace equation1.5.1 p275; Theorem1.6 p276; Poincare duality Theorem2.3 and discussion2.4--2.5 p281.

Valiant--Vazirani, NP is as easy as detecting unique solutions, TCS 47 (1986), 85--93. Theorem 1.1; Corollaries 3.1--3.4, especially Corollary 3.2 p.91.

\begin{proof}
Let $b$ be the least power of two with $b\geq n$, and put $q=2^b$. Then $n\leq b<2n$ unless $n=1$, when $b=1$. The field tower in NN-C024 constructs $K=\mathbb F_q$, its arithmetic, and a basis $e_1,\ldots,e_b$ deterministically in polynomial time. If needed it also supplies a monic irreducible defining polynomial of degree $b$. The absolute trace $\operatorname{Tr}:K\to\mathbb F_2$ is nonzero, by the polynomial-degree proof in NN-C024. Its pairing $(a,x)\mapsto\operatorname{Tr}(ax)$ is nondegenerate: multiplication by any nonzero $a$ is a bijection, so this pairing cannot vanish for every $x$. Linear algebra over $\mathbb F_2$ therefore constructs the dual basis $\beta_i$ with
\[
 \operatorname{Tr}(\beta_i e_j)=\delta_{ij}.
\]
Define the linearized coordinate polynomials
\[
 L_i(X)=\operatorname{Tr}(\beta_i X)
       =\sum_{j=0}^{b-1}(\beta_iX)^{2^j}.
\]
They have degree at most $q/2$, polynomial-size repeated-squaring circuits, and for $x=\sum_j a_j e_j\in K$ satisfy $L_i(x)=a_i\in\mathbb F_2$.

Let $S_\Phi$ be the Boolean satisfaction polynomial in NN-C025. Use all $b$ coordinates to define
\[
 P(X)=1+S_\Phi(L_1(X),\ldots,L_n(X))
                 \prod_{i=n+1}^{b}(1+L_i(X)).
\]
For every $x\in K$, $P(x)$ is zero or one; it is zero exactly when its first $n$ coordinates satisfy $\Phi$ and all remaining coordinates are zero. Hence exactly $\#\operatorname{SAT}(\Phi)$ field elements have $P(x)=0$, with no padding multiplicity. Its formal degree is at most $(km+b)q/2$ and its circuit size is polynomial in $n+m$ for fixed $k$.

Find $\lambda\in K$ with $\operatorname{Tr}(\lambda)=1$ by testing the basis vectors. Choose
\[
 d=(km+b+1)q+1,
 \qquad
 f(X)=\lambda P(X)+(X^q+X)X^{d-q}.
\]
The integer $d$ is odd, its binary length is $O(n+\log(m+2))$, and it exceeds the degree bound for $P$. The polynomial $f$ is monic of degree exactly $d$. All powers shown can be computed by binary powering, so its circuit and coefficient presentation have polynomial bit length. On $K$ the extra term vanishes, and
\[
           \operatorname{Tr}(f(x))=P(x).
\]
The linear map $y\mapsto y^2+y$ on $K$ has kernel $\{0,1\}$ and image equal to the trace-zero hyperplane. The image is contained in that hyperplane by telescoping the trace, and both have dimension $b-1$. Therefore the affine equation $y^2+y=f(x)$ has exactly two $K$-solutions in $y$ when $P(x)=0$, and none otherwise. Its number of affine points is exactly $2\#\operatorname{SAT}(\Phi)$.

Write $d=2g+1$ and set $w=g+1$. Let $f^h(X,Z)=Z^d f(X/Z)$ be its degree-$d$ homogenization. It has a polynomial-size homogeneous arithmetic circuit: at each addition gate multiply the lower formal-degree input by the appropriate power of $Z$, and homogenize multiplication gates multiplicatively. Every needed exponent has $O(\log d)$ bits, so binary powering keeps the overhead polynomial. Define
\[
 C_\Phi=V\bigl(Y^2+YZ^w-Z f^h(X,Z)\bigr)
              \subset\mathbb P_K(1,w,1).
\]
All terms have weighted degree $2w=d+1$. Weighted projectivity does not require constructing an exponentially large ordinary embedding. For existence of such an embedding, the $w$th Veronese ring of $K[X,Y,Z]$ with weights $(1,w,1)$ is generated in degree one by $Y$ and $X^aZ^{w-a}$ for $0\leq a\leq w$. Its Proj is the same weighted projective space and is a projective closed subscheme of ordinary projective space. The latter embedding has exponentially many coordinates and is not part of the polynomial-size output.

The curve is covered by the two weight-one charts $Z\ne0$ and $X\ne0$: a hypothetical point with $X=Z=0$ would have $Y^2=0$ and hence all homogeneous coordinates zero. On $Z=1$ its equation is $y^2+y=f(x)$, whose $y$ derivative is one. On $X=1$, put $z=Z/X$ and $v=Y/X^w$; the equation is
\[
       v^2+v z^w=z^{d+1}f(1/z).
\]
At $z=0$ this has the unique point $v=0$. Its derivative in $z$ is one at that point, because $f$ is monic of degree $d$ and the right side starts with $z$. Thus the curve is smooth everywhere, over an algebraic closure as well, by the local Jacobian criterion. There is exactly one $K$-rational point at infinity. The ambient possible singular point $[0:1:0]$ is off the curve, so it supplies no exception to this chart argument.

The affine polynomial $y^2+y+f(x)$ is geometrically irreducible by the odd-degree proof in NN-C025. The smooth point at infinity lies on the closure of that affine curve: on its smooth one-dimensional chart, $z$ is a nonzero regular function and its zero set is just that point. Thus the projective curve is geometrically integral. This proves
\[
              \#C_\Phi(K)=1+2\#\operatorname{SAT}(\Phi)
\]
with no smoothness promise and no randomized parameter selection.

We compute the genus and the coherent Frobenius directly, so no unverified genus formula is needed. On the intersection of the two affine charts, the coordinate ring has the free $K[x,x^{-1}]$-basis $1,y$. Sections from the $Z=1$ chart have coefficients in $K[x]$ for both basis elements. Sections from the $X=1$ chart have constant coefficients in $K[x^{-1}]$ and $y$-coefficients in $x^{-w}K[x^{-1}]$, since $v=x^{-w}y$. The Cech quotient has no constant part, and its $y$-part has the basis
\[
       e_j=[x^{-j}y],\qquad 1\leq j\leq g.
\]
The same intersection calculation gives $H^0(C_\Phi,\mathcal O)=K$. The two affine opens and their affine intersection compute coherent cohomology by affine acyclicity and the Cech comparison (Stacks, Lemmas30.2.2 and30.2.6); hence $H^1$ has dimension $g=(d-1)/2$. In particular the genus is exponential in the original witness parameter, despite the short equation.

Absolute two-power Frobenius sends
\[
 e_j\longmapsto[x^{-2j}y^2]
       =[x^{-2j}y]+[x^{-2j}f(x)].
\]
The second summand is a Laurent polynomial in $x$ with no $y$ term, and is a Cech coboundary. The first is $e_{2j}$ if $2j\leq g$ and is a coboundary otherwise. Thus its semilinear matrix is the coefficient-independent shift
\[
 e_j\longmapsto
 \begin{cases}e_{2j},&2j\leq g,\\0,&2j>g.\end{cases}
\]
Iterating $b$ times gives the $K$-linear $q$-Frobenius shift $e_j\mapsto e_{qj}$, with out-of-range terms zero. It is nilpotent, with known rank $\lfloor g/q\rfloor$, independently of $f$ and $\Phi$. Applying SGA7 XXII3.1 to this proper curve gives the asserted complete zeta congruence $(1-t)^{-1}$, not just equality of first point-count residues. At fixed $n,m,k$, both the field and this entire based coherent Frobenius description are fixed, although satisfying-assignment counts can differ.

Finally let $\ell$ be any odd prime. The geometric l-adic cohomology is a genuine finite Weil representation. Its $H^1$ is pure of weight one by Deligne1.6. The connected curve has trace one on $H^0$ and trace $q$ on $H^2$, by the duality and degree discussion2.3--2.5. Deligne's trace formula1.5.1 gives the integer identity
\[
 \operatorname{Tr}(\operatorname{Frob}_q\mid H^1_{\rm et})
           =q+1-\#C_\Phi(K)
           =q-2\#\operatorname{SAT}(\Phi).
\]
Purity controls the eigenvalue magnitudes, but computing this integer trace, or its residue modulo four, recovers the count or its parity respectively. This is reduction of an integer trace, not reduction of an odd-adic field modulo four.

The divided count belongs to $\#\mathrm P$ directly on the whole succinct Artin--Schreier class: guess the $b$ coordinate bits of $x\in K$, evaluate the circuit for $f(x)$, and test $\operatorname{Tr}(f(x))=0$. For every monic odd-degree member, it equals $(\#C(K)-1)/2$. Our reduction proves its $\#\mathrm P$-completeness, and its parity is $\oplus\mathrm P$-complete. A deterministic polynomial-time mod-four count procedure would therefore give P=$\oplus$P and NP=RP; exact polynomial-time counting gives the stronger counting and Toda consequences recorded in NN-C006. E instead evaluates the constructed instances from the source $k$-CNF in its original $n$-variable time bound, with polynomial construction overhead. No efficient method for arbitrary high-degree circuit-presented curves is inferred.
\end{proof}
\paragraph{Status and complexity impact.}
Complete deterministic field, curve, Cech and counting proofs; complete deduction-level self-reading passed. Independent complete deduction-proof review by root: passed; imported input theorems are not claimed independently proved. This supplies a deterministic smooth projective finite-field counting object, strengthening the existence side of the geometric gate. It does not derandomize the low-degree rank-one-coherent construction of NN-C024: the curve degree and coherent rank are exponential, and the parity information now occurs modulo four rather than modulo two. The representation is an actual pure l-adic representation attached to a smooth curve, but obtaining its integer trace from the succinct input remains counting-hard. All characteristic-two coherent Frobenius data above are already coefficient-independent, showing by explicit construction why that residue-level realization cannot determine the next integer digit. These are valid reduction barriers, not unconditional complexity separations or a GCT canonical-basis realization.

\paragraph{Geometric and representation dictionary.}
For each fixed-width CNF $\Phi$ on $n\geq1$ variables, the proof explicitly constructs a finite field $K=\mathbb F_q$, $q=2^b$ with $b$ polynomial in the input size, and a monic odd-degree circuit-presented polynomial $f_\Phi(X)\in K[X]$. Put $d=\deg f_\Phi=2g+1$ and $w=g+1$. Its weighted projective completion
\[
 C_\Phi:\quad Y^2+YZ^w=Z^{d+1}f_\Phi(X/Z)
       \subset\mathbb P(1,w,1)
\]
is a smooth geometrically integral projective curve of genus $g$, with one rational point at infinity and
\[
 \#C_\Phi(K)=1+2\#\operatorname{SAT}(\Phi).
\]
The polynomial and the integers $d,g$ have polynomial-size circuit/binary descriptions; $d$ and $g$ themselves may be exponential. The two affine charts in the proof establish smoothness, completion and the Cech basis $e_j=[x^{-j}y]$, $1\leq j\leq g$, for $H^1(C_\Phi,\mathcal O)$. Absolute Frobenius sends $e_j$ to $e_{2j}$, with out-of-range indices zero; $q$-Frobenius sends it to $e_{qj}$. These coherent operators are nilpotent, independently of the encoded satisfying count.

For every prime $\ell\ne2$, the actual $H^1_{\rm et}(C_{\overline K},\mathbb Q_\ell)$ is pure of weight one and has Frobenius trace $q-2\#\operatorname{SAT}(\Phi)$. SGA7 XXII3.1 explains the common full zeta reduction modulo two; Deligne Weil I1.5.1,1.6 and2.3--2.5 supply the integral trace and purity interfaces. This proves a precise contrast between the characteristic-two coherent realization and the next binary digit of the integer etale trace. The divided count $(\#C(K)-1)/2$ is a counting-complete function on the succinct class, and its parity is parity-counting-complete. E evaluates the constructed instances from their original Boolean witnesses; it does not give a polynomial algorithm on arbitrary high-degree circuit-presented curves. No plethysm/Kronecker multiplicity identification or canonical-basis construction is claimed.

\subsection{NN-C027 --- Smooth integral cubic arithmetic schemes carry counting-hard pure Tate multiplicities}
\begin{proposition}[NN-C027]
Let $N=\#\operatorname{SAT}(\Phi)$ for a fixed-width CNF with $n\geq1$ input variables and $m$ clauses. In polynomial time one can construct a smooth affine scheme $X_\Phi/\mathbb Z$, with every geometric fiber integral, given by one sparse cubic equation in $O(n+m)$ coordinates. There are explicit integers $v,h=O(n+m)$, with $v,h\geq1$, for which
\[
 [X_\Phi]=\mathbb L^{v+h}+N\mathbb L^h,
 \qquad
 \#X_\Phi(\mathbb F_q)=q^{v+h}+Nq^h
\]
for every prime power $q$, and
\[
         \zeta_{X_\Phi}(s)=\zeta(s-v-h)\zeta(s-h)^N.
\]
The Grothendieck-class identity is proved by an explicit stratification. Separately, over any field $K$ and any prime $\ell$ invertible in $K$, the actual geometric compactly supported l-adic groups are
\[
 H_c^{2h}=\mathbb Q_\ell(-h)^N,
 \qquad H_c^{2(v+h)}=\mathbb Q_\ell(-v-h),
\]
and all others vanish, as representations of the absolute Galois group of $K$. Over finite fields these particular groups are pure Tate; this is proved, not inferred from their point counts alone.

For every nonreal nontrivial zeta zero $\rho$ in $0<\Re\rho<1$, one has
\[
 \operatorname{ord}_{s=h+\rho}\zeta_{X_\Phi}(s)
              =N\operatorname{ord}_{s=\rho}\zeta(s).
\]
The shifted zero locations are therefore controlled by A, while their multiplicity contains the original exact Boolean count. No use of M is made. E computes the multiplicity in its original $n$-variable time bound. Polynomial-time extraction of that multiplicity from these cubic presentations would compute every $\#\mathrm P$ function in polynomial time.

There is also a literal quadratic-character variant over $B=\mathbb Z[1/2]$:
\[
 X^-_\Phi:\ t^2+1=\sum_i y_i f_i(x),\qquad
 \zeta_{X^-_\Phi}(s)=\zeta^{(2)}(s-v-h)L(s-h,\chi_{-4})^N,
\]
where $\zeta^{(2)}(s)=(1-2^{-s})\zeta(s)$. Its lower group is
$H_c^{2h}=(\chi_{-4}\otimes\mathbb Q_\ell(-h))^N$ and its top group is the same Tate line. At any nonreal nontrivial zero $\rho$ of this Dirichlet L-function, its shifted zeta has order $N\operatorname{ord}_\rho L$. Both variants are smooth with geometrically integral fibers; their fixed-$\mathbb F_3$ point-count parity is $1+N\bmod2$.
\end{proposition}
\paragraph{Exact inputs.}
Prior catalogue CX-E07. constant finite etale Boolean Z-scheme with degree equal to the satisfying-assignment count; background construction only.

E: The Exponential Time Hypothesis is False, supplied LaTeX. Theorem1.1, original-variable counting bound.

A: Alpoge, Remark; designated RH and Dirichlet-GRH input theorems. riemannHypothesis and riemannZeta\_eq\_zero\_re, Main.lean lines4,6; LFunction\_eq\_zero\_re\_of\_isPrimitive, Main.lean line12, as recorded in OURS\_ALPOGE\_LEAN\_MAP\_20261120.md Section3.

SGA1, ExposeII, Section4. Theorem4.10(iii) and Corollary4.12, relative Jacobian criterion over the integer base.

SGA4, ExposeXVII. Theorem5.2.6 and Proposition5.2.8 (!-base change and geometric stalks); Theorem5.4.3 (compact-support Kunneth); localization5.1.16.3.

SGA4, ExposeXVIII. Equation2.8.1 (affine-space top trace), Theorem2.9 (relative trace, base change and normalization); Proposition1.1.6 and equation1.1.3.2 for the projective-line calculation.

NIST Digital Library of Mathematical Functions. Equation25.4.2 (zeta functional equation); equation25.2.11 (Euler product on Re(s)>1).

NN-C006. exact-count and parity complexity implications, including Toda and Valiant--Vazirani locators.

\begin{proof}
First recall the parameter content of the earlier finite Boolean construction, without counting it as a new finding. Compile the CNF into an AND/NOT verifier circuit with $O(n+km)$ gates. Let $x_1,\ldots,x_v$ be all input and gate coordinates. Impose their Boolean equations $x_j^2-x_j=0$, the quadratic gate equations $z-ab=0$ and $z+a-1=0$, and the acceptance equation $z_{\rm out}-1=0$. These are $h=O(v)$ polynomials $f_1,\ldots,f_h\in\mathbb Z[x]$ of degree at most two and constant monomial sparsity. The Boolean algebra splits as a product of copies of $\mathbb Z$ indexed by bit assignments. Each gate or acceptance equation evaluates to zero or to a unit $\pm1$ in each factor. Consequently
\[
        Y=V(f_1,\ldots,f_h)\subset\mathbb A^v_\mathbb Z
           \simeq\coprod_{N}\operatorname{Spec}\mathbb Z.
\]
There is one extension of each accepting input to the gate coordinates. Retain a Boolean polynomial, say $f_1=x_1^2-x_1$, which is nonzero over every field.

Now introduce $y_1,\ldots,y_h,t$ and define
\[
 X=X_\Phi=V\left(t^2-t-\sum_{i=1}^{h}y_i f_i(x)\right)
                  \subset\mathbb A^{v+h+1}_\mathbb Z.
\]
Its single equation has degree at most three, $O(v+h)$ monomials and constant-size integer coefficients. Denote its left side by $G$. In its coordinate ring,
\[
 1=(2t-1)^2-4\sum_i y_i f_i
   =(2t-1)\frac{\partial G}{\partial t}
       +4\sum_i y_i\frac{\partial G}{\partial y_i}.
\]
Thus the relative Jacobian ideal contains one over $\mathbb Z$ itself. The relative hypersurface criterion proves that the projection $\pi:X\to\mathbb A^v_\mathbb Z$ is smooth of relative dimension $h$, and hence $X\to\operatorname{Spec}\mathbb Z$ is smooth of relative dimension $v+h$. This argument treats characteristic two and odd characteristics uniformly.

To prove geometric integrality, work over an algebraically closed field $K$ of any characteristic and view $G$ as a polynomial in $y_1$ over
$R=K[x,t,y_2,\ldots,y_h]$. Its coefficient is $-f_1\ne0$. The other coefficient is $t^2-t-\sum_{i>1}y_i f_i$, a monic polynomial in $t$. It has no common irreducible factor with $f_1$: every such factor of $f_1$ depends only on $x$, whereas the leading coefficient in $t$ is one. Hence $G$ is primitive over the unique-factorization domain $R$. It is linear and irreducible over $\operatorname{Frac}(R)$, so Gauss's lemma makes it irreducible, and therefore prime, in $R[y_1]$. Every geometric fiber is a domain. The monic equation in $t$ also makes the original coordinate algebra free of rank two over $\mathbb Z[x,y]$.

The count and class have an elementary proof valid over every base. Let $U_i$ be the locally closed stratum where $f_1=\cdots=f_{i-1}=0$ and $f_i$ is invertible. On $U_i$, solve the equation for $y_i$. The other $h-1$ coordinates $y_j$ and $t$ are free, giving
\[
             \pi^{-1}(U_i)\simeq U_i\times\mathbb A^h.
\]
On $Y$, the equation is $t(t-1)=0$; the two factors are comaximal over every ring. Thus
\[
             \pi^{-1}(Y)\simeq
                   (Y\times\mathbb A^h)\coprod(Y\times\mathbb A^h).
\]
Here $[\ ]$ denotes the Grothendieck ring with the usual closed-open scissor relation and $\mathbb L=[\mathbb A^1_\mathbb Z]$. Summing the displayed strata gives
\[
 [X]=\mathbb L^h([\mathbb A^v]-[Y])+2\mathbb L^h[Y]
     =\mathbb L^{v+h}+\mathbb L^h[Y]
     =\mathbb L^{v+h}+N\mathbb L^h.
\]
The same stratification proves the point count, with $N$ rational points of $Y$ over every finite field. Equivalently, for a base point with some $f_i(x)\ne0$, every choice of $t$ and the other $h-1$ coordinates gives one point, so its fiber has $q^h$ points. At a point of $Y$, the fiber has $2q^h$ points. No cancellation estimate or character-sum bound is needed.

We next justify the stronger actual-cohomology assertion. This step is separate from the Grothendieck-class calculation. Fix any base field $K$ and a prime $\ell$ invertible in $K$, with rational coefficients $\mathbb Q_\ell$. Work on the geometric base for cohomology and retain the descent data from $K$. The trace, stalk and projector argument does not require a finite field; Frobenius will be used only for the subsequent finite-field interpretation. The standard affine-space calculation is
\[
 R\Gamma_c(\mathbb A^h_{\overline K},\mathbb Q_\ell)
                   =\mathbb Q_\ell(-h)[-2h].
\]
For $h=1$ this follows from Kummer theory and $\operatorname{Pic}(\mathbb P^1)=\mathbb Z$, the projective-line top trace, and the localization sequence at infinity; powers follow from compact-support Kunneth. Exact source interfaces are SGA4 XVIII1.1.3.2 and Proposition1.1.6, XVII5.1.16.3, XVII5.4.3, and XVIII2.8.1. The !-base-change theorem XVII5.2.6 and its geometric-stalk Proposition5.2.8 now show that $R\pi_!\mathbb Q_\ell$ is concentrated in degree $2h$. Write its sole sheaf as $\mathcal T$.

The involution $\iota:t\mapsto1-t$ acts on $\mathcal T$. Since two is invertible in $\mathbb Q_\ell$, the projectors $(1\pm\iota)/2$ split it into $\mathcal T^+\oplus\mathcal T^-$. Off $Y$ the fiber is affine $h$-space, and $\iota$ is an affine automorphism in the displayed coordinates; it acts trivially on its top compact cohomology. On $Y$ it exchanges the two affine-space components. The relative trace of SGA4 XVIII2.9 gives
\[
                  \mathcal T(h)\longrightarrow\mathbb Q_\ell.
\]
On the plus part it is an isomorphism on every geometric stalk: off $Y$ it is the affine-space trace, and on $Y$ it is the sum of the two component traces, restricted to the diagonal line. The latter sends $(a,a)$ to $2a$; its normalized basis is half the sum of the two component classes. This remains invertible for $\ell=2$ because the coefficient field is $\mathbb Q_2$, not $\mathbb Z_2$. A morphism of etale sheaves that is an isomorphism on every geometric stalk is an isomorphism, so this proves a global identification of the plus sheaf, not just equality of ranks.

The minus part vanishes off $Y$. Base change to $Y$ identifies it globally with the anti-diagonal in the two labeled components $t=0$ and $t=1$, namely $\mathbb Q_\ell(-h)$. If $i:Y\hookrightarrow\mathbb A^v$ is the immersion, the canonical localization map $\mathcal T^-\to i_*i^*\mathcal T^-$ is an isomorphism on every geometric stalk, and hence an isomorphism of sheaves. It follows that $\mathcal T^-=i_*\mathbb Q_\ell(-h)$. All maps are defined over $K$ and are Galois-equivariant; when $K$ is finite they commute with Frobenius. Since the original derived image has a single cohomology sheaf, it is canonically that sheaf shifted. We have proved
\[
 R\pi_!\mathbb Q_\ell
  \simeq\bigl(\mathbb Q_\ell\oplus i_*\mathbb Q_\ell\bigr)(-h)[-2h].
\]
Taking compactly supported cohomology on $\mathbb A^v$ gives exactly the two claimed groups, because $Y$ is $N$ split points. Their degrees are distinct because $v\geq1$. Involution acts as minus one on $H_c^{2h}$ and plus one on the top group. These are actual Tate representations, pure of the displayed distinct weights when the base field is finite; no inference from an arbitrary positive polynomial point count is used.

The local zeta function at a prime $p$ follows from the all-extension-field count:
\[
 Z(X_{\mathbb F_p},T)
       =\frac{1}{(1-p^{v+h}T)(1-p^hT)^N}.
\]
Taking the Euler product, initially for $\Re s>v+h+1$, gives the stated arithmetic zeta identity, and the right side gives its meromorphic continuation. A's RH theorem places the nontrivial zeros of the two Riemann-zeta factors on the shifted lines $\Re s=v+h+1/2$ and $\Re s=h+1/2$, with the second factor present only when $N>0$. For any fixed nonreal nontrivial zero $\rho$ of $\zeta$, use only $0<\Re\rho<1$ and $\Im\rho\ne0$. At $s=h+\rho$ the other factor is $\zeta(\rho-v)\ne0$: its argument has negative real part and nonzero imaginary part, and the functional equation together with the zero-free Euler product on $\Re s>1$ excludes a zero there. Hence the displayed order formula holds without A; A is needed here only to specify the shifted critical lines. No simplicity assumption is required, and no algorithm for measuring this analytic order from a cubic equation is supplied by RH.

The quadratic-character variant uses the same $f_i,v,h,N$, over $B=\mathbb Z[1/2]$, and has equation $t^2+1=\sum_i y_i f_i(x)$. Let $G^-=t^2+1-\sum_i y_i f_i$. In its coordinate ring over $B$,
\[
 1=-\frac t2\frac{\partial G^-}{\partial t}
                     -\sum_i y_i\frac{\partial G^-}{\partial y_i}.
\]
Thus the projection $\pi^-:X^-\to\mathbb A^v_B$ is smooth of relative dimension $h$, by the relative hypersurface criterion already sourced above. For geometric integrality, view $G^-$ as a polynomial linear in a coordinate $y_i$ whose coefficient $f_i$ is nonzero. That coefficient has no common irreducible factor with the remaining coefficient, which is monic in $t$. The primitive-linear-polynomial argument above applies in every residue characteristic of $B$.

Off $Y=V(f_1,\ldots,f_h)$, solving for a coordinate $y_i$ again gives affine $h$-space fibers and the same locally closed affine-bundle stratification. Over $Y$ the scheme is
\[
       Y\times_B\operatorname{Spec}B[i]\times_B\mathbb A^h_B.
\]
This is the actual quadratic etale cover with its two conjugate geometric points. For $p$ odd, the number of roots of $t^2+1$ in $\mathbb F_{p^a}$ is $1+\chi_{-4}(p)^a$: a square root of minus one is present precisely when $p^a\equiv1\pmod4$. Since $Y$ has $N$ rational points over every finite field, adding the rational fibers proves
\[
 \#X^-(\mathbb F_{p^a})=p^{a(v+h)}+N\chi_{-4}(p)^a p^{ah}.
\] The corresponding scissor identity is
\[
 [X^-]=\mathbb L^{v+h}
       +N\bigl([\operatorname{Spec}B[i]]-[\operatorname{Spec}B]\bigr)\mathbb L^h
\]
over $B$. Its difference of classes is the quadratic sign class; the entire scheme is nevertheless an honest smooth integral geometric object.

The full compact-support realization is not inferred solely from this class. Use the cohomological projector proof above for the involution $t\mapsto-t$. Over a geometric base point off $Y$, its action on the top compact cohomology of affine $h$-space is the identity. Over $Y$, it interchanges the two components indexed by $i$ and $-i$. Consequently the plus part of $R^{2h}\pi^-_!\mathbb Q_\ell$ is the constant Tate sheaf, by the global trace normalization. The minus part is supported on $Y$ and is the anti-invariant summand of the rank-two permutation sheaf of the quadratic cover, tensored with $\mathbb Q_\ell(-h)$. This is the sheaf of the character $\chi_{-4}$. Geometric-stalk base change, the global relative trace, and the canonical localization map identify these summands as sheaves, exactly as above; there is no unresolved extension between them because the idempotent projectors split it. All other direct-image sheaves vanish because the geometric fibers are unions of affine $h$-spaces.

Taking compact supports on $\mathbb A^v$ therefore gives
\[
 H_c^{2h}(X^-_{\overline{\mathbb F}_p},\mathbb Q_\ell)
       =\bigl(\chi_{-4}\otimes\mathbb Q_\ell(-h)\bigr)^N,
 \qquad
 H_c^{2(v+h)}=\mathbb Q_\ell(-v-h).
\]
Here the character notation denotes the unramified local Frobenius representation at $p$; its eigenvalue is $\chi_{-4}(p)$. Thus the lower group has Frobenius eigenvalues $\chi_{-4}(p)p^h$, of complex absolute value $p^h$, and is pure of weight $2h$. The upper group is pure of weight $2(v+h)$. The degrees differ because $v\geq1$. This explicit calculation uses only classical etale trace and base-change theorems; A is used afterward for the global analytic zero locations.

The local zeta factor is
\[
 \frac{1}{(1-p^{v+h}T)(1-\chi_{-4}(p)p^hT)^N}.
\]
Multiplication over odd primes gives the stated arithmetic zeta identity in its half-plane of absolute convergence and then meromorphically. The omitted prime-two factor of the Tate term is exactly the factor in $\zeta^{(2)}$. Its extra zeros are explicit Euler-factor zeros, not nontrivial zeros of a Dirichlet L-function.

For any nonreal nontrivial zero $\rho$ of $L(s,\chi_{-4})$, write $\sigma=\Re\rho$ and use only $0<\sigma<1$. The Tate factor at $s=h+\rho$ is $\zeta^{(2)}(\rho-v)\ne0$. Indeed $\zeta(\rho-v)\ne0$ by the functional-equation argument above, and the modulus of $2^{-(\rho-v)}$ is $2^{v-\sigma}>1$, so the removed Euler factor is also nonzero. Hence
\[
 \operatorname{ord}_{s=h+\rho}\zeta_{X^-_\Phi}(s)
              =N\operatorname{ord}_{s=\rho}L(s,\chi_{-4}).
\]
This order identity also needs no A; A identifies the shifted critical line. No assumption that the chosen zero is simple is made. Locating its line does not evaluate its multiplicity in this constructed zeta function.

At the fixed prime $p=3$, where $\chi_{-4}(3)=-1$, one obtains
\[
       N=\frac{3^{v+h}-\#X^-_\Phi(\mathbb F_3)}{3^h}.
\]
The point-count function is in $\#\mathrm P$: guess the finitely many ternary coordinates using binary codes, reject unused codes, and check the one cubic equation. It is counting-hard by the displayed numerical reduction. Its parity equals $1+N\bmod2$, so even parity counting over this fixed field is $\oplus\mathrm P$-complete on the promised class. Rational-point existence is always true, because $N\leq2^v<3^v$ gives $\#X^-(\mathbb F_3)>0$. The hard Boolean question is whether the sign-Tate summand has nonzero rank, not whether the entire smooth scheme has a rational point.

E evaluates $N$ from the original CNF with its original $n$-variable bound and polynomial bit overhead for the displayed arithmetic. A makes the nontrivial zero-location statements unconditional under the user's designated input convention; it does not compute $N$. M is unused because the only nontrivial arithmetic character is the primitive quadratic Dirichlet character modulo four.

Finally consider effectivity. A single point count at $q=2$ recovers
\[
          N=\frac{\#X(\mathbb F_2)-2^{v+h}}{2^h}.
\]
Counting those points is in $\#\mathrm P$ by guessing the $v+h+1$ bit coordinates, and is $\#\mathrm P$-hard by this numerical reduction. Since $v\geq1$, its residue modulo $2^{h+1}$ is $2^h(N\bmod2)$, whereas every lower binary digit vanishes. At any odd prime power $q$, instead, $\#X(\mathbb F_q)\equiv1+N\pmod2$. In particular parity of rational-point counting is $\oplus\mathrm P$-complete already for these smooth geometrically integral sparse cubic hypersurfaces over the fixed field $\mathbb F_3$. The multiplicity $N=\dim H_c^{2h}$ is a $\#\mathrm P$ function on the encoded class; its nonvanishing is NP-complete and its parity is $\oplus\mathrm P$-complete. Rational-point existence for $X$ itself is always true, because $\#X\geq2^{v+h}$.

There is also an exact fixed-modulus transfer. Fix an integer $m_0\geq2$ and a prime $p$ not dividing $m_0$. Since $p^h$ is a unit modulo $m_0$,
\[
 N\bmod m_0=(p^h)^{-1}\bigl(\#X(\mathbb F_p)-p^{v+h}\bigr)\bmod m_0.
\]
Conversely the same point-count identity computes the point count modulo $m_0$ from $N\bmod m_0$. Thus the two modular counting functions are polynomial-time equivalent on the encoded class. This is stated for the function returning the count residue; no convention-dependent identification with a class called $\operatorname{Mod}_{m_0}\mathrm P$ is needed. If instead $p$ divides the modulus, division by $p^h$ is not invertible at that precision. In particular modulo $p^r$ with $r\leq h$ the whole count vanishes, whereas knowing the count modulo $p^{h+r}$ and subtracting the top term recovers $N\bmod p^r$ after exact division by $p^h$. The support of the whole point count is always nonzero; it does not supply the separate hard test $N>0$ for the lower cohomological multiplicity.

E computes $N$ through the source formula in $\operatorname{poly}_k(n+m)2^{O_k(n\log\log n/\sqrt{\log n})}$ time, with the sharper exponent when $m=O(n)$. Once $N$ is known, every displayed count and compact zeta-factor description is obtained by polynomial-bit arithmetic; a count at $q$ has $O((v+h)\log q+n)$ bits. Substituting the geometric coordinate count into E would be a weaker parameter statement. The elementary smooth construction is a classical quadric/Artin--Schreier form of the Cayley counting mechanism; the new conclusion here is its explicit arithmetic, pure-sign, and original-witness effectiveness gate, not a claim to have discovered the underlying point-count identity.
\end{proof}
\paragraph{Status and complexity impact.}
Complete integral smoothness, stratification, global cohomological projector, quadratic-character and arithmetic-zeta proofs; complete deduction-level self-reading passed. Independent complete deduction-proof review by root: passed; imported input theorems are not claimed independently proved. The exact obstruction is now a multiplicity in a genuine pure Tate cohomology group of a smooth geometrically integral arithmetic scheme. Its degree is at most three, its defining polynomial is sparse, and its two global zeta descriptions use Riemann-zeta shifts and the quadratic Dirichlet factor covered by A. This proves that those geometric and analytic structures can coexist with a counting-complete multiplicity-extraction problem. It does not prove that polynomial-time extraction is impossible: such an algorithm would give $\mathrm{FP}=\#\mathrm P$ and $\mathrm{PH}=\mathrm P$, while parity extraction would give $\mathrm P=\oplus\mathrm P$ and $\mathrm{NP}=\mathrm{RP}$. M is unnecessary for these objects, but that fact does not remove the general Dedekind witness-field gate in Burgisser or Koiran. No GCT canonical-basis identification, positive polytope formula, or VP/VNP separation follows without an additional effective representation-theoretic construction.

\paragraph{Geometric and representation dictionary.}
Let $f_1,\ldots,f_h\in\mathbb Z[x_1,\ldots,x_v]$ be the sparse quadratic Boolean and verifier equations in the proof, so that $v,h\geq1$, $f_1=x_1^2-x_1$, and $Y=V(f_1,\ldots,f_h)\simeq\coprod_N\operatorname{Spec}\mathbb Z$. Then the sparse cubic arithmetic hypersurface
\[
 X=V\left(t^2-t-\sum_{i=1}^h y_i f_i(x)\right)
             \subset\mathbb A^{v+h+1}_\mathbb Z
\]
is smooth of relative dimension $v+h$ with geometrically integral fibers. Its class is $[X]=\mathbb L^{v+h}+N\mathbb L^h$. For every field $K$ and every invertible prime $\ell$, its actual geometric compact-support groups, with their Galois actions, are
\[
 H_c^{2h}(X_{\overline K},\mathbb Q_\ell)=\mathbb Q_\ell(-h)^N,
 \qquad
 H_c^{2(v+h)}(X_{\overline K},\mathbb Q_\ell)=\mathbb Q_\ell(-v-h),
\]
and all other groups vanish. When $K$ is finite the displayed identifications are Frobenius-equivariant; they do not require finite $K$ to hold as Galois representations. Its arithmetic zeta is
$\zeta_X(s)=\zeta(s-v-h)\zeta(s-h)^N$.
These are separate proved assertions: affine-bundle stratification proves the scissor identity, while !-base change, the global trace and the involution $t\mapsto1-t$ prove the actual sheaf decomposition. Their exact interfaces are SGA4 XVII5.2.6/5.2.8 and XVIII2.8.1/2.9. The relative Jacobian identity proves smoothness over the full integer base, using SGA1 II4.10(iii)/4.12.

Over $B=\mathbb Z[1/2]$, replace the equation by $t^2+1=\sum_i y_i f_i(x)$. The resulting $X^-$ is also smooth with geometrically integral fibers. Its lower compact-support group is $(\chi_{-4}\otimes\mathbb Q_\ell(-h))^N$, its upper group is the same Tate line, and
\[
 \zeta_{X^-}(s)=\zeta^{(2)}(s-v-h)L(s-h,\chi_{-4})^N.
\]
Here the sign is the actual anti-invariant summand of the regular permutation representation $1\oplus\chi_{-4}$ of $\operatorname{Spec}\mathbb Z[i,1/2]\to\operatorname{Spec}B$. The first construction has a separate algebraic involution; the second realizes a quadratic Galois sign. Neither asserts an injective realization of the full split-zero semiring.

The proof constructs these data from a Boolean formula with $N$ satisfying assignments, with $v,h$ and the cubic description polynomial in that formula's size. Consequently extracting the lower cohomology multiplicity $N$ from this succinct geometric input is counting-complete, its parity is parity-counting-complete, and its nonvanishing is NP-complete. The corresponding shifted nonreal nontrivial zeta-zero orders are $N$ times a fixed nonzero analytic multiplicity; this identity is unconditional, while A locates the shifted critical lines. E evaluates $N$ with its original-witness bound. Thus actual purity, a concrete sign representation, smooth integral geometry and A-controlled analytic zero locations all coexist with this precise multiplicity-extraction problem. No unconditional time lower bound or effective GCT canonical-object identification is inferred.

\subsection{NN-C028 --- Embedding comparison for actual pure Tate cohomology is PP-complete}
\begin{proposition}[NN-C028]
There is a polynomial-time construction from a pair of Boolean verification circuits to two smooth geometrically integral affine schemes $X_1,X_2/\mathbb Z$, each defined by one sparse cubic equation. They have equal relative dimension $V+H$, and
\[
 H_c^{2H}(X_{j,\overline{\mathbb F}_p},\mathbb Q_\ell)
          =\mathbb Q_\ell(-H)^{N_j},\qquad
 H_c^{2(V+H)}=\mathbb Q_\ell(-V-H),
\]
with all other groups zero and $N_j$ the corresponding acceptance count. For any fixed prime $p$ and any fixed $\ell\ne p$, deciding whether a Frobenius-equivariant linear injection
\[
 H_c^{2H}(X_1,\mathbb Q_\ell)
       \hookrightarrow H_c^{2H}(X_2,\mathbb Q_\ell)
\]
exists is PP-complete on this encoded family. Hardness persists when $N_1,N_2>0$, so the Frobenius eigenvalue support, the cohomological degree support and the two weights agree. The same comparison holds simultaneously at every prime, and for the literal quadratic-character variants over $\mathbb Z[1/2]$ at every odd prime.

This is existence of a linear map between actual geometric cohomology representations. It is not a claim that the map is induced by a morphism, correspondence, or containment between the schemes. Under A the global zeta factors have the same shifted critical lines; their multiplicity comparison is still the PP-complete operation.
\end{proposition}
\paragraph{Exact inputs.}
NN-C027. smooth integral sparse cubic construction and actual compact-support Tate decomposition; both arithmetic zeta identities.

NN-G015. work/gct/chapter.tex subsectionNN-G015; support-versus-count comparison is background, not re-registered.

NN-C006. Toda main-theorem locator and resulting PH collapse.

Aaronson, Quantum Computing, Postselection, and Probabilistic Polynomial-Time (2005). Author manuscript Section1, printedp1, definition of PP by strict-majority acceptance; the reduction below is proved directly.

SGA4, ExposesXVII--XVIII. XVII Theorem5.2.6 and Proposition5.2.8; XVIII equation2.8.1 and Theorem2.9; used through the complete global sheaf proof in NN-C027.

E: The Exponential Time Hypothesis is False, supplied LaTeX. Theorem1.1, original-CNF-variable count and bounded-density specialization.

\begin{proof}
Start with two verification circuits. The gate construction of NN-C027 gives systems $f_{j,1},\ldots,f_{j,h_j}$ on $v_j$ Boolean input and gate coordinates, with
\[
 Y_j=V(f_{j,1},\ldots,f_{j,h_j})\simeq
             \coprod_{N_j}\operatorname{Spec}\mathbb Z.
\]
Each gate extension is unique. The equations have degree at most two and constant monomial sparsity, and each system retains $x_1^2-x_1$ among its equations. A fixed dummy coordinate constrained to zero can be introduced first if necessary to make every coordinate list nonempty.

The geometric parameters can be matched without multiplying either count. Put $V=\max(v_1,v_2)$. For each missing coordinate append a new variable $z$ and the two equations $z^2-z=0$, $z=0$. These force its unique value zero, so $Y_j$ and $N_j$ are unchanged. Let the resulting numbers of equations be $h'_1,h'_2$, and put $H=\max(h'_1,h'_2)$. Pad the shorter list by repeated copies of its retained Boolean equation until both lists have exactly $H$ entries. Repeating an equation does not change its common-zero scheme. All of this uses only polynomially many coordinates, equations and constant-size coefficients. Apply NN-C027 to the padded lists:
\[
 X_j:\quad t^2-t=\sum_{i=1}^H y_i f_{j,i}(x).
\]
Its proof applies unchanged to redundant equations and gives the displayed actual cohomology. In particular the lower Frobenius operator is the scalar $p^H$ on a vector space of dimension $N_j$. Every linear map between these two lower groups therefore commutes with Frobenius. Such a map can be injective exactly when
\[
                         N_1\leq N_2.
\]
This proves the assertion simultaneously at all primes; no explicitly listed exponential-size cohomology basis is required in the input. In the quadratic-character version the common scalar is $\chi_{-4}(p)p^H$, so the same equivalence holds at every odd prime. The equal top Tate lines can also be included if one asks for an injection of the full graded compact-support representation.

For hardness, take a Boolean circuit $C$ on $n\geq1$ bits, with $a$ accepting assignments. On a tagged input $(b,x)\in\{0,1\}^{n+1}$ define two polynomial-size verifiers. The right verifier accepts when $b=0$ and $C(x)=1$, or when $b=1$ and $x=0^n$. Its count is $N_R=a+1$. The left verifier accepts when $b=0$ and $x_1=0$, or when $b=1$ and $x$ is one of the two distinct strings $0^n,e_1$. Its count is
\[
             N_L=2^{n-1}+2.
\]
Both counts are positive. The injection from the left lower group to the right exists precisely when
\[
  2^{n-1}+2\leq a+1
       \quad\Longleftrightarrow\quad a>2^{n-1}.
\]
Strict-majority circuit acceptance is PP-complete directly from the definition: unroll a polynomial-time probabilistic verifier on its fixed input, using its independent fair coin bits as circuit inputs, and pad the coin count if needed. The circuit accepts on exactly the machine's accepting coin strings. Thus the entire reduction to the two sparse cubics is polynomial and proves PP-hardness, including the positive-rank restriction. An input circuit with no coin bits can be given an ignored bit without changing its majority status.

For completeness, here is a direct PP-membership construction that handles equality rather than silently replacing $\leq$ by $<$. Let $P,Q$ be the two acceptance predicates, with counts $A,B$, padded to a common witness length $w$ by requiring every added bit to be zero. Add a new bit $r$ and define
$P'(r,x)=[r=0]P(x)$, $Q'(r,x)=[r=0]Q(x)$.
Their common domain has size $M=2^{w+1}$, while their counts remain $A,B$. On a selector bit $b$ and this domain, accept when
\[
          (b=0\ \hbox{and }Q'=1)
            \quad\hbox{or}\quad(b=1\ \hbox{and }P'=0).
\]
There are $M+B-A$ accepting strings among $2M$. The string $b=0,r=1,x=0^w$ is always rejecting. Add a fresh free bit $u$, doubling all paths, and override exactly the one path $u=0,b=0,r=1,x=0^w$ to accept. The number of accepting paths becomes
\[
                   2M+2(B-A)+1
\]
among $4M$. This exceeds half exactly when $B-A\geq0$. The whole test is a polynomial-size Boolean predicate. It proves PP membership on the encoded family, and hence PP-completeness. The circuit-and-equation format can be checked syntactically; the assertion is not a promised PP upper bound for arbitrary sparse cubic schemes whose provenance is absent.

The associated arithmetic zeta functions are
\[
             \zeta_{X_j}(s)
                 =\zeta(s-V-H)\zeta(s-H)^{N_j}.
\]
When $N_1,N_2>0$ the types of nontrivial factors, cohomological degrees and eigenvalue supports agree. At any fixed nonreal nontrivial zero $\rho$ of $\zeta$, NN-C027 gives
\[
 \operatorname{ord}_{s=H+\rho}\zeta_{X_j}(s)
                 =N_j\operatorname{ord}_\rho\zeta.
\]
Consequently the same embedding comparison is comparison of those two analytic orders. The order identity is unconditional; A fixes the common shifted critical lines. For the quadratic variant replace the lower factor by $L(s-H,\chi_{-4})$ and remove the upper prime-two Euler factor as in NN-C027. The same conclusion holds at its nonreal nontrivial zeros.

E decides the comparison by computing the two original counts and comparing their binary integers. For fixed-width source CNFs it keeps the original variable parameters, regardless of the padding of the geometric schemes. For a general verification circuit of size $s$ on $n$ inputs, the parsimonious gate-to-3-CNF encoding has $O(n+s)$ variables and a linear number of clauses, so Theorem1.1's bounded-density form gives a bound $\operatorname{poly}(n+s)2^{O((n+s)/\sqrt{\log(n+s)})}$ for that source circuit. No bound solely in the original coin count $n$ is inferred for arbitrary circuits with many gates.
\end{proof}
\paragraph{Status and complexity impact.}
Complete actual-cohomology comparison, parameter matching and strict-majority proofs; complete deduction-level self-reading passed. Independent complete deduction-proof review by root: passed; imported input theorems are not claimed independently proved. This turns NN-G015's abstract support-versus-dominance example into a comparison of actual pure geometric representations of smooth integral cubic schemes. The representation type and analytic critical lines are fixed; deciding which multiplicity is larger is PP-complete. A polynomial-time injection-existence procedure on this succinct family would give $\mathrm P=\mathrm{PP}$, and hence $\mathrm{PH}=\mathrm P$ by Toda. This is a conditional implication, not an unconditional lower bound. E supplies a subexponential evaluator of the comparison, and A does not evaluate the multiplicities. The result does not transfer PP-hardness to the permanent/determinant orbit-closure multiplicities or establish a GCT multiplicity obstruction.

\paragraph{Geometric and representation dictionary.}
The algebraic-geometric assertion is that, for the explicitly padded Boolean systems $Y_j\simeq\coprod_{N_j}\operatorname{Spec}\mathbb Z$, the two cubic schemes
\[
 X_j=V\left(t^2-t-\sum_{i=1}^H y_i f_{j,i}(x)\right)
          \subset\mathbb A^{V+H+1}_\mathbb Z
\]
are smooth with geometrically integral fibers, and their actual lower compact-support realizations are $\mathbb Q_\ell(-H)^{N_j}$. In the category of rational l-adic Frobenius representations, an injection between these realizations exists if and only if $N_1\leq N_2$. This equivalence holds for every prime with identical geometric parameters, and for the anti-invariant quadratic-character realizations at every odd prime. The scheme construction, relative trace and projectors are those proved in NN-C027, with exact SGA4 XVII5.2.6/5.2.8 and XVIII2.8.1/2.9 interfaces. Parsimonious padding proves that the geometric dictionary preserves the relevant counts and common Tate twist.

Thus even on this actual pure Tate subcategory, a succinct defining-equation presentation can make the existence of a subrepresentation PP-complete. The explicit matrices, once supplied, would be scalar and their rank comparison elementary; obtaining their multiplicities from the equations is the hard operation. Their zeta factors have A-controlled shifted critical lines, but the orders at a fixed nonreal lower-factor zero give precisely the same comparison. This is a theorem about linear injections of cohomology realizations, not about scheme morphisms, algebraic correspondences or orbit-closure containment. No fullness statement for a realization functor is assumed.

\subsection{NN-C029 --- E constructs subexponential ordinary ideal certificates and affine trivializations}
\begin{proposition}[NN-C029]
Fix a density constant $C$. Let $f_1,\ldots,f_h\in\mathbb Z[x_1,\ldots,x_v]$, $h\leq Cv$, be the Boolean, AND/NOT-gate and acceptance equations of NN-C027. In particular all $x_j^2-x_j$ are retained, each equation has degree at most two, scope at most three and constant coefficient size, and its Boolean values belong to $\{0,1,-1\}$. If the system is unsatisfiable, E deterministically constructs division-free arithmetic circuits over $\mathbb Q$ for integral polynomials $a_i(x)$ satisfying the ordinary polynomial identity
\[
                  \sum_{i=1}^h a_i(x)f_i(x)=1,
          \qquad \deg a_i\leq v.
\]
The total circuit size and construction time are
\[
 S(v)=\operatorname{poly}(v)\,
 2^{O_C(v(\log\log v)^2/\sqrt{\log v})},
\]
and every rational constant has polynomial bit length. Thus $\sum_i a_i(x)z_i$ is a Hilbert-like IPS certificate with a constructible $S(v)$-size proof over $\mathbb Q$. Its coefficient polynomials are integral; the circuit representation is allowed rational constants.

For fixed-width $k$-CNFs with $m\leq Cn$ clauses there is a stronger specialization on the original $n$ input variables: writing $g_i$ for the product of the clause's falsifying literals, E constructs integral $A_i$ with $\sum_i A_i g_i=1$, without Boolean axioms, in size and time $S(n)$, with constants also depending on $k$. The first-violated-clause identity itself is classical Grochow--Pitassi Lemma4.2; the deduction here is the E-based circuit bound and its ordinary-ideal extension to the quadratic verifier presentation.

In the empty case, E also constructs $S(v)$-size rational arithmetic circuit descriptions of mutually inverse integral polynomial maps
\[
 X=\left\{t^2-t=\sum_i y_i f_i(x)\right\}
       \simeq\mathbb A^{v+h}_{\mathbb Z}.
\]
Their degrees are at most $v+2$ and $v+6$. For a verifier on $n$ input bits with $s$ gates, the parameter is $v=n+s$, up to harmless constant additions. No original-$n$ bound for arbitrary long verifier circuits is asserted.

There is an explicit large-characteristic transfer. Put
$B(v)=\max(2,\lceil\log_2 v\rceil)$. The multiplier circuits can be chosen with constants in $\mathbb Z[1/B(v)!]$. Hence they give circuits of the same size over every field of characteristic $p>B(v)$, with construction bit overhead polynomial in $\log p$. The same statement holds for C016's universal partition polynomials and for the sparse-CNF specialization with $v$ replaced by $n$. This original localization argument does not cover fixed small characteristic as the input size tends to infinity. NN-C033 subsequently closes that limitation by a modified integral-swap-matrix and extension-node construction; NN-C034 then constructs integral circuits over Z itself.
\end{proposition}
\paragraph{Exact inputs.}
NN-C016. complete coefficient-independent symbolic partition construction, fixed arity and linear factor density.

E: The Exponential Time Hypothesis is False, supplied LaTeX. Proposition3.4 over every commutative Q-algebra; Lemma4.6 replacement-row denominator (R-1)!Q!; Lemmas4.7--4.8 and equation parameter-QR; AppendixB Hermite matrix in subsection rational-readout; coefficient-independent measurement separation.

Grochow--Pitassi, Circuit Complexity, Proof Complexity, and Polynomial Identity Testing: The Ideal Proof System. Author JACM manuscript, Definition1.1 and Hilbert-like terminology, printed draftp1:4; Lemma4.2 and complete proof, printed draftpp1:26--1:28, equations(4)--(6).

EGA IV, fourth part, Publications Mathematiques IHES32 (1967). Printed p.111,18.2.1--18.2.2; base-change equation18.2.1.2; primary PDF page110 (zero-based109).

NN-S016. ordinary integral Bezout identity and explicit affine-space coordinate formulas; work/complexity/c027\_degree\_addendum.tex and work/root/chapter8.tex.

NN-C027. quadratic AND/NOT verifier equations with all Boolean axioms and the associated smooth integral cubic.

\begin{proof}
Bounded values of the parameters can be handled by direct expansion. We first prove the ordinary-ideal assertion for the quadratic verifier list. For a bit vector $b$ put
\[
 e_b(x)=\prod_{j:b_j=1}x_j\prod_{j:b_j=0}(1-x_j),
 \qquad
 a_i^{(0)}(x)=\sum_{b\in\{0,1\}^v}e_b(x)f_i(b)
                  \prod_{j<i}(1-f_j(b)^2).
\]
The factors $1-f_j(b)^2$ are exactly the indicators of $f_j(b)=0$. At any bit vector $b$, unsatisfiability guarantees a first index $i$ with $f_i(b)\ne0$. Its coefficient in the displayed sum is $f_i(b)\in\{1,-1\}$, and all other first-violation coefficients vanish. Therefore
\[
                    \sum_i a_i^{(0)}(b)f_i(b)=1.
\]
Each $a_i^{(0)}$ is an integral multilinear polynomial of degree at most $v$. The Boolean equations themselves have zero values and hence zero preliminary multipliers.

The preliminary multipliers are symbolic bounded-arity partition functions. Their summation variables are the $v$ bits $b$. There are at most $h$ prefix factors with tables $1-f_j(b)^2$, one local factor with table $f_i(b)$, and $v$ unary factors with respective weights $x_l$ at $b_l=1$ and $1-x_l$ at $b_l=0$. Every scope has size at most three and there are $O(v)$ factors and incidences. The fact that the unreduced expression $1-f_j^2$ can have degree four is harmless: its Boolean table, interpolated on at most three variables, has degree at most three. The coefficient algebra is $\mathbb Q[x_1,\ldots,x_v]$, with the $x_l$ treated as formal coefficients rather than extra summation variables.

Apply NN-C016 separately to these $h$ partition functions and specialize its independent factor weights to the indicated constants or linear polynomials. These are constant-size substitutions, so they preserve the bound $S(v)$. The factor count, rather than an unproved bound on each variable's incidence, supplies the hypothesis: the total number of incidences is $O(v)$, and C016 explicitly branches on high-incidence variables. Multiplying by $h=O(v)$ changes only the polynomial prefactor. Its construction makes no identity tests on the coefficient algebra and introduces no division by a polynomial in the $x_l$. Its only inversions are of explicit rational interpolation constants, whose bit lengths are polynomial. This yields circuits for every $a_i^{(0)}$ within the claimed time and size.

An ordinary polynomial identity still requires Boolean-ideal correction; evaluating correctly on the cube alone is insufficient. Here locality makes that correction cheap without globally expanding the circuit. Let $I_i$ be the scope of $f_i$ and $t_i=|I_i|\leq3$. Since $a_i^{(0)}$ is multilinear, its local interpolation is
\[
 a_i^{(0)}(x)=\sum_{\sigma\in\{0,1\}^{I_i}}
        A_{i,\sigma}(x_{\overline I_i})e_\sigma(x_{I_i}),
 \qquad A_{i,\sigma}=a_i^{(0)}\big|_{x_{I_i}=\sigma}.
\]
At most eight circuit specializations compute these polynomials, and their degrees are at most $v-t_i$. In the constant number of local variables the polynomial
\[
       e_\sigma f_i-f_i(\sigma)e_\sigma
\]
has zero multilinear remainder. Its degree is at most $t_i+2$. Explicit integral monomial division therefore gives
\[
 e_\sigma f_i-f_i(\sigma)e_\sigma
      =\sum_{j\in I_i}r_{i,\sigma,j}(x_{I_i})(x_j^2-x_j),
      \qquad \deg r_{i,\sigma,j}\leq t_i.
\]
For completeness, this division uses only
\[
 x^\alpha-x^{\epsilon}=\sum_j
  \left(\prod_{l<j}x_l^{\epsilon_l}\right)
  (x_j^{\alpha_j}-x_j^{\epsilon_j})
  \left(\prod_{l>j}x_l^{\alpha_l}\right),
 \quad \epsilon_l=\min(\alpha_l,1),
\]
and $x_j^a-x_j=(x_j^2-x_j)\sum_{r=0}^{a-2}x_j^r$ for $a\geq2$.
There are constantly many local monomials and coefficients of bounded size. Thus all $r_{i,\sigma,j}$ have constant-size integral circuits. Empty scopes give no correction.

Write $\mathrm{ML}$ for the unique Boolean multilinear remainder and put
\[
 Q_{i,j}=\sum_{\sigma\in\{0,1\}^{I_i}}
                       A_{i,\sigma}r_{i,\sigma,j}.
\]
The local identities imply
\[
 a_i^{(0)}f_i-\mathrm{ML}(a_i^{(0)}f_i)
              =\sum_{j\in I_i}Q_{i,j}(x)(x_j^2-x_j).
\]
The asserted remainder really is multilinear: its local factor is
$\sum_\sigma f_i(\sigma)e_\sigma$, with the corresponding coefficients $A_{i,\sigma}$ involving only the other variables. Its values are those of $a_i^{(0)}f_i$. Each $Q_{i,j}$ has degree at most $(v-t_i)+t_i=v$. Summing over $i$, the sum of the remainders is the unique multilinear polynomial whose Boolean values are all one, hence is identically one over $\mathbb Z$. Consequently
\[
 1=\sum_i a_i^{(0)}f_i-
             \sum_{j=1}^v\left(\sum_i Q_{i,j}\right)(x_j^2-x_j).
\]
Subtract the indicated $Q$ sum from the multiplier of each retained Boolean equation. This produces the promised integral $a_i$ and an ordinary identity in $\mathbb Z[x]$. Constantly many specializations and local products per equation, followed by $O(v)$ additions, keep the total circuit size and construction time within $S(v)$. In particular, no exponential multilinear expansion is being silently treated as a polynomial-time operation. The preliminary coefficient norm is bounded by $3^v$ in total, and the local correction introduces only a fixed factor per equation, so the resulting integer coefficients have $O(v+\log(v+1))$ bits; integrality is a property of the output polynomials, not of every rational circuit gate.

This construction can be made unconditional as an algorithm that either reports satisfiability or produces the certificate. Each local relation $f_i(b)=0$ is a constant-scope Boolean table, convertible to at most eight forbidden-assignment clauses. The resulting fixed-width CNF has the same $v$ variables and $O(v)$ clauses. E first decides it within its bounded-density running time, which is absorbed by $S(v)$. When it is satisfiable the algorithm reports that outcome; when it is not, the displayed construction applies. The circuit $C(x,z)=\sum_i a_i(x)z_i$ satisfies $C(x,0)=0$ and $C(x,f(x))=1$, exactly Definition1.1 of IPS, and is linear in the placeholders. Its actual polynomial degree is at most $v+1$. If a syntactic variable-degree bound is desired, retain homogeneous components $0,\ldots,v$ at every gate of each multiplier circuit: addition is componentwise and multiplication uses the truncated convolution. There is no variable division, so discarded higher components cannot contribute to a later lower component. Since the output degree is at most $v$, their sum is exactly the same multiplier. This uses $O(v^2 S(v))$ gates and no identity test, with syntactic degree at most $v$ when rational constants have variable-degree zero. It is not the different constant-free formal-degree convention in which constant leaves have degree one.

For the stronger sparse-CNF specialization, delete tautological clauses and repeated literals. An empty clause has the immediate certificate $1\cdot1=1$. Otherwise let $I_i$ be the distinct-variable scope of clause $i$, let $\alpha_i$ be its unique falsifying assignment, and set $g_i=e_{\alpha_i}(x_{I_i})$. Define
\[
 A_i(x_{\overline I_i})=
    \sum_{b\in\{0,1\}^{\overline I_i}}
       e_b(x_{\overline I_i})
       \prod_{j<i}\bigl(1-g_j(b,\alpha_i)\bigr).
\]
Upon multiplying by $g_i$, this is precisely the sum of the full idempotents $e_c(x)$ over assignments whose first violated clause is $i$. For an unsatisfiable CNF these sets partition the entire cube, so
\[
              \sum_i A_i g_i=\sum_c e_c
                   =\prod_{l=1}^n(x_l+(1-x_l))=1
\]
as an ordinary integral polynomial identity. In particular $\deg A_i\leq n-|I_i|$ and every product $A_i g_i$ has degree at most $n$. This is the classical identity of Grochow--Pitassi Lemma4.2, equations(4)--(6), specialized by fixing the falsifying local assignment. Its E consequence is as follows. Each $A_i$ is a partition function on $n-|I_i|$ summation bits, with at most $m+n$ factors of arity at most $k$. For fixed $k,C$ and $n>2k$, its variable count is at least $n/2$ and the factor density remains bounded. C016 therefore constructs all of them in size and time $S(n)$. Small $n$ is handled directly. No Boolean correction, auxiliary gate variable, or expansion in the original $n$ variables is required in this specialization. For dense CNFs this particular argument does not preserve the same original-$n$ bound; compiling into verifier gates changes the relevant parameter to their total number. NN-C036 subsequently supplies an original-$n$ integral circuit certificate for arbitrary-density fixed-width CNFs by slower disjoint sparsification, with its explicitly weaker exponent.

We next prove the asserted large-characteristic transfer, including a uniform denominator bound across every high-incidence branch. In a branch with $u$ remaining summation variables, E Lemma4.8 uses
\[
 \lambda=\lceil\log_2u\rceil,\quad
 \nu=\lceil\log_2\lambda\rceil,\quad
 Q=\left\lfloor\frac{\lambda}{128\nu}\right\rfloor,
\]
and every successful layout has $R<Q$. Since $u\leq v$, both $Q$ and $R$ are at most $B(v)$. The high-incidence branching introduces no new summation variables. For fixed density and arity, only finitely many input sizes fail the layout's numerical checks; direct expansion handles them using only integral constants. The branch sizes are $v-O(v/\sqrt{\log v})$, so this fallback affects only a fixed finite range and preserves the asymptotic bound. Thus the single localization
\[
                       A_v=\mathbb Z[1/B(v)!]
\]
works for every branch, preliminary multiplier and local correction. In the sparse-CNF specialization each summation dimension is at most $n$, so $A_n$ works uniformly there.

Here is why all arithmetic-circuit constants may be chosen in that localization. E Lemma4.6 gives the exact common denominator $(R-1)!Q!$ for every replacement row. Its prime divisors are at most $B(v)$, so these row coefficients lie in $A_v$. Taylor substitution and Boolean differences introduce only binomial integers and signs. Factor interpolation on the Boolean cube uses only zero, one and minus one. All subsequent sums and products therefore remain in $A_v$ until the final Hasse-Hermite interpolation.

That interpolation also needs only this localization. Let the multiplicity of each node be $J$ and put $M=QJ$. The evaluation-of-jets map is
\[
 A_v[T]\big/\!\left(\prod_{t=1}^{Q}(T-t)^J\right)
       \longrightarrow\prod_{t=1}^{Q} A_v[z]/(z^J),
 \qquad P\longmapsto(P(t+z))_{t=1}^{Q}.
\]
For two distinct nodes, their nonzero difference is an integer of absolute value less than $Q$, and is a unit in $A_v$. The corresponding linear ideals are comaximal, and so are their $J$th powers: expanding a relation $a+b=1$ to exponent $2J-1$ expresses one as a sum of an element of each power. The Chinese remainder theorem therefore makes this map an isomorphism. Monic division gives the free basis $1,T,\ldots,T^{M-1}$ on the left; the right uses the local powers $1,z,\ldots,z^{J-1}$. In these bases its matrix is exactly E's Hasse-Hermite matrix
\[
  \mathsf V_{(t,j),u}=\binom uj t^{u-j}\quad(u\geq j),
               \qquad \mathsf V_{(t,j),u}=0\quad(u<j).
\]
Its inverse thus has entries in $A_v$. Compute those final entries over $\mathbb Q$ as before and put them in the circuit as constants; intermediate elimination pivots are not circuit gates. Their polynomial bit-height bound is unchanged. This reasoning does not divide by $j!$ and imposes no bound $J<p$.

Consequently the C016 circuit is a circuit over $A_v$ for a polynomial with integer coefficients in its independent weights. Equality with that universal partition polynomial holds in $\mathbb Q[y]$, and hence in $A_v[y]$ by the injective inclusion into $\mathbb Q[y]$. For every prime $p>B(v)$ there is a homomorphism $A_v\to\mathbb F_p$. Apply it to every constant and gate: equality survives, and integer-weight specializations, local corrections and homogeneous-component constructions survive as well. The resulting output is exactly the reduction of each integral multiplier, not merely a polynomial agreeing with it on a finite set. Reducing polynomial-bit numerators and denominators and inverting their nonzero residues uses polynomial bit cost per constant in $v+\log p$. The compiled circuit uses only prime-field constants and is valid over every extension field of that characteristic. This proves the size and construction claims. No procedure choosing a prime, and no assertion for a fixed prime below the threshold, is hidden in the argument.

Finally apply the ordinary identity to the smooth cubic of NN-C027. Retain $x$ and define
\[
 z_i=y_i+a_i(x)(2t-t^2).
\]
On $X$, the scalar product $\sum_i f_i z_i$ equals
$(t^2-t)+(2t-t^2)=t$. Conversely, for free coordinates $(x,z)$ set
\[
 t=\sum_i f_i(x)z_i,
 \qquad y_i=z_i+a_i(x)(t^2-2t).
\]
Then $\sum_i f_i y_i=t+(t^2-2t)=t^2-t$, and direct substitution shows both composites are identities. These are integral polynomial maps over $\mathbb Z$. The first has degree at most $v+2$; the second has degree at most $v+6$, because $t$ has degree at most three. Shared circuit substitution adds only polynomially many gates to the multiplier circuits. This proves the effective geometric assertion. The rational syntax in this original proof describes integral polynomial maps; it is not a license to reduce a rational constant modulo a prime dividing its denominator. The later integral compilation in NN-C034 supplies a different circuit valid after arbitrary base change.
\end{proof}
\paragraph{Status and complexity impact.}
Complete construction of ordinary ideal multipliers, local correction, sparse-CNF specialization and integral affine trivialization. Independent complete deduction-proof review by root: passed; imported input theorems are not claimed independently proved. This closes a specific algebraic-certificate gate left by NN-C009: for bounded-density local verifier systems, E yields actual ordinary ideal multipliers with subexponential rational arithmetic circuits. It strengthens the transcript conclusion of NN-C008 and makes NN-S016's empty-case affine trivialization effective in this circuit model. It gives neither polynomial-size certificates nor a polynomial-time verifier for arbitrary supplied IPS circuits; verification still presents polynomial identity tests. It therefore implies no $\mathrm{NP}=\mathrm{coNP}$ collapse and no VP=VNP conclusion. The exponential monomial-size polynomial-calculus lower bounds in NN-C021 remain compatible with these shared circuits. The denominator audit here gives circuits in every characteristic greater than $\max(2,\lceil\log_2 v\rceil)$. NN-C033 subsequently supplies a different construction over every prime field, including constant-free circuits of polynomial formal degree. Semantic integrality alone did not justify reducing the original rational constants at denominator primes. NN-C034 subsequently closes the integral-gate question over $\mathbb Z$ by cyclotomic interpolation and supplies constant-free circuits of polynomial formal degree.

\paragraph{Geometric and representation dictionary.}
The algebraic-geometric theorem is constructive in the following precise representation. For the split Boolean coordinate algebra
\[
 B=\mathbb Z[x_1,\ldots,x_v]/(x_j^2-x_j)
                 \simeq\prod_{b\in\{0,1\}^v}\mathbb Z,
\]
let $Y=V(f_1,\ldots,f_h)$ be the closed subscheme cut out by the local verifier equations. If $Y$ is empty, E produces rational arithmetic circuits representing integral polynomials $a_i$ of degree at most $v$ whose ordinary ideal combination is one in $\mathbb Z[x]$. The first-violation construction partitions the component idempotents of $B$; the constant-scope division above lifts that partition to the ambient affine coordinate ring. It thereby constructs an integral polynomial isomorphism
\[
 V\left(t^2-t-\sum_i y_i f_i(x)\right)
        \longrightarrow\mathbb A^{v+h}_{\mathbb Z},
 \quad (x,y,t)\longmapsto(x,y+a(2t-t^2)),
\]
with the explicitly proved inverse and degree bounds. Its circuit description has the displayed $S(v)$ bound. Moreover that description is defined over the explicit open arithmetic base $\operatorname{Spec}\mathbb Z[1/B(v)!]$: the Hasse-jet Chinese-remainder isomorphism proves the interpolation denominators are invertible there. Thus the same effective presentation of the ideal identity and affine trivialization base-changes to every field of characteristic $p>B(v)$. The integral morphisms exist at all primes; only their claimed arithmetic-circuit presentation has this restricted base.

The symbolic preliminary multipliers are finite-algebra traces of products of local weights, so C016's trace construction and base change apply, with exact classical locator EGA IV4, printedp111,18.2.1--18.2.2 and equation18.2.1.2. The passage from quotient identities to ordinary ideal multipliers is proved here by local interpolation, without invoking an unproved effectivity theorem. In the sparse-CNF case, the classical first-violation partition already gives the ordinary identity without the Boolean equations. These are integral geometric statements with localized arithmetic-circuit presentations; no efficient circuit presentation at every prime, motivic equivalence theorem, or GCT representation-theoretic lower bound is inferred.

\subsection{NN-C030 --- Pairwise homogeneity gives a deterministic smooth finite-field parity reduction}
\begin{proposition}[NN-C030]
The randomized finite-field specialization in NN-C024 can be replaced by a deterministic polynomial-time construction on the actual incidence-rail encoding of NN-C023. Given a Boolean formula, construct a smooth projective hypersurface over an explicitly presented finite field $\mathbb F_{2^b}$, where $b=O(r^2)$ and its factored degree-$r$ equation has polynomial description length. The top coherent cohomology has its known one-dimensional basis, and its $2^b$-power Frobenius scalar is exactly the original satisfying-assignment parity $c$.

More precisely, let $r=2M+1\geq3$ and let
\[
 F=z_0P(z_{1,0},z_{1,1},\ldots,z_{M,0},z_{M,1})
\]
be nonzero over $\mathbb F_2$, with $P$ homogeneous of degree two in every displayed pair. Put
\[
 H=z_0^r+\sum_{e=1}^M(z_{e,0}^r+z_{e,1}^r),
 \quad d=\operatorname{ord}_r(2),\quad D=dr.
\]
Every parameter $a$ for which $V(H+aF)$ is singular lies in $\mathbb F_{2^D}$. Choose a power of two $b>D$ with $b\leq2D$ and an explicitly constructed field element $\theta$ of absolute degree $b$. Then $V(H+\theta F)$ is smooth. Its top absolute Frobenius has semilinear coefficient $\theta c$, while its field-linear Frobenius scalar is $c$.

Consequently scalar extraction, and equivalently point-count parity on this encoded class, is parity-counting complete under deterministic polynomial-time reductions. A polynomial-time algorithm for either would give $\mathrm P=\oplus\mathrm P$ and $\mathrm{NP}=\mathrm{RP}$. E gives the inherited subexponential evaluator. This closes the deterministic finite-field smooth-specialization gate in NN-C024; it does not construct canonical GCT objects.
\end{proposition}
\paragraph{Exact inputs.}
NN-C023. incidence-rail factorization F=z0 P, with each rail pair of degree two and parity coefficient c.

NN-C024. complete quadratic Artin--Schreier field-tower construction and semilinear-versus-q-linear Frobenius computation.

SGA1, ExposeII, Section4. Theorem4.10(iii) and Corollary4.12, the homogeneous hypersurface Jacobian criterion used through NN-C024.

Katz, Algebraic Solutions of Differential Equations, Invent.Math.18(1972),1--118. Rank-one Special Case2.3.7.17, printedp36; exact Cech calculation supplied in NN-C023.

SGA 7 II, Expose XXII. Theorem 3.1 (3.1.1); Corollary 3.2 (3.2.1).

Deligne, La conjecture de Weil I, PMIHES43(1974),273--307. Theorem1.6, printedp276, smooth projective finite-field purity.

Valiant--Vazirani, NP is as easy as detecting unique solutions, TCS 47 (1986), 85--93. Theorem 1.1; Corollaries 3.1--3.4, especially Corollary 3.2 p.91.

E: The Exponential Time Hypothesis is False, supplied LaTeX. Theorem1.1, inherited original-formula evaluator.

\begin{proof}
We first prove the singular-parameter assertion for every nonzero $P$ with the stated pairwise homogeneity; no factorization or complexity assumption is needed for this geometric lemma. Work in a fixed algebraic closure of $\mathbb F_2$. The odd integer $r$ is one in this field, so the partial derivatives of $H$ are the corresponding $(r-1)$st powers. The zero parameter gives the smooth Fermat hypersurface. Suppose $H+aF$ has a singular projective point for some $a$.

At such a point, $z_0$ cannot be zero. If it were, all rail partial derivatives of $F=z_0P$ would vanish, and the zero-gradient equations would force every rail coordinate to be zero. This is not a projective point. Normalize $z_0=1$. The $z_0$ derivative now gives
\[
                              1+aP=0.
\]
Thus $a\ne0$, $P\ne0$, and $aP=1$ at the normalized point. If both coordinates of any rail pair were zero, pairwise degree two would force $P=0$. This too is impossible.

Euler's identity in one rail pair reads
\[
 z_{e,0}\frac{\partial P}{\partial z_{e,0}}
 +z_{e,1}\frac{\partial P}{\partial z_{e,1}}=2P=0.
\]
Multiply its two zero-gradient equations by their corresponding rail coordinates and add. The displayed identity yields
\[
                         z_{e,0}^r=z_{e,1}^r.
\]
Neither coordinate can be zero, since then both would be zero. Therefore
\[
              s_e=z_{e,0}\ne0,\qquad
              t_e=z_{e,1}/z_{e,0},\qquad t_e^r=1.
\]
Since $r$ is odd, $2$ is invertible modulo $r$ and its multiplicative order $d$ exists, with $d\leq\varphi(r)\leq r-1$. Put $q_0=2^d$. As $r$ divides $q_0-1$, every $r$th root of unity lies in the field $K=\mathbb F_{q_0}$. Thus all $t_e$ lie in $K$.

Pairwise homogeneity gives an identity
\[
 P(s_1,s_1t_1,\ldots,s_M,s_Mt_M)
                 =\left(\prod_e s_e^2\right)C(t_1,\ldots,t_M),
\]
where $C(t)=P(1,t_1,\ldots,1,t_M)$ has coefficients in $\mathbb F_2$. Each polynomial $z_{e,0}\partial P/\partial z_{e,0}$ is also homogeneous of degree two in every pair. Hence its normalized value is $(\prod_j s_j^2)C_e(t)$ for a polynomial $C_e$ over $\mathbb F_2$. Since $P\ne0$, $C(t)\ne0$. Multiplying the $z_{e,0}$ zero-gradient equation by $z_{e,0}$ and using $aP=1$ shows
\[
            s_e^r=
            \frac{z_{e,0}\,\partial P/\partial z_{e,0}}{P}
                  =\frac{C_e(t)}{C(t)}\in K^*.
\]

Every such $s_e$ lies in $K'=\mathbb F_{q_0^r}$, with no Kummer extension oracle required. Indeed $s_e^{r(q_0-1)}=1$. The integer $r(q_0-1)$ divides $q_0^r-1$: after dividing by $q_0-1$, the remaining sum $1+q_0+\cdots+q_0^{r-1}$ is divisible by $r$ because $q_0\equiv1\pmod r$. Thus $s_e^{q_0^r-1}=1$. The roots of $X^{q_0^r}-X$ form the unique field with $q_0^r$ elements in the algebraic closure, so $s_e\in K'$. Also $K\subset K'$, because $d$ divides $dr$. All coordinates of the normalized singular point therefore lie in $K'$, and
\[
                         a=1/P\in K'=\mathbb F_{2^{dr}}.
\]
This proves the claimed confinement of all bad parameters. It is stronger than a bound on the number or total degree of bad parameters: their individual field of definition is contained in one explicitly bounded finite field. The proof uses the multihomogeneous rail structure; it makes no such claim for arbitrary degree-$r$ forms.

We next construct a parameter outside this bad field with polynomial bit cost. Compute $d$ by repeated doubling modulo $r$ until the residue one returns; this takes at most $r$ steps and polynomial time in $r$. Compute $D=dr$ and choose the least power of two $b>D$. Then $b\leq2D\leq2r^2$. There is no need to construct the possibly non-power-of-two-degree field $\mathbb F_{2^D}$ itself.

Use the explicit quadratic tower proved in NN-C024: beginning with $\mathbb F_2$, at each degree $j=1,2,4,\ldots,b/2$, find a basis element of trace one in the current field and adjoin a root of $Y^2+Y+u$. Such an element exists because the absolute trace is a nonzero linear functional, and the polynomial is irreducible because $\operatorname{Tr}(y^2+y)=0$. Basis evaluation, squaring and arithmetic have polynomial cost in $b$. Let $\theta$ be the final adjoined root. It lies outside the preceding field of degree $b/2$. Its absolute degree divides $b$; every proper divisor of the power of two $b$ divides $b/2$, and every element of such degree belongs to that preceding field. Hence $\theta$ has absolute degree exactly $b$. In particular $\theta\notin\mathbb F_{2^D}$, since an element of the latter has degree at most $D<b$. The polynomial $H+\theta F$ is therefore smooth by the singular-parameter lemma.

The tower and its basis are an explicit field presentation of polynomial length. If one monic irreducible polynomial is required, compute the first linear dependence of $1,\theta,\ldots,\theta^b$ by linear algebra in the tower basis. Its degree is $b$ by the preceding argument, so it is the minimal polynomial; the same computations give the change to a power basis. This remains polynomial in $b=O(r^2)$. The parameter is nonzero because it has degree $b>1$. No random choice, smoothness test, discriminant expansion, or enumeration of the bad field is used.

We verify that the encoding always meets the geometric lemma. In NN-C023 each incidence rail pair occurs once in its clause factor and once in its variable-consistency factor. Every monomial in each of those factors uses exactly one member of the pair. Thus their product $P$ has degree exactly two in that pair, and it is independent of the extra coordinate $z_0$. The factors are nonzero over $\mathbb F_2$, so $P\ne0$. There are $2M+1=r$ total coordinates, and $F=z_0P$ is homogeneous of total degree $r$. The general hard instances therefore satisfy every hypothesis above, even when their target coefficient $c$ is zero.

For the elementary exceptional source inputs that NN-C023 handled by a non-rail cubic, first compute their parity directly. An empty clause has parity zero; a formula with no declared variables has its single assignment explicitly testable. Replace such an instance by one of two fixed one-variable formulas: the positive unit clause $x$, with parity one, or the two contradictory unit clauses $x$ and $\neg x$, with parity zero. Both have nonempty incidence sets and the same rail construction just proved. This replacement is made only for cases whose answer is already directly known. For ordinary inputs, unused variables can be made incident by the tautological factors in NN-C023, which preserve the exact count. Thus no oracle for the hard parity is used in any edge case.

The equation $H+\theta F$ has polynomial circuit description: the rail product is already succinct, and each pure $r$th power is computed by repeated squaring and multiplication. It is nonzero because $F$ is divisible by $z_0$ whereas the rail terms of $H$ are not. Its hypersurface in $\mathbb P^{r-1}_{\mathbb F_{2^b}}$ is projective, and the zero-gradient criterion proves geometric smoothness. The nontrivial coherent group, in degree $r-2$, is the known line of NN-C023. Its absolute Frobenius is semilinear, with matrix coefficient
\[
          [z_0\prod_e z_{e,0}z_{e,1}](H+\theta F)=\theta c.
\]
The pure powers in $H$ make no contribution to this squarefree monomial. Iterating semilinear Frobenius $b$ times gives the field-linear scalar
\[
                   (\theta c)^{1+2+\cdots+2^{b-1}}
                      =(\theta c)^{2^b-1}=c.
\]
This uses $\theta\ne0$ and $c\in\{0,1\}$. Thus the deterministic smooth specialization preserves the hard entry exactly, rather than only its nonvanishing with some probability.

SGA7 II XXII3.1--3.2 give the same full congruence as in NN-C024,
\[
 \#X(\mathbb F_{2^b})\equiv1+c\pmod2,
 \qquad Z(X/\mathbb F_{2^b},T)\equiv\frac{1-cT}{1-T}\pmod2.
\]
Deligne's Theorem1.6 applies to every produced hypersurface, so these are genuine smooth projective finite-field realizations with pure l-adic cohomology. That theorem is an imported classical input, not a purity conclusion inferred from the congruence.

The construction is now a deterministic polynomial-time reduction from parity-3-SAT. On the encoded family scalar membership in $\oplus\mathrm P$ follows from the factor-choice verifier for $c$ in NN-C023. Point-count parity also lies in $\oplus\mathrm P$ directly: guess a projective coordinate vector in the explicit field, require its first nonzero coordinate to equal one, and verify the equation. This normalization gives exactly one representative per rational point, and $rb=O(r^3)$ guessed bits suffice. The complement by one in the displayed formula is an elementary parity-output conversion. Hence the asserted parity-completeness holds under deterministic reductions. A polynomial-time scalar or parity algorithm gives $\mathrm P=\oplus\mathrm P$; Valiant--Vazirani then gives $\mathrm{NP}=\mathrm{RP}$, as recorded in NN-C006. This strengthens the randomized-promise consequence in NN-C024.

E computes $c$ from the original source formula, keeping its original variable parameter and its stated density dependence. The field construction and the displayed geometric description add only polynomial overhead in the formula size. Neither the $rb$ projective-point witness length nor the full l-adic Betti number is substituted into that sharper original-formula evaluator. A and M are unused in this construction.
\end{proof}
\paragraph{Status and complexity impact.}
Complete singular-parameter confinement, deterministic field construction and preserved-Frobenius proof. Independent complete deduction-proof review by root: passed; imported input theorems are not claimed independently proved. The open finite-field specialization gate of NN-C024 is closed for its actual hard rail family. Smoothness, properness, finite base field, a specified rank-one coherent basis and a deterministic construction now coexist with a parity-complete Frobenius entry. This original construction has field degree at most quadratic in the succinct hypersurface parameter. NN-C031 subsequently removes the growing-field limitation by incidence padding over fixed F4; the present bound is retained as the historical general-parameter result. No polynomial bound on total l-adic rank is asserted. The result strengthens the conditional algorithmic consequence to $\mathrm P=\oplus\mathrm P$ and $\mathrm{NP}=\mathrm{RP}$, rather than merely a randomized reduction. It proves no unconditional complexity lower bound, no VP/VNP separation, and no realization of a GCT canonical basis.

\paragraph{Geometric and representation dictionary.}
Let $P\in\mathbb F_2[z_{e,0},z_{e,1}]$ be nonzero of degree two in each of $M$ pairs, and put $r=2M+1$, $F=z_0P$. In the Fermat pencil $H+uF$, every singular geometric fiber occurs at a parameter in the one finite field $\mathbb F_{2^{r\operatorname{ord}_r(2)}}$. The proof is the pairwise Euler identity, normalization at $z_0=1$, and the resulting root-of-unity and power equations above. A parameter of absolute degree greater than that bounded field degree therefore gives a smooth projective fiber; a quadratic Artin--Schreier tower supplies such a parameter with degree $O(r^2)$ and a polynomial-length presentation.

For the incidence-rail polynomials, the proper smooth fiber has $H^0(\mathcal O)=\mathbb F_{2^b}$ and one other nonzero coherent group, a canonically based line in degree $r-2$. Absolute Frobenius has semilinear coefficient $\theta c$, and the linearized field Frobenius has scalar $c$, the original parity count. The complete Cech proof is NN-C023 and the classical smooth formula is Katz2.3.7.17, p36. SGA1 II4.10(iii)/4.12 is the smoothness interface, SGA7 II XXII3.1--3.2 gives the proper zeta congruence, and Deligne Weil I1.6 gives purity of the actual l-adic realizations. These provide an exact deterministic finite-field geometric dictionary for the hard coefficient. A purity or existence theorem does not evaluate it, and the construction supplies no canonical identification with GCT multiplicities.

\subsection{NN-C031 --- Forced-variable padding puts the smooth rank-one parity gate over fixed F4}
\begin{proposition}[NN-C031]
Fix once and for all
\[
 K=\mathbb F_4=\mathbb F_2[\omega]/(\omega^2+\omega+1).
\]
There is a deterministic polynomial-time reduction from parity-3-SAT to a family of smooth projective hypersurfaces over this fixed field. Each has a polynomial-size homogeneous arithmetic-circuit equation, and its top coherent cohomology is a canonically based one-dimensional space. The $4$-power Frobenius scalar on that line is the original satisfying-assignment parity $c$. Point-count parity is $1+c$.

The geometric statement is uniform: let $k\geq3$ be odd, $r=2^k-1$, $M=(r-1)/2$, and let $P$ be any nonzero polynomial over $\mathbb F_2$ homogeneous of degree two in every one of $M$ variable pairs. Then
\[
 X_P=V\left(z_0^r+\sum_{e=1}^{M}(z_{e,0}^r+z_{e,1}^r)
                         +\omega z_0P\right)
                       \subset\mathbb P^{r-1}_K
\]
is smooth projective. If $c$ is the squarefree coefficient of $z_0P$, its top coherent line has absolute-Frobenius coefficient $\omega c$ and $4$-Frobenius scalar $c$. The hard subclass has $P$ in the succinct rail-factor format, so its scalar and point-count parity are $\oplus\mathrm P$-complete under deterministic polynomial-time reductions. Polynomial-time extraction gives $\mathrm P=\oplus\mathrm P$ and $\mathrm{NP}=\mathrm{RP}$. No varying base field is needed.
\end{proposition}
\paragraph{Exact inputs.}
NN-C030. complete singular-parameter confinement lemma for pairwise degree-two P: every bad a lies in F\_(2\textasciicircum{}(r ord\_r2)).

NN-C023. incidence-rail coefficient encoding; tautological factors for unused declared variables; one-dimensional coherent basis.

SGA1, ExposeII, Section4. Theorem4.10(iii) and Corollary4.12; smooth hypersurface interface used in NN-C030.

Katz, Algebraic Solutions of Differential Equations, Invent.Math.18(1972),1--118. Rank-one Special Case2.3.7.17, printedp36; full direct Cech proof in NN-C023.

SGA 7 II, Expose XXII. Theorem 3.1 (3.1.1); Corollary 3.2 (3.2.1).

Deligne, La conjecture de Weil I, PMIHES43(1974),273--307. Theorem1.6, printedp276.

Valiant--Vazirani, NP is as easy as detecting unique solutions, TCS 47 (1986), 85--93. Theorem 1.1; Corollaries 3.1--3.4, especially Corollary 3.2 p.91.

E: The Exponential Time Hypothesis is False, supplied LaTeX. Theorem1.1; evaluator retains the source formula parameter before padding.

\begin{proof}
First prove the uniform smoothness statement. For $r=2^k-1$, the multiplicative order of $2$ modulo $r$ is exactly $k$. The exponent $k$ works, while for $0<i<k$ the positive integer $2^i-1$ is strictly smaller than $r$ and therefore cannot be divisible by $r$. NN-C030 confines every bad parameter in the Fermat pencil $H+u z_0P$ to
\[
                            \mathbb F_{2^{rk}}.
\]
Both $r$ and $k$ are odd. The element $\omega$ has degree two over $\mathbb F_2$, so it cannot lie in that odd-degree extension: otherwise the tower law would make two divide $rk$. Therefore the fiber at the fixed parameter $\omega$ is smooth. The same reasoning applies to its conjugate $\omega^2$. Projectivity follows from the homogeneous equation. The form is nonzero, since its rail pure-power terms are not divisible by $z_0$, whereas $z_0P$ is.

The coherent calculation from NN-C023 applies over $K$ without change. There is $H^0(\mathcal O)=K$, one other nonzero coherent group in degree $r-2$, and its canonical boundary/Cech basis has absolute-Frobenius coefficient equal to the full squarefree coefficient of the defining form. The Fermat terms do not contribute to that coefficient. Thus the coefficient is $\omega c$, with $c\in\mathbb F_2$. Iterating this semilinear map twice yields the $K$-linear $4$-power Frobenius with scalar
\[
                        (\omega c)^{1+2}=\omega^3c^3=c.
\]
There is no semilinear-versus-linear ambiguity: absolute Frobenius sends $ae$ to $a^2\omega c\,e$, while its second iterate sends $ae$ to $a c\,e$ because $a^4=a$ in $K$.

We next arrange this special degree with polynomial padding while preserving the original count. Start with the incidence-rail construction of NN-C023, with $M_0\geq1$ incidences and $r_0=2M_0+1$. Choose the least odd integer $k\geq3$ such that
\[
                  r=2^k-1\geq r_0,
                 \qquad M=(r-1)/2.
\]
This is computed from the binary length of $r_0+1$ and a parity check. The new degree is linear in the old one: if $k>3$, minimality gives $2^{k-2}-1<r_0$, hence $r<4(r_0+1)-1$; the bounded case $k=3$ also satisfies a constant-factor bound. Consequently $M-M_0=O(M_0)$.

Append exactly $M-M_0$ fresh Boolean variables, each with its own positive unit clause. Each new variable has exactly one allowed value, so every old satisfying assignment has exactly one extension. The number of satisfying assignments, and in particular its parity, is unchanged. Each such gadget adds one clause, one variable and exactly one incidence; no old occurrence is modified. Thus the new incidence count is exactly $M$. In the rail polynomial its clause factor and its two-term variable-consistency factor together have degree two in the new rail pair and squarefree coefficient one. All old factors remain unchanged. This verifies both the required multihomogeneity and the coefficient preservation directly at the polynomial level.

As in NN-C030, treat the elementary exceptional source instances before forming the rails. An empty clause has known parity zero; an instance with no declared variables has its one assignment directly testable. Replace it by the positive unit formula $x$ for parity one or the contradictory unit formula $x\wedge\neg x$ for parity zero. These are fixed instances with nonempty incidence sets. For all other inputs, tautological factors make every otherwise unused declared variable incident without restricting it. Tautological clauses and repeated literals can be simplified first. After these steps each local clause factor is nonzero and the product $P$ is nonzero over $\mathbb F_2$. No decision of the hard parity is used to choose a padding gadget.

The construction now has precisely $r$ coordinates and a degree-$r$ equation over the fixed two-bit field $K$. Every local clause factor has at most eight monomials; each variable-consistency factor has two. The number of factor incidences and their arithmetic-circuit descriptions is polynomial in the source formula size. Each pure power $z^r$ takes $O(\log r)$ multiplications by repeated squaring, so the entire sum of Fermat terms and the rail product has a polynomial-size arithmetic circuit. Its degree is itself polynomial in the source size. This assertion concerns the circuit/factor presentation; fully expanding its coefficients can still require exponentially many monomials.

The actual smooth hypersurface therefore satisfies
\[
 \#X_P(\mathbb F_4)\equiv1+c\pmod2,
 \qquad
 Z(X_P/\mathbb F_4,T)\equiv\frac{1-cT}{1-T}\pmod2,
\]
by SGA7 II XXII3.1--3.2, exactly as in NN-C024. Deligne's Theorem1.6 supplies purity of its l-adic cohomology because every output is smooth and projective. The proof uses the singular-parameter confinement to establish smoothness first; it does not infer smoothness or purity from a point-count congruence.

For a precise complete problem, use the syntactically specified incidence-rail factor format of NN-C023, with nonempty scopes, the displayed consistency factors, and an incidence count giving $r=2^k-1$ for odd $k\geq3$. These conditions are checkable in polynomial time, and fixed straight-line instructions generate the equation $H+\omega z_0P$ from that format. No identity test against an arbitrary purported equivalent circuit is part of the input validation. Every such input is smooth by the uniform theorem. Nondeterministically choose one monomial from each factor and check that the exponents sum to the full squarefree vector. The parity of the number of accepting choices is $c$, proving $\oplus\mathrm P$ membership for the Frobenius scalar on this class. The preceding padded reduction proves hardness. This is a complete problem on a specified geometric presentation family, not a general promise that every smooth hypersurface has a one-dimensional coherent top group.

Point-count parity belongs to $\oplus\mathrm P$ independently of the special coefficient format. Guess its $r$ coordinates in $K$, require the first nonzero coordinate to be one, and check the circuit equation. The chosen normalization gives exactly one vector for each projective rational point. There are $2r$ guessed bits, and all field operations have constant bit cost. The displayed parity formula converts this output to $c$ by adding one. Thus point-count parity on this smooth fixed-field family is also parity-counting complete. An algorithm correct on arbitrary smooth projective inputs in particular handles these guaranteed smooth instances, so polynomial time would give $\mathrm P=\oplus\mathrm P$. Valiant--Vazirani then gives $\mathrm{NP}=\mathrm{RP}$ with the scope already detailed in NN-C006.

E computes the parity by counting the source formula before the forced-variable padding. Its original variable count and clause-density hypotheses therefore remain the relevant parameters, with only polynomial overhead for building the padded geometric input. Running E instead on all geometric coordinates would lose that sharper statement. The field is the one fixed $\mathbb F_4$ throughout, and neither A nor M is used.
\end{proof}
\paragraph{Status and complexity impact.}
Complete fixed-field smoothness, exact padding, coherent-Frobenius and parity-completeness proofs. Independent complete deduction-proof review by root: passed; imported input theorems are not claimed independently proved. The remaining varying-field qualification of NN-C030 is removed: the hard rank-one coherent Frobenius entry is realized deterministically on smooth projective hypersurfaces over fixed $\mathbb F_4$. The padding is linear and parsimonious at the Boolean count level. Neither total l-adic rank nor the fully expanded defining polynomial is claimed polynomially bounded; the geometric equation is succinct. This yields a strictly smaller effective gate than arbitrary field construction or random specialization. It supplies actual pure finite Weil realizations for the counting gadgets without evaluating the hard coefficient or identifying a GCT canonical basis.

\paragraph{Geometric and representation dictionary.}
For every odd $k\geq3$, every nonzero pairwise-degree-two $P$ over $\mathbb F_2$ in $M=2^{k-1}-1$ pairs defines the smooth projective hypersurface
\[
 V\left(z_0^{2^k-1}+\sum_e(z_{e,0}^{2^k-1}+z_{e,1}^{2^k-1})
                        +\omega z_0P\right)/\mathbb F_4.
\]
The singular-parameter field from NN-C030 has odd degree $k(2^k-1)$ and therefore excludes the fixed quadratic parameter $\omega$. Its top coherent line is explicitly trivialized by the boundary/Cech construction. Absolute Frobenius has semilinear coefficient $\omega c$; $4$-power Frobenius has scalar $c$. The lower coherent groups are only the degree-zero field, and the proper zeta congruence is $(1-cT)/(1-T)$ modulo two. Exact classical interfaces are Katz2.3.7.17, SGA1 II4.10(iii)/4.12, SGA7 II XXII3.1--3.2 and Deligne Weil I1.6.

The forced-unit-clause padding turns any source formula into this geometric class while preserving its count exactly. Thus the field, smoothness and coherent basis are all supplied effectively, and the sole unknown scalar is parity-counting complete. This is an actual smooth projective finite-field construction, with pure l-adic cohomology, and a precise computational operation inside it. It does not assert a canonical representation-theoretic realization of permanent/determinant multiplicities.

\subsection{NN-C032 --- An explicit fixed-F2 smooth projective model has four-dimensional coherent cohomology with hard Frobenius}
\begin{proposition}[NN-C032]
From the fixed-$\mathbb F_4$ hypersurface $X=V(G)$ of NN-C031, with $G=H+\omega F$ in $r$ coordinates, one can deterministically construct a smooth geometrically integral projective scheme
\[
             W\subset\mathbb P^{r^2-1}_{\mathbb F_2},
              \qquad W\simeq\operatorname{Res}_{\mathbb F_4/\mathbb F_2}X.
\]
It has a polynomial-size homogeneous circuit-equation presentation consisting of $O(r^4)$ quadratic equations and $O(r)$ degree-$r$ equations. Put $d=r-2$, which is odd. Its coherent cohomology is
\[
 H^0(W,\mathcal O)=\mathbb F_2,\quad
 H^d(W,\mathcal O)=\mathbb F_2^2,\quad
 H^{2d}(W,\mathcal O)=\mathbb F_2,
\]
and every other group vanishes. In explicitly descended bases, absolute Frobenius is respectively
\[
                   1,\qquad
              c\begin{pmatrix}0&1\\1&0\end{pmatrix},
                   \qquad c,
\]
where $c$ is the original satisfying-assignment parity. Thus the entire coherent realization has dimension four, its basis is supplied, and one of its one-dimensional Frobenius scalars is parity-counting complete over the fixed prime field.

Moreover
\[
 \#W(\mathbb F_2)=\#X(\mathbb F_4),\qquad
 \#W(\mathbb F_2)\equiv1+c\pmod2,\qquad
 Z(W/\mathbb F_2,T)\equiv\frac{1-cT}{1-T}\pmod2.
\]
These yield deterministic parity-counting-complete problems on the encoded family. Every output is smooth projective and has pure l-adic cohomology. The presentation uses a projective system of equations; it is not asserted to be a hypersurface over $\mathbb F_2$.
\end{proposition}
\paragraph{Exact inputs.}
NN-C031. deterministic fixed-F4 smooth projective hypersurface, canonical coherent basis and absolute Frobenius coefficient omega*c.

Stacks Project, Constructions of Schemes. Lemma27.13.6, Tag01WD, Segre embedding; the exact two-by-two-minor chart computation is proved below.

Stacks Project, Criteria for Representability. Section97.11, Tag05Y8; equation97.11.0.2, Lemma97.11.1 (fppf sheaf) and Lemma97.11.2 (base change).

Stacks Project, Descent. Lemma35.23.29, Tag02VL: smoothness is fpqc local on the base.

Stacks Project, Cohomology of Schemes. Lemma30.5.2, Tag02KH, flat base change; especially its affine-base statement(2) and Cech proof.

Stacks Project, Varieties. Lemma33.29.1, Section Tag0BEC: coherent Kunneth for quasi-compact schemes with affine diagonal; cup-product normalization and complete proof.

Stacks Project, Varieties and Commutative Algebra. Lemma33.25.3, Tag056S: smooth over a field implies regular; Lemma10.106.2, Tag00NP: a regular local ring is a domain.

SGA 7 II, Expose XXII. Theorem 3.1 (3.1.1); Corollary 3.2 (3.2.1).

Deligne, La conjecture de Weil I, PMIHES43(1974),273--307. Theorem1.6, printedp276.

Valiant--Vazirani, NP is as easy as detecting unique solutions, TCS 47 (1986), 85--93. Theorem 1.1; Corollaries 3.1--3.4, especially Corollary 3.2 p.91.

E: The Exponential Time Hypothesis is False, supplied LaTeX. Theorem1.1, original-source evaluator.

\begin{proof}
Write $K=\mathbb F_4$ and let $\sigma$ be its nontrivial automorphism, so $\sigma(\omega)=\omega^2=\omega+1$. Use $r^2$ homogeneous coordinates over $\mathbb F_2$: diagonal coordinates $a_i$, and coordinates $b_{ij},c_{ij}$ for every $i<j$. Form the matrix with entries in $K[a,b,c]$
\[
 M_{ii}=a_i,\qquad
 M_{ij}=b_{ij}+\omega c_{ij},\qquad
 M_{ji}=b_{ij}+\omega^2 c_{ij}\quad(i<j).
\]
Coefficientwise conjugation, with the coordinate indeterminates held fixed, satisfies $\sigma(M)=M^{\mathsf t}$. For a polynomial $A$ over $K$ in these coordinates, write uniquely $A=A_0+\omega A_1$ with $A_0,A_1$ over $\mathbb F_2$.

Define $W$ by the following homogeneous equations over $\mathbb F_2$: both coefficient components of every two-by-two minor of $M$, and both components of
\[
                     G(M_{1j},\ldots,M_{rj})
                       \quad(1\leq j\leq r).
\]
This is a closed projective scheme. There are $2\binom r2^2=O(r^4)$ minor equations and $2r$ column equations, with degrees two and $r$ respectively; redundant or zero equations may be retained.

We identify this scheme after the faithfully flat extension $K/\mathbb F_2$, including its scheme structure. On each off-diagonal pair, the coordinate change
\[
        (b_{ij},c_{ij})\longmapsto(M_{ij},M_{ji})
\]
has determinant $\omega^2-\omega=1$, so over $K$ the entries of $M$ are independent matrix coordinates. Imposing both $\mathbb F_2$ coefficient components of a polynomial is equivalent over $K$ to imposing that polynomial and its coefficientwise conjugate. For the minors the conjugates are other minors of the transpose, already in the list. Thus their equations define the ordinary projective rank-one matrix scheme.

Here is the exact scheme-theoretic chart check. On a projective chart where $M_{ij}$ is nonzero, normalize it to one. The minors force
\[
                         M_{kl}=M_{kj}M_{il}.
\]
The remaining entries of column $j$ and row $i$ are free coordinates, giving the standard affine chart of
$\mathbb P^{r-1}\times\mathbb P^{r-1}$ under the Segre map $(z,w)\mapsto zw^{\mathsf t}$. These charts cover the projective minor scheme, so no extra embedded component is being discarded. This is the explicit minor form of the Segre embedding in Stacks27.13.6.

Under $M=zw^{\mathsf t}$, the column equations are
\[
                      w_j^rG(z)=0\quad(1\leq j\leq r).
\]
Their conjugates are the row equations
\[
                      z_i^rG^\sigma(w)=0\quad(1\leq i\leq r),
             \qquad G^\sigma=H+\omega^2F.
\]
Locally on the product of projective spaces some $z_i$ and some $w_j$ are units, so the generated sheaf of ideals is exactly $(G(z),G^\sigma(w))$. This proves an isomorphism of schemes
\[
                         W_K\simeq X\times_K X^\sigma.
\]
In these matrix coordinates the descent automorphism conjugates coefficients and transposes $M$; it therefore interchanges the two factors with their conjugation data. This is precisely the descent datum of restriction of scalars for the quadratic extension. More formally, $K\otimes_{\mathbb F_2}K\simeq K\times K$, so the restriction-of-scalars functor base-changes to $X\times X^\sigma$, with that same factor exchange. The functor is an fppf sheaf and commutes with base change by Stacks97.11.1--97.11.2. The displayed compatible isomorphism over the faithfully flat cover thus descends to an identification of its sheaf with the represented sheaf of $W$. This proves the claimed Weil-restriction identification; no general projectivity theorem for restriction of scalars is needed because the projective equations are already explicit.

Both $X$ and $X^\sigma$ are smooth by NN-C031. Their product is smooth, and smoothness descends under the finite faithfully flat extension by Stacks35.23.29. Hence $W$ is smooth over $\mathbb F_2$, of dimension $2(r-2)=2d$. Over an algebraic closure, each factor has only the ground field as its degree-zero coherent cohomology, by NN-C023's Cech calculation. The same is true of the product by coherent Kunneth. It is therefore connected, since a decomposition into two open-and-closed pieces would give a nontrivial idempotent in its global functions. A smooth scheme over a field has regular local rings by Stacks33.25.3, and regular local rings are domains by Stacks10.106.2; distinct irreducible components consequently cannot intersect. Thus this connected smooth geometric scheme is integral. This proves geometric integrality of $W$ as well.

We now compute its full coherent Frobenius, taking care with descent and semilinearity. Over $K$, let $e$ be the canonical basis of $H^d(X,\mathcal O_X)$ from NN-C023 and let $e'$ be its conjugate basis on $X^\sigma$. Coherent Kunneth (Stacks33.29.1) and flat base change (Stacks30.5.2) show that the only cohomology groups of $W_K$ are
\[
 K\quad\text{in degree }0,\qquad
 Ke\oplus Ke'\quad\text{in degree }d,\qquad
 K(e\mathbin\cup e')\quad\text{in degree }2d.
\]
Here $e,e'$ in the middle line mean their pullbacks from the two factors. The Kunneth isomorphism is the pullback-and-cup-product map, so it is compatible with Frobenius and the factor exchange, including its graded sign.

The descent automorphism is semilinear on coefficients and sends $e$ to $e'$ and $e'$ to $e$. Thus the following middle vectors are fixed:
\[
                    v_1=e+e',\qquad
                    v_2=\omega e+\omega^2e'.
\]
They are independent over $K$, because the determinant of their coefficients is $\omega^2-\omega=1$. Their fixed coefficient combinations are exactly their $\mathbb F_2$ span. The top vector $w=e\cup e'$ is also fixed: exchange gives $e'\cup e=(-1)^{d^2}e\cup e'$, and this sign is one in characteristic two. Hence $v_1,v_2,w$ are explicit descended bases over $\mathbb F_2$. This also follows by taking the invariant vectors in the flat-base-change isomorphism; it does not require computing any unknown Frobenius entry.

Let $\mathfrak F$ denote absolute Frobenius on coherent cohomology. On the two factors, NN-C031 gives
\[
                \mathfrak F(e)=\omega c\,e,
                \qquad\mathfrak F(e')=\omega^2c\,e',
 \qquad \mathfrak F(a\xi)=a^2\mathfrak F(\xi).
\]
Under flat base change from $\mathbb F_2$ to $K$, this absolute Frobenius corresponds to Frobenius on $H^*(W,\mathcal O_W)$ tensored with squaring on $K$, not with the identity on $K$. Absolute Frobenius is natural for the displayed isomorphism with the product. Restricting its semilinear action to the descended $\mathbb F_2$ vectors therefore computes the original $\mathbb F_2$-linear Frobenius on $W$ exactly.
Applying these formulas to the descended middle basis gives
\[
       \mathfrak F(v_1)=c v_2,\qquad
       \mathfrak F(v_2)=c v_1.
\]
For example, the coefficient of $e$ in the second identity is
$\omega^2\omega c=c$, and the other coefficient is
$(\omega^2)^2\omega^2c=c$. Cup-product compatibility gives
\[
                 \mathfrak F(w)=\omega\omega^2c^2w=cw.
\]
On degree zero it is the identity over $\mathbb F_2$. These computations prove exactly the three displayed matrices. As a separate consistency check, the square of the middle matrix is $cI$, agreeing with $4$-power Frobenius on both factors; the top scalar has square $c$ as required. The argument supplies a basis and a symbolic matrix description in terms of $c$, not a polynomial-time evaluation of that unknown bit.

There is also an elementary exact rational-point bijection, which does not use a trace formula. If $[z]\in X(K)$, define
\[
                            M=z\sigma(z)^{\mathsf t}.
\]
It is Hermitian and rank one. Replacing $z$ by $\lambda z$ for $\lambda\in K^*$ leaves it unchanged, since $\lambda\sigma(\lambda)=\lambda^3=1$. Every column is a scalar multiple of $z$, so the column equations vanish. Its diagonal entries belong to $\mathbb F_2$, and every off-diagonal pair has the unique displayed $(b_{ij},c_{ij})$ coordinates in $\mathbb F_2$. Thus it defines an $\mathbb F_2$ point of $W$.

Conversely, an $\mathbb F_2$ point gives a nonzero Hermitian rank-one matrix over $K$. It has a nonzero diagonal entry. If every diagonal were zero, the minor relation
$M_{ij}M_{ji}=M_{ii}M_{jj}$ would force every off-diagonal entry to vanish, since $M_{ij}M_{ji}$ is its nonzero norm whenever $M_{ij}\ne0$. This would give the zero matrix. Choose $i$ with $M_{ii}=1$, and let $z$ be column $i$. The minors give $M=z\sigma(z)^{\mathsf t}$, with $z_i=1$. The column equation implies $G(z)=0$. The column space recovers $[z]$ uniquely. These two constructions are inverse, proving
\[
                             \#W(\mathbb F_2)=\#X(K).
\]
Projective scaling on the $\mathbb F_2$ matrix coordinates introduces no additional quotient, since $\mathbb F_2^*=\{1\}$.

For the full zeta congruence, apply SGA7 II XXII3.1 to the just-computed coherent Frobenius. The degree $d$ is odd and $2d$ is even. The middle determinant is
\[
       \det\left(I-Tc\begin{pmatrix}0&1\\1&0\end{pmatrix}\right)
                      =1-c^2T^2=(1-cT)^2
\]
in $\mathbb F_2[T]$. Dividing by the degree-zero and top determinants gives
\[
     Z(W/\mathbb F_2,T)\equiv
         \frac{(1-cT)^2}{(1-T)(1-cT)}
                   =\frac{1-cT}{1-T}\pmod2.
\]
The point-count congruence is also immediate from the exact bijection and NN-C031. The full series congruence is obtained from the coherent determinants, not inferred merely from parity of the first point count. Deligne's theorem applies to the smooth projective $W$, giving the actual pure l-adic realizations independently of this characteristic-two calculation.

All constructions are effective with polynomial output length. A $K$ arithmetic gate can be replaced by constantly many $\mathbb F_2$ gates computing its two coordinate components: for example
\[
 (a+b\omega)(u+v\omega)
                  =(au+bv)+(av+bu+bv)\omega.
\]
The $O(r^4)$ minors have constant-size circuits, and substituting each of the $r$ columns into $G$ multiplies its circuit size by at most $r$, up to a constant coordinate-arithmetic factor. There are only $r^2$ projective coordinates. No expansion of $G$ into all monomials is required, and no factorization or Galois-descent oracle is hidden in these formulas.

The source-to-$W$ construction is a deterministic polynomial-time reduction. Extracting the top Frobenius scalar is exactly extracting $c$, so it is $\oplus\mathrm P$-hard and is in $\oplus\mathrm P$ on the syntactic input family through the factor-choice verifier of NN-C023. Point-count parity is directly in $\oplus\mathrm P$: guess the $r^2$ binary projective coordinates, reject the zero vector, and check the homogeneous circuit equations. Every nonzero binary vector is already a unique projective representative. Its output is $1+c$, which has the same parity-counting complexity; complementation preserves $\oplus\mathrm P$ by adding one distinguished accepting computation. Thus the point-parity decision language is also complete under deterministic many-one reductions. A polynomial-time algorithm for either geometric task gives $\mathrm P=\oplus\mathrm P$ and $\mathrm{NP}=\mathrm{RP}$ as before.

E evaluates $c$ through the original formula before its geometric padding and descent. The extra construction cost is polynomial in that source size. A and M are unused. The geometric input is a prescribed homogeneous circuit-equation family with its mechanically generated format; no polynomial-time recognition theorem for arbitrary smooth projective presentations is being asserted.
\end{proof}
\paragraph{Status and complexity impact.}
Complete explicit projective descent, scheme-structure, coherent-Frobenius, point-bijection and complexity proofs. Independent complete deduction-proof review by root: passed; imported input theorems are not claimed independently proved. The hard coherent-Frobenius gate now occurs over the fixed prime field, in a smooth geometrically integral projective scheme with a supplied four-dimensional coherent basis. Neither random specialization, growing coefficient fields, singular fibers, unknown coherent rank nor unknown basis is the missing operation. The single bit $c$ is still parity-counting complete. The total l-adic rank and expanded equation size are not claimed polynomial. This identifies a concrete remaining scalar-extraction gate inside an actual pure finite Weil realization; it does not supply GCT's canonical-basis functor or an unconditional circuit lower bound.

\paragraph{Geometric and representation dictionary.}
Let $X=V(H+\omega F)/\mathbb F_4$ be a hypersurface of NN-C031. The displayed Hermitian-matrix equations define a projective $\mathbb F_2$ scheme $W$ whose base change is the Segre product $X\times X^\sigma$ and whose descent datum is factor exchange. Thus $W$ represents the Weil restriction. It is smooth and geometrically integral, with coherent cohomology only in degrees $0,d,2d$, of ranks $1,2,1$. In the explicit descended bases $1$, $(e+e',\omega e+\omega^2e')$, and $e\cup e'$, its Frobenius operators are $1$, $c\left(\begin{smallmatrix}0&1\\1&0\end{smallmatrix}\right)$, and $c$. Its exact rational-point bijection with $X(\mathbb F_4)$ and its proper zeta congruence are proved above.

The scheme-theoretic interfaces are the explicit Segre chart proof with Stacks27.13.6, restriction-of-scalars sheaf/base-change Lemmas97.11.1--97.11.2, smoothness descent Lemma35.23.29, coherent flat base change Lemma30.5.2 and coherent Kunneth Lemma33.29.1. SGA7 II XXII3.1--3.2 supplies the coherent zeta congruence and Deligne Weil I1.6 supplies purity. These are exact locators, with no original Grothendieck archive cote claimed. The formula for the descended middle matrix is proved by semilinear arithmetic and cup-product compatibility, not guessed from point counts. The underlying count reduction and E evaluator are unchanged; the new construction removes the fixed-field obstruction while preserving an explicit small coherent realization.

\subsection{NN-C033 --- Integral swap matrices remove the fixed-characteristic symbolic-circuit gate}
\begin{proposition}[NN-C033]
Fix arity $w$ and factor density $C$. For every prime $p$ and every factor hypergraph on $n$ Boolean variables with $m\leq Cn$ factors of scope at most $w$, its universal partition polynomial
\[
 Z_H(y)=\sum_{x\in\{0,1\}^n}\prod_f y_{f,x|I_f}
                      \in\mathbb F_p[y]
\]
has a deterministically constructible, division-free arithmetic circuit over $\mathbb F_p$ of size
\[
 S(n)=\operatorname{poly}(n)\,
       2^{O_{w,C}(n(\log\log n)^2/\sqrt{\log n})}.
\]
Construction bit time is $S(n)\operatorname{poly}(\log p)$. The equality is equality of formal polynomials, with arbitrary independent weights, in every characteristic including two. The output circuit uses only prime-field constants. A small auxiliary field is used during construction and eliminated by explicit coordinate simulation.

Consequently all ordinary ideal-multiplier and affine-trivialization constructions of NN-C029 hold over every field of characteristic $p$, with the same variable parameter and subexponential bound. In particular, an unsatisfiable bounded-density local Boolean verifier has a Hilbert-like IPS certificate over $\mathbb F_2$ of that size, with ordinary polynomial identity $\sum_i a_i f_i=1$ and multiplier degrees at most $v$. For a fixed-width $O(n)$-clause CNF the original-$n$ specialization still needs no Boolean axioms. The output multipliers are exactly reductions of the explicitly specified integral first-violation and local-correction polynomials.

After polynomial overhead, the prime-field circuits can be made constant-free with syntactic formal degree at most $\max(m,1)$ for $Z_H$, at most $\max(v,1)$ for the multipliers, and at most $v+1$ for the IPS certificate. Here constant-free means that the only constant leaves are $0,1,-1$, and formal degree gives each constant or variable leaf degree one, addition the maximum, and multiplication the sum. This is a positive-characteristic statement. This prime-field argument itself does not construct integral-gate circuits over $\mathbb Z$; that remaining step is subsequently proved in NN-C034.
\end{proposition}
\paragraph{Exact inputs.}
E: The Exponential Time Hypothesis is False, supplied LaTeX. Lemma3.2 complete integer-measurement-fiber proof; Lemma3.3 ring-valued Boolean differences; Proposition3.4 Taylor routine; Lemmas4.1--4.3 swap matrix and row support; equation swap-matrix-definition; Lemmas4.5--4.8 layout, height and parameters.

NN-C016. bounded-arity universal partition polynomial, high-incidence branching and coefficient-independent circuit unrolling.

NN-C029. ordinary integral ideal-multiplier correction, sparse-CNF first-violation identity, homogeneous components and affine coordinate maps.

EGA IV, fourth part, Publications Mathematiques IHES32 (1967). Printed p.111,18.2.1--18.2.2; base-change equation18.2.1.2; primary PDF page110 (zero-based109).

Elementary finite-field construction and Chinese remainder theorem. Complete existence, deterministic bounded-size construction and Hasse-jet argument proved below; no fast irreducible-polynomial oracle is assumed.

\begin{proof}
We modify the actual common routine of E, rather than reducing its rational circuit modulo $p$. The key is an integral strengthening of its auxiliary swap matrix.

Use E's notation: commuting integral permutation involutions $\Pi_1,\ldots,\Pi_Q$, their product $V_0$, and
\[
 W_0=\frac12\sum_{s=1}^Q(I-\Pi_s),\qquad
 \mathcal I=\{A_0,\ldots,A_0+R-1\},\qquad
 M_0=(-1)^{A_0}V_0\sum_{j=0}^{R-1}
                      (-2)^j\binom{W_0-A_0I}{j}.
\]
Although this displayed construction uses rational matrices, $M_0$ is in fact integral. In the formal matrix power-series ring over $\mathbb Q$, one has
\begin{equation}\label{eq:C033-integral-generating}
 \sum_{j\geq0}(-2)^j\binom{W_0-A_0I}{j}U^j
       =(1-2U)^{-A_0}\prod_{s=1}^Q\bigl(I+U(\Pi_s-I)\bigr).
\end{equation}
To prove the identity, use the common rational eigenbasis in E Lemma4.1. If exactly $a$ involutions act by minus one, the right side acts by
$(1-2U)^{-A_0}(1-2U)^a=(1-2U)^{a-A_0}$, whose generalized binomial expansion is exactly the left side. Equality on this basis proves the matrix identity coefficient by coefficient. The scalar series $(1-2U)^{-A_0}$ has integral coefficients: for $A_0>0$ its coefficient of $U^j$ is $2^j\binom{A_0+j-1}{j}$; for $A_0=0$ the entire series is $1$. Every factor on the right is an integral matrix polynomial. Therefore each coefficient on the left is integral, and so is $M_0$.

This proof also gives a direct integral construction by truncating the product at degree $R-1$. It takes polynomially many operations in the auxiliary dimension $D$, $Q$ and $R$. Since $A_0,R\leq Q$, the coefficients have polynomial bit length, for example the norm bounds of E Lemma4.6 imply $O(Q\log(Q+2))$ bits for entries of $M_0$. The stronger conclusion is integrality, not a new asymptotic height assumption. The generating identity is proved over $\mathbb Q$ only to identify this particular integral matrix; after the integer coefficients have been formed, reduction or scalar extension is an ordinary ring homomorphism. No occurrence of $1/2$ or $1/j!$ is evaluated in characteristic two.

Keep the integer measurement matrix $C^{\mathcal B}$ of E Lemma4.1, the same incidence labels and copies, and the same integral matrices $C_{t,g}$. Two facts of E's construction remain available. First, $M_0[u,y]\ne0$ implies at most $R-1$ mismatches between $a_y(s)$ and $d_u(s)$, by Lemma4.2. Second,
\[
   C^{\mathcal B}\xi=0,\quad \xi\in\mathbb Z^{\mathcal B}
                  \quad\Longrightarrow\quad M_0\xi=\xi
                       \quad\hbox{in }\mathbb Z^{\mathcal B}.
\]
Indeed the hypothesis also holds over $\mathbb Q$, where Lemma4.1 proves the conclusion; integrality then makes it an exact integer equality. We do not claim that the reduction of $M_0$ fixes the entire kernel of the reduced matrix over $\mathbb F_p$. That stronger assertion is unnecessary and is not obtained by reducing a rational kernel calculation.

Choose an explicit extension field $L/\mathbb F_p$ containing at least $Q$ distinct nonzero elements, denoted $\alpha_1,\ldots,\alpha_Q$. If $p>Q$, take $L=\mathbb F_p$ and $\alpha_t=t$. If $p\leq Q$, let $e$ be least with $p^e>Q$. Then
\[
                 p^e\leq pQ\leq Q^2,\qquad e=O(\log(Q+1)).
\]
An explicit monic irreducible polynomial of degree $e$ can be found deterministically by enumerating the $p^e$ candidates and testing division by every monic polynomial of positive degree at most $\lfloor e/2\rfloor$. A reducible polynomial has a proper factor of degree at most half its degree, so this test is exact, and the total enumeration has cost polynomial in $Q$ and $\log p$.

For completeness, existence does not require a randomized construction theorem. In an algebraic closure the roots of $T^{p^e}-T$ are distinct and closed under the field operations, so they form a field with $p^e$ elements. Every element of smaller degree has degree a proper divisor of $e$, by the period of Frobenius, and hence belongs to the union of the root sets of $T^{p^d}-T$ for proper divisors $d$ of $e$. For $e\geq2$, the size of that union is at most
$\sum_{d=1}^{\lfloor e/2\rfloor}p^d<p^e$. Some element therefore has degree $e$, and its minimal polynomial is an irreducible candidate. The case $e=1$ is immediate. Represent $L$ by the resulting quotient with basis $1,\theta,\ldots,\theta^{e-1}$, and take any $Q$ distinct nonzero coordinate vectors. All arithmetic in this field has bit cost polynomial in $Q$ and $\log p$.

For a pair $(u,y)$ with $M_0[u,y]\ne0$, write
$E_{u,y}=\{s:a_y(s)\ne d_u(s)\}$ and replace E's scalar polynomial by
\[
       \varphi_{u,y}(T)=\prod_{s\in E_{u,y}}(1-T/\alpha_s).
\]
Form the new rows over $L$ by
\[
 \vec P_u(T)=\delta_u(I-M_0)
          +\sum_yM_0[u,y]\varphi_{u,y}(T)\delta_y.
\]
Integer matrix entries in this formula are reduced into $L$. The rows have degree at most $R-1$ and satisfy $\vec P_u(0)=\delta_u$. At $T=\alpha_t$, every term with $a_y(t)\ne d_u(t)$ vanishes. All denominators here are nonzero field elements chosen explicitly.

The separation property needed by the routine still holds on its actual enumeration domain. Suppose two Boolean assignments differ only inside an off-target group and have equal \emph{integer} measurements $C_{t,g}x_g$. Their difference, copied into each auxiliary array, is an integer vector in the kernel of $C^{\mathcal B}$ and supported in that off-target group. The term $I-M_0$ annihilates this integer vector by the exact integer equality above. Every surviving $\varphi$ term uses the target group and hence also annihilates the difference. Reducing these equalities into $L$ proves that $\langle\vec P_{f,v}(\alpha_t),x\rangle$ is constant on each such integer measurement fiber. The same reasoning applies over every commutative $L$-algebra. It uses an equality of integer vectors before scalar extension, not an assumed saturation of a reduced kernel.

We spell out why this weaker fiber statement suffices for the entire symbolic routine. E Lemma3.2 enumerates tuples of integer measurement values in $[-D,D]$, scans Boolean assignments to find one representative in each nonempty local fiber, and then scans the same fibers again. To evaluate a factor assigned to group $g$, it combines the actual group assignment with fixed representatives in all other groups. Every intermediate assembled vector is Boolean. Successively replacing an off-target group's representative by another assignment in the same integer fiber leaves that factor unchanged by the property just proved. Distributivity then gives the exact product of group sums and sum over measurement tuples. The proof never needs invariance on arbitrary vectors over $L$, and it never replaces integer measurement equality by equality modulo $p$. Thus its vector count, storage bound and coefficient-independent control remain exactly those in Lemma3.2.

Let $\mathcal A=L[y]$, with the factor weights retained as formal independent indeterminates, and use the standard multilinear factor interpolation
\[
 \Phi_f(z)=\sum_{a\in\{0,1\}^{I_f}}y_{f,a}
            \prod_{i:a_i=1}z_i\prod_{i:a_i=0}(1-z_i).
\]
Its value at a Boolean local input is the corresponding indeterminate as a formal polynomial, in every characteristic. Define the polynomial $H(T)$ using the new replacement rows as in E Section3.1. Then $H(0)=Z_H(y)$ exactly in $\mathcal A$, and $\deg_T H\leq dR$ with $d\leq wm=O(n)$.

Set $q=\lfloor dR/Q\rfloor$ and $J=q+1$. At node $\alpha_t$ write
$\Delta P(z)=P(\alpha_t+z)-P(\alpha_t)$. Its entries are divisible by $z$. Introduce the independent Boolean-substitution indeterminates $X_i$ used in E's auxiliary polynomial $\mathcal G_x(X,z)$. Every occurrence of $X_i$ carries a factor $z$, so $\deg_X[z^j]\mathcal G_x\leq j$. This is a polynomial degree statement over $\mathcal A$ and has no characteristic restriction.

E Lemma3.3 is explicitly an identity over every commutative ring. Its proof replaces positive powers of each $X_i$ by $X_i$ and performs subset Mobius inversion with coefficients $1,-1$. Apply it coefficientwise to $[z^j]\mathcal G_x$, for $j<J$. For each fixed substitution $X=\mathbf1_U$, the corrections involving $\Delta P$ are independent of the enumerated Boolean assignment $x$. Hence the integer-fiber separation just proved computes the corresponding restricted sum with $x_i=1$ for $i\in E$. The resulting identity is the exact Hasse jet
\[
                      H(\alpha_t+z)\bmod z^J
                         \quad\hbox{in }\mathcal A[z]/(z^J).
\]
All equalities are formal polynomial identities in the weight indeterminates $y$. We have not compared numerical evaluations at finitely many $y$ values.

Finally the ideals $(T-\alpha_t)^J$ are pairwise comaximal in $L[T]$. Their differences are nonzero units. The Chinese remainder theorem gives an isomorphism
\[
 L[T]\big/\!\prod_{t=1}^Q(T-\alpha_t)^J
              \longrightarrow\prod_{t=1}^Q L[z]/(z^J),
       \qquad H\longmapsto(H(\alpha_t+z))_t.
\]
Both sides have dimension $QJ$, with the monomial and jet bases. Its invertible matrix has entries
$\binom uj\alpha_t^{u-j}$, using integer binomial coefficients reduced into $L$. This is Hasse interpolation: it divides by no $j!$ and requires no inequality $J<p$. Compute its inverse over the explicit finite field and use only its final entries as scalar constants. Tensoring this matrix identity with the coefficient algebra $\mathcal A$ recovers $H$ and $H(0)$ formally. Any internal field inversions used to find the constants disappear from the resulting arithmetic circuit. Since $QJ>dR$, the recovered polynomial is the original $H$.

The number of enumerated vectors is the same as E Proposition3.4:
\[
 2Qr\left(\sum_{s=0}^{\min(q,n)}2^s\binom ns\right)
                       (2D+1)^{2\tau n}2^{2n/r}.
\]
The combinatorial layouts and their integer measurements are unchanged. For bounded arity and linear total incidence, branch on at most $n/\sqrt{\log_2n}$ variables of incidence above
$\max(1,\lceil M_{\rm inc}/n\rceil)\lceil\sqrt{\log_2n}\rceil$, retaining declared scopes. On each branch E Lemmas4.7--4.8 apply with the same $Q,R,r,D,\tau$ and the same three exponential costs. The field and matrix computations introduce polynomial overhead only. Bounded exceptional input sizes are handled by direct expansion. Summing all branches yields $S(n)$ operations in the coefficient algebra, with construction bit cost $S(n)\operatorname{poly}(\log p)$.

Unroll these operations into a shared division-free circuit over $L$. Every branch, representative and layout choice depends only on discrete scopes, Boolean bits and integer measurement data. No test of whether a symbolic polynomial is zero occurs. Replace each $L$-valued gate by its $e$ coordinates over $\mathbb F_p$. Addition is componentwise and multiplication uses the explicit structure constants for the quotient basis, at polynomial overhead in $e$. A weight input $y$ is represented by $(y,0,\ldots,0)$. The final computed element is precisely the universal partition polynomial, whose coefficients already lie in $\mathbb F_p$. Its coordinate of one is therefore exactly $Z_H(y)$ and every other coordinate is the zero formal polynomial. Taking that coordinate supplies the required prime-field circuit; no identity test is needed to know the other coordinates vanish. This also proves agreement with the reduction of the original integral universal partition polynomial, not merely agreement as a function on a finite field.

NN-C029's preliminary first-violation multipliers are constant-size substitutions into these universal partition polynomials. Its at-most-eight local specializations and integral monomial division establish the ordinary identity over $\mathbb Z$; performing those same constructions over $\mathbb F_p$ reduces each polynomial identity coefficientwise. Consequently the final circuits compute the reductions of the specified integral multipliers and satisfy $\sum_i a_if_i=1$ as a formal identity. The stronger sparse-CNF identity and affine coordinate formulas transfer for the same reason. For example the mutually inverse maps remain
\[
 z_i=y_i+a_i(x)(2t-t^2),\qquad
 t=\sum_i f_i(x)z_i,\qquad
 y_i=z_i+a_i(x)(t^2-2t),
\]
with coefficients reduced in the field. In characteristic two the terms containing $2t$ vanish, and the directly verified identities still hold. Original-$n$ versus $n+s$ verifier parameters are exactly those already proved in NN-C029.

To obtain the constant-free formal-degree claim, apply homogeneous-component truncation through the known output degree $D_0$ in the polynomial indeterminates. It has polynomial overhead and computes the same output because no division by an indeterminate is present. Every degree-zero component is a prime-field constant and can be evaluated explicitly. A multiplication with such a scalar can be implemented by binary doubling and addition of its other input, using $O(\log p)$ addition gates; this does not increase formal degree. Each nonconstant multiplication combines positive-degree components, so induction bounds its formal degree by the sum of their variable degrees. Constant outputs themselves can be built from one by additions, with formal degree one. Replacing all scalar multiplications in this manner, and using constant leaves only $0,1,-1$, gives formal degree at most $\max(D_0,1)$ and polynomial overhead in $\log p$. Take $D_0=m$ for the universal partition polynomial, $D_0=v$ for multipliers, and $D_0=v+1$ for the Hilbert-like IPS polynomial. This argument uses reduction in the prime field; it does not convert the original rational circuits into integral-gate circuits over $\mathbb Z$.
\end{proof}
\paragraph{Status and complexity impact.}
Complete integral-matrix, modified finite-characteristic routine, formal-output descent and ideal-certificate proofs. Independent complete deduction-proof review by root: passed; imported input theorems are not claimed independently proved. This closes the fixed-small-characteristic gap left by NN-C016 and the large-characteristic-only addendum of NN-C029. It gives formal symbolic partition circuits, ordinary ideal certificates and empty-case affine trivializations over fixed F2 at E's subexponential source-parameter bound. The proof does not merely reduce rational interpolation constants modulo two: it changes the nodes and uses integral Boolean measurement fibers. Constant-free prime-field circuits of polynomial formal degree follow after homogeneous-component processing. NN-C034 subsequently supplies integral-gate circuits over Z as well. Polynomial-size IPS certificates, efficient PIT verification and a VP/VNP collapse remain unproved. Exponential monomial-PC lower bounds are still compatible with the resulting shared circuits.

\paragraph{Geometric and representation dictionary.}
For every prime $p$, let $B=\mathbb F_p[y]$ and let
\[
 \mathscr A=B[x_1,\ldots,x_n]/(x_i^2-x_i),\qquad
       a=\prod_f\Phi_f(x).
\]
This is a finite free split algebra of rank $2^n$, with the Boolean orthogonal idempotents as a basis. Its multiplication trace $\operatorname{Tr}_{\mathscr A/B}(a)$ is the universal partition polynomial, because multiplication is diagonal in that basis. The theorem constructs a prime-field arithmetic circuit for this relative trace in size $S(n)$, naturally under arbitrary substitutions of the parameters. The auxiliary scalar extension is finite separable and free; the proof computes the trace after that extension and recovers its original scalar coordinate by an explicit basis. This descent is an identity of regular functions on the full parameter scheme, not pointwise interpolation on its rational points.

For an unsatisfiable local verifier, the same construction supplies an ordinary ideal identity in $\mathbb F_p[x]$ and polynomial maps identifying its associated smooth cubic with affine space over every extension field of characteristic $p$. The maps and inverses are explicitly represented by subexponential circuits and have the degree bounds from NN-C029. Thus the characteristic-two geometric trivialization gate is also closed for this family. The exact classical finite-locally-free trace locator is EGA IV4,18.2.1--18.2.2, printed p111, especially base-change equation18.2.1.2. The integral-matrix identity, change of interpolation nodes, integer-fiber separation and coordinate descent are proved here; no additional Grothendieck archive locator is asserted. This family-level effective trace theorem supplies neither a canonical GCT multiplicity realization nor general polynomial-time geometric recognition.

\subsection{NN-C034 --- Cyclotomic interpolation gives integral constant-free symbolic circuits over every ring}
\begin{proposition}[NN-C034]
Fix factor arity $w$ and density $C$. The full integral universal partition polynomial
\[
 Z_H(y)=\sum_{x\in\{0,1\}^n}\prod_f y_{f,x|I_f}
                              \in\mathbb Z[y]
\]
of every $n$-variable, $m\leq Cn$-factor instance has a deterministically constructible division-free arithmetic circuit with integer constants, of size and bit construction time
\[
 S(n)=\operatorname{poly}(n)\,
       2^{O_{w,C}(n(\log\log n)^2/\sqrt{\log n})}.
\]
All integer constants have polynomial bit length. The circuit can be made constant-free, with only $0,1,-1$ as constant leaves, at polynomial overhead and with polynomial syntactic formal degree. For the universal polynomial one may take formal degree at most $\max(m,1)$ when subtraction gates are allowed and have the same degree rule as addition.

More generally, let $m\leq Cn$ factors be explicit monomial lists $\Phi_f=\sum_j y_{f,j}x^{\alpha(f,j)}$ of full bit length $L$, with independent coefficient blocks, declared scope at most $(\lceil\log_2n\rceil)^a$ and total input-degree budget $\sum_f\max_j|\alpha(f,j)|\leq Cn$, for fixed $C,a$. Their symbolic sum-product has size and bit construction time $\operatorname{poly}(n,L)2^{O_{C,a}(n(\log\log n)^2/\sqrt{\log n})}$, with the same integral and constant-free conclusions and formal weight degree at most $\max(m,1)$. This includes fixed-degree polylogarithmically sparse factors and preserves the original $n$ Boolean variables.

As a concrete application, the symbolic permanent of an $r\times r$ matrix with at most $Cr(\lceil\log_2r\rceil)^a$ allowed entries in total, for fixed $C,a$, has such an integral constant-free circuit in the original matrix-order parameter $r$, of formal degree at most $2r$. No individual row- or column-degree restriction is needed.

This is an identity over $\mathbb Z$, so the very same circuit computes the universal partition polynomial after base change to \emph{every} commutative ring with identity. No characteristic choice, finite-field construction, modulus factorization or rational denominator remains in that circuit.

Consequently the ordinary integral Nullstellensatz multipliers of NN-C029, and its mutually inverse integral affine-space coordinate maps in the empty case, have deterministically constructible integral-gate and constant-free circuits at the same subexponential parameter bound, with polynomial formal degree. For local verifiers the parameter is $v=n+s$ including wires; for fixed-width linear-clause CNFs the stronger original-$n$ certificate still omits Boolean axioms. These are ordinary polynomial identities over $\mathbb Z$, not only separate certificates over each finite field.

On the fixed-density 3-CNF family of NN-C021, this yields an exponential ratio between ordinary monomial-PC proof size and circuit-IPS proof size over the same field $\mathbb F_2$: the former is $2^{\Omega(n)}$, while the latter is at most $S(n)=2^{o(n)}$, with $n$ the original CNF variable count. The circuit certificates retain actual degree at most $n+1$ and polynomial formal degree.
\end{proposition}
\paragraph{Exact inputs.}
NN-C033. integral swap-matrix generating identity and the complete integer-measurement-fiber Taylor routine over a coefficient ring.

NN-C016 and NN-C029. universal bounded-arity partition polynomial; ordinary integral ideal multipliers, sparse-CNF identity and affine coordinate maps.

NN-C021. fixed-density mod-three Tseitin 3-CNF family with exponential monomial-PC lower bounds over F2; full repaired source reduction and size-degree deduction.

E: The Exponential Time Hypothesis is False, supplied LaTeX. Section3.1 independent-incidence total degree; complete Lemmas3.2--3.3, Proposition3.4 and Lemmas4.1--4.8; same discrete layouts, exact parameter checks and three exponential costs.

Servedio--Wan, Computing sparse permanents faster, author manuscript18 June2005. Section2, printed p3, first display: Rysers inclusion-exclusion identity; original definition p1.

Erdos, Beweis eines Satzes von Tschebyschef, Acta Litt. Sci. Szeged5(1932),194--198. Section5, printed pp197--198: for every integer n>1 a prime lies strictly between n and2n; Q=1 handled directly.

EGA IV, fourth part, Publications Mathematiques IHES32 (1967). Printed p.111,18.2.1--18.2.2; base-change equation18.2.1.2; primary PDF page110 (zero-based109).

Integral interpolation construction below. Explicit two-prime cyclotomic quotient, integral unit inverses, Hasse cardinal polynomials and coefficient-size proof; proposed by root and developed here.

\begin{proof}
The argument uses the integral matrix and exact integer measurement fibers proved in NN-C033. The new step is to carry out all interpolation inside an explicit finite free integral algebra with sufficiently many nodes whose differences are integral units. Coordinate simulation will then return to $\mathbb Z$.

Let $Q\geq2$ be the number of interpolation nodes in one layout. Choose primes
\[
                   Q<p<2Q<q<4Q.
\]
Bertrand's theorem, in the exact primary source cited above, supplies them. Trial division of every candidate up to $4Q$ finds them deterministically in polynomial time in $Q$; no fast primality theorem is needed. For $Q=1$ take $p=2,q=3$. Put $N=pq$ and
\[
 \Phi_N(T)=\frac{(T^N-1)(T-1)}{(T^p-1)(T^q-1)},\qquad
 a=(p-1)(q-1),\qquad A=\mathbb Z[\zeta]=\mathbb Z[T]/(\Phi_N(T)).
\]
Here the notation for $A$ means this explicit quotient; no computation of a ring of integers is required. We verify the polynomial facts used about it.

Over $\mathbb C$, the denominator's roots occur in the numerator with at least the same multiplicities: the common root one has multiplicity two, and the remaining $p$th and $q$th roots occur once. Thus the quotient is a monic polynomial of degree $a$. Polynomial division by a monic integer polynomial takes place over $\mathbb Z$, and its complex remainder is zero, proving $\Phi_N\in\mathbb Z[T]$. Its roots are exactly the primitive $N$th roots, each with multiplicity one, because $N$ is the product of the two distinct primes. Its value at one is
\[
                    \Phi_N(1)=\frac{N}{pq}=1,
\]
by canceling the two factors $T-1$, or equivalently by evaluating the resulting geometric sums. Also $\Phi_N$ divides $T^N-1$ over $\mathbb Z$, so $\zeta^N=1$ in $A$. The algebra $A$ is free over $\mathbb Z$ with the explicit basis
$1,\zeta,\ldots,\zeta^{a-1}$, of rank $a=O(Q^2)$. In particular its scalar inclusion is injective. Irreducibility of $\Phi_N$ is not needed anywhere in the argument.

If $\gcd(d,N)=1$, taking $d$th powers permutes the primitive roots. Therefore $\Phi_N(T)$ divides $\Phi_N(T^d)$ for a nonnegative representative of $d$ modulo $N$. Monic division again shows the divisibility over $\mathbb Z$. Hence $\Phi_N(\zeta^d)=0$ in $A$. Define the integral polynomial
\[
                     U(T)=\frac{\Phi_N(T)-1}{T-1}.
\]
It is integral because its numerator vanishes at one. The identity
\[
                   (1-\zeta^d)U(\zeta^d)=1
                 \qquad(\gcd(d,N)=1)
\]
proves that $1-\zeta^d$ is an integral unit with a displayed inverse.

Take the $Q$ nodes $\alpha_s=\zeta^s$, $1\leq s\leq Q$. Each is a unit because $\zeta^N=1$. For $i\ne j$, the nonzero integer $i-j$ has absolute value less than $Q<p,q$, so it is coprime to $N$. Thus
$\alpha_i-\alpha_j=\zeta^j(\zeta^{i-j}-1)$ is also a unit, with an inverse obtained from the displayed formula and a power of $\zeta$. All these nodes and inverses are elements of the original integral quotient $A$; no localization has occurred. The quotient may ramify after reduction at $p$ or $q$, but this has no effect on these unit identities.

The integral swap matrix $M_0$ from NN-C033 now gives rows
\[
 \varphi_{u,y}(T)=\prod_{s\in E_{u,y}}(1-T/\alpha_s),\qquad
 P_u(T)=\delta_u(I-M_0)+\sum_yM_0[u,y]\varphi_{u,y}(T)\delta_y
                             \quad\hbox{in }A[T]^{1\times D}.
\]
Division by $\alpha_s$ here means multiplication by its explicit integral unit inverse. Their degree and zero-at-mismatched-node properties are unchanged. The exact integer-kernel identity for $M_0$ gives invariance on the same integer Boolean measurement fibers after mapping integers into $A$. This is all that the fiber proof of E Lemma3.2 uses. The ring-valued Boolean-difference identity of Lemma3.3, fixed substitutions, and Taylor truncations therefore compute all jets
\[
                  h_i(z)=H(\alpha_i+z)\bmod z^J
                              \quad\hbox{in }A[y][z]/(z^J)
\]
with the same operation bound as NN-C033. Every equality is formal in the independent weights $y$, and $H(0)=Z_H(y)$ as an element of $A[y]$.

We give the entire final Hasse interpolation explicitly, so that there is no unexamined inverse-matrix or denominator step. Put
\[
 B_i(T)=\prod_{j\ne i}
             \left((T-\alpha_j)(\alpha_i-\alpha_j)^{-1}\right)^J,
 \qquad b_i(z)=B_i(\alpha_i+z).
\]
Then $b_i(0)=1$. Its inverse modulo $z^J$ is the polynomial
$c_i(z)=\sum_{k=0}^{J-1}c_{i,k}z^k$ determined by
\[
 c_{i,0}=1,\qquad
 c_{i,k}=-\sum_{l=1}^k[z^l]b_i(z)\,c_{i,k-l}\quad(1\leq k<J).
\]
This recurrence uses only integer-algebra additions and multiplications. Define
\[
 \widetilde H(T)=\sum_{i=1}^Q B_i(T)
      \left([h_i(z)c_i(z)]_{<J}\right)_{z=T-\alpha_i}.
\]
Its degree is at most $(Q-1)J+J-1=QJ-1$. At node $\alpha_i$ its jet is $h_i$, because $b_ic_i=1\bmod z^J$; at every other node the factor $B_i$ vanishes to order at least $J$. The ideals $(T-\alpha_i)^J$ are pairwise comaximal in $A[T]$, since the differences are units. Their intersection equals their product, generated by the monic polynomial $\prod_i(T-\alpha_i)^J$ of degree $QJ$. A polynomial of smaller degree in this ideal is zero, even over a ring with zero divisors, by monic polynomial division. Therefore the jets uniquely determine a polynomial of degree less than $QJ$. As $\deg H<QJ$, one has $\widetilde H=H$ as a formal identity in $A[y][T]$.

In particular $H(0)$ is computed by the displayed formula using only fixed integral scalars and the jet entries. All interpolation constants are in $A$. This uses neither Gaussian pivots nor division by a derivative factorial, and remains valid after every coefficient-ring homomorphism. No assertion about equality of numerical functions is substituted for polynomial equality.

Here are explicit size and coefficient-height bounds for the scalar computations. The rank $a$ and $N$ are $O(Q^2)$. Since the roots of $\Phi_N$ have complex absolute value one, the sum of absolute values of its coefficients is at most $2^a$ by the elementary-symmetric bounds $\binom ak$. Monic division therefore constructs it in polynomial bit time. Division of $\Phi_N-1$ by $T-1$ gives $U$ of coefficient norm at most $a2^{a+1}$. Evaluating its finitely many needed powers of $\zeta$, and reducing to the monic basis, involves only polynomially many powers and operations of degree polynomial in $Q$.

More quantitatively, let $\|u\|$ be the sum of absolute coordinates of an element $u\in A$ in the monic basis, and let
\[
 K=\max\left(1,\max_{0\leq i,j<a}\|\zeta^{i+j}\|\right).
\]
Then $\|uv\|\leq K\|u\|\|v\|$. Repeated monic reduction of powers of degree at most $2a-2$, using coefficient bounds $2^a$, gives $\log_2 K=O(a^2+1)$. The unit inverses above have coordinate bit height polynomial in $Q$: reduce exponents first modulo $N$, using $\zeta^N=1$, and then reduce the polynomial of degree below $N$ in the monic basis. This requires at most $N$ monic reductions, each with coefficient growth bounded by $1+\|\Phi_N\|$, so its logarithmic norm is polynomial in $Q$.

The coefficients of $B_i(\alpha_i+z)$ are formed from $O(QJ)$ factors with these bounded norms. Extend $\|\cdot\|$ to polynomials by summing coefficient norms. The same multiplication inequality applies, so these coefficients have logarithmic norm polynomial in $Q,J$. Let $B\geq1$ bound the norms of the nonconstant coefficients of $b_i$, uniformly in $i$. If $C_k=\|c_{i,k}\|$, the recurrence gives
\[
 C_0=1,\qquad C_k\leq KB\sum_{l=0}^{k-1}C_l,
          \qquad C_k\leq KB(1+KB)^{k-1}\quad(k\geq1).
\]
Thus $\log(1+C_k)=O(J\log(2+KB))$, still polynomial in $Q,J$. The cardinal-polynomial formula then gives polynomial bit bounds for every final scalar coefficient as well. All scalar computations take polynomial bit time: they involve polynomially many operations on arrays of polynomial length and on integers of these polynomial bit heights. The same estimates apply to rows, shifted Taylor coefficients and a local factor evaluator, because their degrees and scope sizes are polynomially bounded. No integer of superpolynomial bit length is being hidden in a unit inverse.

Use the same high-incidence branches, labels, integer measurements and layout parameters as in NN-C033. Its proof gives $S(n)$ operations over $A[y]$. Each fixed-layout scalar table and each local coefficient evaluation has polynomial bit cost; repeated local evaluations are included in the $S(n)$ operation bound. Store each symbolic result as an arithmetic-gate reference. Replacing an $A$-valued gate by its $a$ coordinates gives an integer arithmetic circuit: addition is coordinatewise and multiplication uses the integer structure constants for the monic quotient basis. There are at most polynomially many coordinate gates per original operation. Their constants have the polynomial bit bounds just proved. Hence the total size and construction bit time remain $S(n)$.

The coordinate descent is exact. The polynomial initially specified as $Z_H(y)$ has integer coefficients, and every step proved that the computed element equals its image in $A[y]$. Since $A$ is free with basis beginning in one, that image has first coordinate $Z_H(y)$ and all other coordinates zero. Take the first output coordinate. This gives the original integral polynomial with integer-gate syntax, without computing its exponentially many coefficients and without a symbolic identity oracle. Its image under any map $\mathbb Z\to R$ is therefore the universal partition polynomial over the arbitrary target ring $R$.

We finally justify the constant-free and formal-degree assertions without assuming that semantic integrality alone eliminates denominators. The integral construction just obtained already has only integer scalar constants of polynomial bit length. It may be arranged homogeneous in the factor-weight blocks: every local factor is linear in its own block, a group product uses exactly its assigned factors, every group sum retains that block degree, and the final Taylor and interpolation sums have total weight degree $m$. In constructing a local evaluator, precompute its scalar Taylor coefficients and then form the linear combinations of its factor weights; all these scalar coefficients have the height bounds above. Empty group products contribute one, and their sums are integer counts at most $2^n$, also of polynomial bit length. Thus all degree-zero operations may be performed as explicit scalar preprocessing. After coordinate simulation, each remaining product either combines positive-degree homogeneous pieces or multiplies one piece by a known integer scalar. Its degree bookkeeping is known from this construction and requires no identity test.

For an integer scalar of bit length $L$, multiplication of a symbolic input by that scalar can be implemented with $O(L)$ doubling and addition/subtraction gates. These gates do not increase formal degree. Integer constant outputs themselves can be built from $1$ and $-1$ by such additions, with formal degree one. Replacing all scalar multiplications this way leaves only products of positive-degree pieces. Induction then bounds formal degree by the positive weight degree, and hence by $m$; the constant case uses degree one. This gives the stated constant-free bound at polynomial overhead. If one's convention permits only addition and multiplication, maintain both signs of each value; the same construction has formal degree at most $2\max(m,1)$, which is still polynomial.

Apply these integral circuits to the exact universal-partition specializations in NN-C029. Preliminary multipliers, their constantly many local specializations, and their fixed-scope integral division corrections now all use integral gates. Their ordinary identity over $\mathbb Z$ is the already-proved first-violation and local-remainder identity, with exactly the same specified multiplier polynomials. Each substituted weight is a constant or a degree-one expression, and each local correction has bounded degree and size, so the constant-free formal-degree bound stays $O(v)$, even without identifying it with the sharper actual degree bound $\deg a_i\leq v$. The IPS polynomial adds its placeholder variables with one multiplication per term. The explicit affine coordinate maps use only these multipliers and the displayed low-degree formulas of NN-C029, preserving subexponential circuit size and polynomial formal degree. This proves all assertions over $\mathbb Z$ and hence after arbitrary base change.

\paragraph{Extension to sparse polynomial factors with growing scopes.}
The same integral construction applies beyond constant arity.
Fix integers $C,a\ge1$.  Suppose there are $n\ge2$ Boolean variables
and $m\le Cn$ factors, each given by an explicit monomial list
\[
 \Phi_f(x)=\sum_{j=1}^{\ell_f}y_{f,j}x^{\alpha(f,j)},
 \qquad d_f=\max_j|\alpha(f,j)|,
 \qquad \sum_f d_f\le Cn.
\]
The coefficient symbols $y_{f,j}$ are independent; repeated
monomials in a list are permitted.  An empty list means the zero
factor and is handled immediately.  Let $L$ be the full bit length
of these lists, including their variable indices and binary
exponents.  Suppose each declared scope $I_f$ has size at most
$\lambda^a$, where $\lambda=\lceil\log_2 n\rceil$.
The degree bound is imposed on these explicitly supplied
polynomials before Boolean reduction.  Then
\[
 Z_{\Phi}(y)=\sum_{x\in\{0,1\}^n}\prod_f\Phi_f(x)
\]
has an integral division-free circuit of size and deterministic
bit construction time
\[
 \operatorname{poly}(n,L)\,
 2^{O_{C,a}(n(\log\log n)^2/\sqrt{\log n})}.
\]
Its integer constants have polynomial bit length in $n,L$.
It admits the same constant-free compilation of formal degree
at most $\max(m,1)$, with subtraction gates assigned the
addition degree rule.  The output is the formal polynomial
over $\mathbb Z[y]$, so this single circuit identity transfers
to every commutative coefficient ring.

In particular, fixed-degree factors with a polylogarithmic
number of monomials satisfy these hypotheses after changing
the fixed constants: if the degree is at most $k$ and the
list length at most $(\log n)^b$, the scope is at most
$k(\log n)^b$, and the total degree is at most $kCn$.
This preserves the original number $n$ of Boolean variables.
There are no auxiliary verifier variables, and no truth table
on the entire growing scope is constructed.

We prove the exact quantitative interface.  Append one
dummy factor equal to one, with declared scope $\{i\}$, for
each variable $i$.  These factors contribute no coefficient
symbols and have input degree zero; they only ensure that
every variable has a declared incidence.  If $M_0$ is the
resulting total scope length, then
\[
              n\le M_0\le n+Cn\lambda^a.
\]
Put
\[
 B_0=\max(1,\lceil M_0/n\rceil)\lceil\sqrt\lambda\rceil,
 \qquad H=\{i:\deg_{\rm inc}(i)>B_0\}.
\]
Double-counting incidences gives
$|H|\le M_0/B_0\le n/\sqrt\lambda$.
Sum over all $2^{|H|}$ assignments to these variables.  On a
branch, substitute the chosen bits in the explicit monomial
lists and remove them from each declared scope.  Keep the
remaining declared scopes even if some monomial becomes
zero.  This uses only the syntactic list and chosen bits,
with no test of a formal polynomial identity.

Let $n'=n-|H|$.  For sufficiently large $n$ one has
$n/2\le n'\le n$.  Each surviving variable has incidence at
most $B_0$, every scope still has size at most $\lambda^a$,
and the dummy factors ensure $n'\le M'\le B_0n'$ for the
remaining total incidence $M'$.  Factors with empty scope
can be multiplied outside the branch's counting routine;
their polynomials in the $y$ variables are explicit linear
forms.  This removal does not change $M'$.

Choose a fixed integer $\beta$ larger than $a+2$, increasing
the threshold of the exceptional finite range to absorb
$C$.  With $\lambda'=\lceil\log_2 n'\rceil$, both the maximum
scope and $B_0$ are at most $(\lambda')^\beta$ for all
sufficiently large $n$.  E Lemma~4.8, with
$\nu'=\lceil\log_2\lambda'\rceil$, therefore supplies
\[
 r=\left\lfloor\frac{\sqrt{\lambda'}}{(\nu')^2}\right\rfloor,
 \quad Q=\left\lfloor\frac{\lambda'}{128\nu'}\right\rfloor,
 \quad R=(\beta+3)\lceil\sqrt{\lambda'}\rceil,
 \quad \tau=\frac1{r\lambda'},\quad D=r^{2Q}.
\]
Its exact numerical checks and E Lemma~4.7 construct the
layout in polynomial time.  Input sizes failing those
checks form a finite range depending only on the fixed
constants; direct expansion handles that range within
the asserted asymptotic bound.

The independent-incidence degree in E Section~3.1 is bounded
by the sum of the factor degrees, because replacing the
same Boolean variable in different factors by distinct
indeterminates does not increase any factor's total degree.
Consequently, on every branch,
\[
       d\le\sum_f d_f\le Cn\le 2Cn'.
\]
This is the crucial degree hypothesis; the possibly larger
number $M'$ of incidences is not substituted for $d$.
The interpolation polynomial has degree at most $dR$.
Taking $q=\lfloor dR/Q\rfloor$ and $J=q+1$ gives $QJ>dR$.
E Proposition~3.4 has the explicit enumeration bound
\[
       2Qr\,C_q(2D+1)^{2\tau n'}2^{2n'/r},\qquad
       C_q=\sum_{s=0}^{\min(q,n')}2^s\binom{n'}s.
\]
Put $\eta=R/Q$.  The source's finite-difference estimate
$\eta^qC_q\le(1+2\eta)^{n'}$ yields
\[
 \log_2 C_q
 \le q\log_2(Q/R)+n'\log_2(1+2R/Q)
 =O_C\!\left(n'\frac RQ\log\frac{2Q}{R}\right).
\]
Together with E Lemma~4.8 this bounds all three exponential
costs by
\[
 O_{C,a}\!\left(\frac{n'(\nu')^2}{\sqrt{\lambda'}}\right).
\]
The additional branching cost has logarithm
$|H|\le n/\sqrt\lambda$.  Since $n'=\Theta(n)$, their sum
is $O_{C,a}(n(\log\log n)^2/\sqrt{\log n})$.

It remains to verify the factor and coefficient costs,
which an arity statement alone would not justify.  After
the fixed Boolean substitutions in E Lemma~3.3, evaluate
each listed monomial directly in the truncated algebra
$A[z]/(z^J)$, using its binary exponents or at most
$\sum_i\alpha(f,j)_i\le Cn$ multiplications.  Sum these
values with their formal coefficients $y_{f,j}$.
This takes polynomially many algebra operations in
$n,J,L$; it never enumerates the $2^{|I_f|}$ local
assignments.  Before multiplication by the formal
coefficients, each such substituted monomial is a
scalar Taylor polynomial.  The rows and unit inverses
have the polynomial coordinate heights already proved
for the cyclotomic construction.  A product of at most
$Cn$ such scalar inputs, followed by Taylor truncation,
still has polynomial bit height by the same multiplicative
coordinate-norm bound.  Adding at most $L$ listed terms
increases logarithmic norm by at most $O(\log(L+1))$.
All local scalar preprocessing is therefore polynomial
in $n,L$, and $Q,R,J,D$ are polynomially bounded in $n$.

Use the integral swap matrix, integer measurement fibers,
unit-separated cyclotomic nodes and explicit Hasse
interpolation proved above.  Their correctness depends on
the layout and the degree estimate, not on constant scope
size.  They give formal identities over $A[y]$ at each
stage.  The same coefficient-independent choices of
representatives and branches allow circuit unrolling,
and taking the coordinate of one in $A$ descends the
formal output to $\mathbb Z[y]$.  Each genuine factor
remains linear in its own coefficient block.  Group
products, sums, empty-scope multiplication and the final
sum over branches retain the same block degree.  The
constant-free compilation therefore has formal degree
at most $\max(m,1)$, just as in the constant-arity case.
This finishes the construction and its bit bound.

The hypotheses delimit the extension.  A general truth
table on $\operatorname{polylog} n$ variables can have
interpolation degree equal to its scope size; $O(n)$ such
factors need not have total degree $O(n)$.  In that case
the actual $q=\lfloor dR/Q\rfloor$ must be retained and
the displayed bound has not been proved.  Nor does a
coefficient-dependent compression of a numerical equation
system automatically preserve this universal formal
polynomial in independent parameters.

\paragraph{A sparse permanent application in the matrix-order parameter.}
Fix $a\ge1$ and $r\ge2$.  Let $\Omega\subseteq[r]\times[r]$ be a prescribed
matrix support with at most $(\lceil\log_2 r\rceil)^a$ entries in
each row, and give the allowed entries independent variables
$a_{ij}$.  The matrix has zero elsewhere.  Its permanent polynomial
has a deterministically constructible integral constant-free circuit
of size and bit construction time
\[
 \operatorname{poly}(r)\,
       2^{O_a(r(\log\log r)^2/\sqrt{\log r})},
\]
with formal degree $O(r)$.  No column-degree restriction is imposed.
The same claim allows a fixed multiplicative constant in the
row-support bound, by increasing the fixed scope exponent and
handling bounded sizes directly.  The parameter is the matrix
order $r$, rather than the number of allowed matrix entries.

Import Ryser's classical inclusion--exclusion identity in the
exact form reproduced in Servedio--Wan, \emph{Computing sparse
permanents faster}, author manuscript dated June~18, 2005,
Section~2, printed p.~3, first display:
\[
 \operatorname{per}(A)=(-1)^r\sum_{S\subseteq[r]}(-1)^{|S|}
                          \prod_{i=1}^r\sum_{j\in S}a_{ij}.
\]
The primary manuscript is
\url{https://www.cs.columbia.edu/~atw12/papers/perm.pdf}.
It states the identity for real or complex matrices.  Both sides
are polynomials with integer coefficients in independent matrix
entries, so agreement for all complex entries is an identity in
$\mathbb Z[a_{ij}]$ and hence after every ring base change.
This uses that a polynomial vanishing on all complex tuples is
the zero polynomial; it does not reprove the classical formula.

For the characteristic vector $x$ of $S$ one has
$(-1)^{|S|}=\prod_{j=1}^r(1-2x_j)$.  Consequently
\[
 \operatorname{per}(A)=(-1)^r
 \sum_{x\in\{0,1\}^r}
       \left(\prod_{j=1}^r(1-2x_j)\right)
       \left(\prod_{i=1}^r\sum_{j:(i,j)\in\Omega}a_{ij}x_j\right).
\]
This is a sum-product on exactly $r$ Boolean variables with $2r$
linear factors.  Its independent-incidence degree is at most
$2r$, each row factor has polylogarithmic scope, and the unary
sign factors have scope one.  The list description has polynomial
size.  A zero row gives the zero permanent immediately, and the
empty matrix has permanent one.  Apply the growing-scope theorem
of NN-C034 to the nontrivial cases, specializing the unary factor
coefficients to $1,-2$.  These are fixed integer constants, with
$-2=(-1)+(-1)$ available in the constant-free syntax.  Its integral
descent proves the full permanent identity; it does not merely
compute its values modulo a prime.  The factor-block formal degree
is at most $2r$, and the outer sign can be implemented by subtraction
from zero.  This proves the claimed circuit and construction bounds.

An allowed column may occur in every row.  The total-incidence
branching already proved in NN-C034 removes at most
$r/\sqrt{\log r}$ such high-incidence variables before applying
the layout, so unbounded column degree introduces no extra witness
parameter.  In particular, the conclusion covers supports with
$r\operatorname{polylog} r$ allowed entries at the original matrix
order.  It makes no claim for an arbitrary dense support: its
row scopes can have size $r$, outside the hypothesis.

\paragraph{Only a total support bound is necessary.}
Fix $C,a\ge1$ and $r\ge2$.  For a prescribed support
$\Omega\subseteq[r]\times[r]$ with
\[
                     |\Omega|\le Cr(\lceil\log_2r\rceil)^a,
\]
the full permanent polynomial has a deterministically
constructible integral constant-free circuit of size and
bit construction time
\[
 \operatorname{poly}(r)\,
       2^{O_{C,a}(r(\log\log r)^2/\sqrt{\log r})},
\]
and formal degree at most $2r$, with the subtraction convention
used above.  Thus no bound on any individual row or column
degree is needed.  The original parameter is still the
matrix order $r$.

Here is the full reduction to the preceding row-scope result.
Let $\lambda=\lceil\log_2r\rceil$, set
$T=\lceil\lambda^{a+2}\rceil$, and let $H$ be the set of
rows containing more than $T$ allowed entries.  If $h=|H|$,
then counting allowed entries gives
\[
                          h\le Cr/\lambda^2.
\]
For sufficiently large $r$, this is at most $r/2$.
Expand the definition of the permanent by first specifying
the columns matched to the rows of $H$.  It gives the exact
integral polynomial identity
\[
 \operatorname{per}(A)=
 \sum_{\substack{\pi:H\hookrightarrow[r]\\(i,\pi(i))\in\Omega}}
       \left(\prod_{i\in H}a_{i,\pi(i)}\right)
       \operatorname{per}\bigl(A_{[r]\setminus H,\,[r]\setminus\pi(H)}\bigr).
\]
This is simply the partition of the permutations in the
permanent definition according to their restrictions to
$H$; each full permutation occurs in exactly one summand.
There is no sign and no multiplicity factor.  If $H$ is
empty the sum has its one empty restriction, and its
weight is one.  Enumerating at most $r^h$ functions and
retaining the allowed injections constructs the sum
deterministically with polynomial overhead per term.

Every minor has order $r'=r-h\ge r/2$.  Its remaining
row supports have size at most $T$, and deleting columns
cannot increase them.  Since
$\lceil\log_2r'\rceil\ge\lambda-1$, one has
\[
                   T\le(\lceil\log_2r'\rceil)^{a+3}
\]
for sufficiently large $r$.  The preceding row-scope
permanent construction therefore applies to every minor,
with a fixed changed scope exponent.  Zero rows can be
detected from the prescribed support and give zero terms.
The finitely many exceptional matrix orders, including
orders zero and one, are computed directly.

The logarithm of the number of restrictions is at most
\[
                   h\log_2r=O_C(r/\log r),
\]
which is absorbed by
$O_{C,a}(r(\log\log r)^2/\sqrt{\log r})$.
The minor construction has that same bound because
$r'=\Theta(r)$.  Multiplication by the $h$ selected entry
variables and summation of all terms add only the stated
overhead.  Each minor circuit has formal degree at most
$2r'$, so each term has formal degree at most
$2r'+h\le2r$, and summation takes the maximum.  All scalar
constants retain their polynomial bit lengths.  This
proves the total-support theorem as an integral formal
identity, and therefore over every commutative coefficient
ring by base change.

There is a concrete proof-size separation on the family already audited in NN-C021. Use its mod-three Tseitin 3-CNF contradictions, with $n$ denoting the number of variables of the resulting CNF and with $O(n)$ clauses. Every ordinary polynomial-calculus refutation over $\mathbb F_2$ has explicitly written monomial size at least $2^{\kappa n}$ for some fixed $\kappa>0$. The sparse-CNF specialization just proved constructs a constant-free Hilbert-like IPS circuit over $\mathbb F_2$ of size $S(n)=2^{o(n)}$ for the same clause polynomials, without introducing verifier wires or Boolean axioms into that IPS identity. If $M_{\rm PC}(n)$ and $C_{\rm IPS}(n)$ are the minimum sizes in these two specified models, then for all sufficiently large $n$,
\[
       \frac{M_{\rm PC}(n)}{C_{\rm IPS}(n)}
           \geq\frac{2^{\kappa n}}{S(n)}
           \geq 2^{\kappa n/2}.
\]
This is the precise exponential ratio. The numerator counts explicit monomials throughout an ordinary PC derivation; the denominator counts gates in a circuit IPS certificate. The latter certificate has actual polynomial degree at most $n+1$ and a constant-free circuit of formal degree $O(n)$. These degree statements neither identify the two size measures nor contradict the linear PC degree lower bound. The deduction uses the classical monomial-PC lower bound and E's designated algorithmic input, and implies no polynomially bounded proof system or $\mathrm{NP}=\mathrm{coNP}$ collapse.
\end{proof}
\paragraph{Status and complexity impact.}
Complete integral cyclotomic-node, Hasse interpolation, height, symbolic descent and certificate proof. Independent complete deduction-proof review by root: passed; imported input theorems are not claimed independently proved. The remaining integral-gate obstacle in NN-C029 and NN-C033 is closed for the stated symbolic families, including polynomial-list factors with polylogarithmic scopes and linear total input degree. List length and coefficient bit cost are explicit; generic growing-scope truth tables need not satisfy the degree budget. The construction is over Z itself, with polynomial-bit integer constants and a constant-free presentation of polynomial formal degree, so it transfers to every coefficient ring through one universal identity. It makes actual integral ideal multipliers and empty-case geometric coordinate maps effective at E's subexponential bound. It does not give polynomial-size circuits or IPS certificates, deterministic polynomial-time PIT, NP=coNP, or $VP^0=VNP^0$. Dense permanent encodings and long verifier circuits retain their actual edge/wire parameter overhead. The finite free cyclotomic algebra is an explicit computational extension, not an assumed canonical GCT object.

\paragraph{Geometric and representation dictionary.}
Let $B=\mathbb Z[y]$ and
\[
 \mathscr A=B[x_1,\ldots,x_n]/(x_i^2-x_i),\qquad
                       a=\prod_f\Phi_f(x).
\]
The finite locally free multiplication trace $\operatorname{Tr}_{\mathscr A/B}(a)$ is the universal partition polynomial. The theorem supplies an integral arithmetic circuit for this regular function on $\operatorname{Spec}B$ in size $S(n)$, with a constant-free presentation of polynomial formal degree. The auxiliary extension $B\to A\otimes_{\mathbb Z}B$ is finite free and faithfully flat. It contains the required distinct unit-separated interpolation sections; these differences remain units after every base change, including ramified residue characteristics. Exact interpolation and coordinate projection descend the computed regular function to $B$.

Fix $C,a\ge1$.  Let $m\le Cn$ explicit polynomial lists
$\Phi_f=\sum_j y_{f,j}x^{\alpha(f,j)}$ have description
length $L$, declared scopes of size at most
$(\lceil\log_2n\rceil)^a$, and total degree budget
$\sum_f\max_j|\alpha(f,j)|\le Cn$.  Put $B=\mathbb Z[y]$ and
\[
 \mathscr A=B[x_1,\ldots,x_n]/(x_i^2-x_i),\qquad
                       u=\prod_f\Phi_f(x).
\]
The finite locally free algebra $\mathscr A$ has rank $2^n$.
Its multiplication trace
$\operatorname{Tr}_{\mathscr A/B}(u)=Z_\Phi(y)$, as is
seen in the orthogonal Boolean-idempotent basis.  This
regular function on $\operatorname{Spec}B$ has an integral
circuit of size and bit construction time
$\operatorname{poly}(n,L)2^{O_{C,a}(n(\log\log n)^2/\sqrt{\log n})}$,
with polynomial-bit integer constants and a constant-free
compilation of formal degree at most $\max(m,1)$, with $n$
the original number of Boolean coordinates.  For every map
$B\to R$ the same circuit computes the base-changed
multiplication trace over $R$; coordinate descent from
the finite free cyclotomic extension proves the identity
before any specialization.  The exact classical trace
and base-change interface is EGA IV4, 18.2.1--18.2.2,
printed p.~111, including equation 18.2.1.2.  This is an
effective theorem for these specified regular trace
functions, not a claimed polynomial-time algorithm for
general finite algebra traces or GCT multiplicities.

For a support $\Omega\subseteq[r]\times[r]$ satisfying
$|\Omega|\le Cr(\lceil\log_2r\rceil)^a$, put
$B=\mathbb Z[a_{ij}:(i,j)\in\Omega]$ and
$\mathscr A=B[x_1,\ldots,x_r]/(x_j^2-x_j)$.  The regular
function
\[
 \operatorname{Tr}_{\mathscr A/B}\left(
     (-1)^r\prod_j(1-2x_j)
                 \prod_i\sum_{j:(i,j)\in\Omega}a_{ij}x_j\right)
                              =\operatorname{per}(A)
\]
has an integral constant-free circuit of size
$\operatorname{poly}(r)2^{O_{C,a}(r(\log\log r)^2/\sqrt{\log r})}$
and formal degree at most $2r$, valid after every base
change.  Ryser's identity identifies this signed trace,
while the heavy-row decomposition constructs it using
smaller trace functions with polylogarithmic row scopes.

For the split torus with characters indexed by $\Omega$,
let $\mathcal V_\Omega$ have one basis vector of character
$\prod_i a_{i,\sigma(i)}$ per allowed perfect matching.
Its effective character is the same permanent polynomial.
The restriction partition above gives the actual direct-sum
decomposition
\[
 \mathcal V_\Omega\simeq
  \bigoplus_{\pi}
    \chi_\pi\otimes\mathcal V_{\Omega_\pi},
 \qquad \chi_\pi=\prod_{i\in H}a_{i,\pi(i)},
\]
where $\Omega_\pi$ is the remaining minor support and its
characters are included in the original torus.  Applying
the integral minor circuits to these summands gives the
stated character circuit; positivity of the decomposition
does not remove the virtual differences inside the minor
circuits.  The precise classical trace and character-ring
interfaces remain EGA IV4, 18.2.1--18.2.2, printed p.~111,
and the SGA3/SGA6 locators already stated in NN-S019.

For the unsatisfiable local-verifier scheme, the same circuits give an ordinary global ideal identity $\sum_i a_if_i=1$ over $\mathbb Z$ and the explicit mutually inverse maps identifying its smooth counting cubic with affine space. The maps have the actual degree bounds from NN-C029, integral coefficient circuits of size $S(v)$, and polynomial formal degree after removing all large constant leaves. Thus these are effective integral geometric trivializations, valid after arbitrary base change, with no denominator-prime exclusion. On the empty Boolean equation schemes of NN-C021, these circuit ideal certificates over $\mathbb F_2$ have size $2^{o(n)}$, while deriving their unit ideal by ordinary monomial polynomial calculus requires size $2^{\Omega(n)}$. This compares specified presentations of the unit-ideal proof, not different underlying zero schemes or different cohomological degrees. The exact classical multiplication-trace and base-change locator is EGA IV4,18.2.1--18.2.2, printed p111, equation18.2.1.2. All node, inverse, flat scalar-extension and polynomial-identity calculations are proved above; no unverified Grothendieck archive cote or GCT canonical-basis functor is invoked.

\subsection{NN-C035 --- Exact zero-set compression yields formal weighted CNF circuits in characteristic two}
\begin{proposition}[NN-C035]
Fix a clause width $k\ge2$. For any $k$-CNF $F$ on $n$ declared variables with $m$ clauses, the full satisfying-assignment polynomial
\[
 A_F(z)=\sum_{x\models F}\prod_{i=1}^n z_{i,x_i}
                       \in\mathbb F_2[z_{1,0},z_{1,1},\ldots,z_{n,0},z_{n,1}]
\]
has a deterministically constructible constant-free arithmetic circuit of size and bit construction time
\[
 (n+m)^{O_k(1)}\,
      2^{O_k(n(\log\log n)^2/\sqrt{\log n})}.
\]
Its formal degree is $O_k(n)$. The parameter is the original number of Boolean variables; no verifier wires are added. This is equality of formal polynomials, valid after base change to every commutative ring of characteristic two.

More generally, one may replace the product of the $n$ displayed unary weights by at most $Cn$ explicitly listed polynomial factors in $x$, with independent coefficient parameters, polylogarithmic scopes, and total input-degree budget at most $Cn$, for fixed $C$ and scope exponent. The same bound holds with a polynomial factor in their description length and with the constants depending on these fixed parameters.

In particular, for every graph $G$ on $n$ vertices, with no sparsity or maximum-degree restriction, its full multivariate independence polynomial
\[
                       I_G(y)=\sum_{S\text{ independent}}\prod_{i\in S}y_i
                                      \in\mathbb F_2[y_1,\ldots,y_n]
\]
has such a circuit in the original vertex parameter. The two-rail polynomial $A_F$ is formally nonzero exactly when $F$ is satisfiable, even though its value at all weights equal to one records only the parity of the satisfying count.
\end{proposition}
\paragraph{Exact inputs.}
E: The Exponential Time Hypothesis is False, supplied LaTeX. Lemma6.3 disjoint sparsification on the original declared variables; Lemmas6.1--6.2 and Proposition6.4, especially the exact no-false-survivor conclusion.

NN-C001. deterministic compression of quasilinear sparse F2 equation lists to at most6n equations, preserving their full Boolean zero set.

NN-C034. reviewed growing-scope extension: explicit polynomial lists, polylogarithmic scopes, linear total input degree, and one formal integral circuit identity over every coefficient ring.

EGA IV, fourth part, Publications Mathematiques IHES32 (1967). Printed p.111,18.2.1--18.2.2; base-change equation18.2.1.2; primary PDF page110 (zero-based109).

\begin{proof}
We first retain the exact scope of the imported compression theorem. Its rows are selected by numerical zero counts for the supplied Boolean equations. They depend on $F$, but they do not depend on any of the formal weight parameters introduced here. NN-C001 proves equality of the whole Boolean zero set, not just equality of its cardinality. That stronger assertion is what permits symbolic weights.

Handle an empty clause by outputting zero. A formula with no clauses has the explicit circuit $\prod_i(z_{i,0}+z_{i,1})$ in the two-rail case. Bounded exceptional sizes may be expanded directly. For the general case put $\lambda=\lceil\log_2 n\rceil$ and apply E Lemma6.3 with the rational parameter
\[
                         \sigma=1/\lceil\sqrt\lambda\rceil.
\]
The lemma produces a tree of at most $\operatorname{poly}_k(n)2^{O(n/\sqrt{\log n})}$ nodes on the same $n$ declared variables. Its leaf satisfying sets partition the satisfying set of $F$. In every leaf each variable has occurrence at most
\[
 B=O_k\bigl([\sigma^{-1}\log(2/\sigma)]^{k-1}\bigr)
                                  =(\log n)^{O_k(1)}.
\]
The input normalization costs $\operatorname{poly}_k(n,m)$ and the numerical parameter has $O(\log\log(n+2))$ bits. Thus every nontrivial leaf has $n(\log n)^{O_k(1)}$ clauses, while the number of Boolean variables is still $n$.

For a leaf $F_t$, take the usual clause-violation polynomial
\[
                  v_C(x)=\prod_{\ell\in C}(1-\ell(x))
                                         \quad\hbox{over }\mathbb F_2,
\]
where the affine literal polynomial has value one exactly when that literal is true. A clause has at most $k$ literals, so its expanded violation polynomial has degree at most $k$ and at most $2^k$ monomials. It vanishes exactly when the clause is satisfied. Apply NN-C001 to these equations. The resulting list
\[
                   p_{t,1},\ldots,p_{t,b_t},\qquad b_t\le6n,
\]
has degree at most $k$, polylogarithmic monomial sparsity and hence polylogarithmic scopes, and satisfies
\[
 x\models F_t\quad\Longleftrightarrow\quad
              p_{t,j}(x)=0\text{ for every }j
                     \qquad(x\in\{0,1\}^n).
\]
Its construction takes $\operatorname{poly}_k(n)2^{O_k(n\log\log n/\sqrt{\log n})}$ bit operations. This uses only the numerical equations, so all coefficient-parameter symbols remain untouched.

Now work in the polynomial coefficient algebra $\mathbb F_2[z]$. For a Boolean assignment, each $p_{t,j}(x)$ is either zero or one. Therefore
\[
                e_t(x)=\prod_{j=1}^{b_t}(1+p_{t,j}(x))
\]
is exactly the indicator of $F_t$ on the Boolean cube. Use the unary interpolation factors
\[
                   W_i(x_i)=(1-x_i)z_{i,0}+x_i z_{i,1}.
\]
The sum-product
\[
              Z_t(z)=\sum_{x\in\{0,1\}^n}e_t(x)\prod_iW_i(x_i)
\]
therefore equals $\sum_{x\models F_t}\prod_i z_{i,x_i}$ as a formal polynomial over $\mathbb F_2$. This assertion follows term by term before any evaluation of $z$.

There are at most $7n$ factors. The sum of their degrees in the Boolean inputs is at most $(6k+1)n$, and every scope is polylogarithmic. Each factor is explicitly listed, and the list lengths are polynomial in $n$. These are precisely the hypotheses of the growing-scope part of NN-C034. To use its universal coefficient form, first place independent formal coefficients on each monomial occurrence of these factors. Construct the universal integral circuit supplied there. Then apply the coefficient-ring homomorphism to $\mathbb F_2[z]$ which sends the constraint coefficients to their known zero-one values and the weight coefficients to the displayed linear parameters. Thus the descended circuit computes $Z_t$ as a formal polynomial. The intermediary integral universal identity is being specialized to characteristic two; the characteristic-two indicator equality is not asserted as an integer equality.

The block compilation in NN-C034 gives a constant-free circuit of formal degree $O_k(n)$ after these specializations: each genuine coefficient input becomes a constant or a linear expression in a weight parameter, and only $O_k(n)$ factor blocks are used. In characteristic two subtraction is addition. The resulting size and construction time per leaf are bounded by $\operatorname{poly}_k(n)2^{O_k(n(\log\log n)^2/\sqrt{\log n})}$. All representation and control choices are discrete or depend on the known clause coefficients; there is no symbolic identity test and no evaluation/interpolation in the weight parameters.

Sum these circuits over the sparsification leaves. Disjointness is essential: every satisfying assignment occurs in exactly one leaf, and every nonsatisfying assignment in none. Consequently the sum is exactly $A_F(z)$, with no multiplicity introduced by the decomposition. Adding the leaf-count exponent $O_k(n/\sqrt{\log n})$, and the smaller compression-construction exponent, preserves the displayed bound. An additive sum of the leaf circuits does not increase formal degree.

For the general weighted statement, replace $\prod_iW_i$ by the supplied factor product. The compressed constraint factors add at most $6kn$ to its total input degree and at most $6n$ to its factor count. The maximum scope remains a fixed power of $\log n$, and direct list evaluation has polynomial cost in the stated description length. NN-C034 therefore applies with changed fixed constants. The same indicator identity and disjoint leaf sum prove the formal weighted equality. In this version distinct assignments can yield the same weight monomial, whose coefficient correctly counts their contributions modulo two.

For the graph application take
\[
                         F_G=\bigwedge_{uv\in E(G)}(\neg x_u\vee\neg x_v).
\]
Its satisfying assignments are exactly the characteristic vectors of independent vertex sets. It has $n$ variables and $|E(G)|$ width-two clauses; no edge variable is introduced. Specializing $z_{i,0}=1$ and $z_{i,1}=y_i$ gives $I_G(y)$. The construction time is polynomial in $n+|E(G)|$ times the displayed original-$n$ exponential.

Finally, different assignments give different two-rail monomials, each with coefficient one. Hence $A_F$ is the zero formal polynomial if and only if the satisfying set is empty. The specialization $z_{i,b}=1$ collapses all these monomials to the same constant and records only their sum modulo two. These are different readout operations. The construction is mildly subexponential, not polynomial, so even a polynomial-time PIT algorithm applied to its output would not by itself yield a polynomial-time SAT algorithm. Nor does a circuit over $\mathbb F_2$ for a zero-one-coefficient polynomial automatically provide an integral circuit for its chosen zero-one coefficient lift: reduction loses multiples of two, and no coefficientwise lifting operation is supplied here. NN-C036 separately constructs the integral polynomial for this CNF family by slower disjoint sparsification and exact integer clause indicators, with a weaker exponent.
\end{proof}
\paragraph{Status and complexity impact.}
Complete characteristic-two symbolic-weight deduction proved. Independent complete deduction-proof review by root: passed; imported input theorems are not claimed independently proved. This closes the bounded-density restriction for formal weighted fixed-width CNF counting in characteristic two while retaining the original witness variables. Numerical compression is legitimately reused because it preserves the entire Boolean solution set and ignores weight parameters. The result is stronger than a parity value at a single specialization and weaker than a polynomial-size circuit theorem. This characteristic-two construction does not itself give a characteristic-zero circuit for the chosen zero-one coefficient lift; NN-C036 subsequently supplies a direct integral construction for this family at a weaker exponent. Neither result gives general PIT derandomization or a VP/VNP collapse. The dense-graph independence-polynomial example has an explicit original-vertex bound, unlike an encoding that charges every graph edge as a witness bit.

\paragraph{Geometric and representation dictionary.}
Let $Y_F$ be the finite etale subscheme of $\operatorname{Spec}(\mathbb F_2[x_1,\ldots,x_n]/(x_i^2-x_i))$ cut out by the clause-violation equations. Put $B=\mathbb F_2[z_{i,b}]$ and $\mathscr A_F=\mathcal O(Y_F)\otimes_{\mathbb F_2}B$. It is finite locally free over $B$, with one rank-one summand for every satisfying assignment. For
\[
                           w=\prod_i\bigl((1-x_i)z_{i,0}+x_i z_{i,1}\bigr),
\]
the multiplication trace $\operatorname{Tr}_{\mathscr A_F/B}(w)$ is $A_F(z)$ in the component-idempotent basis. For fixed clause width, this regular function has a deterministically constructible constant-free circuit of size $(n+m)^{O_k(1)}2^{O_k(n(\log\log n)^2/\sqrt{\log n})}$ and formal degree $O_k(n)$. This statement remains true after every base change to a characteristic-two coefficient ring. The same trace theorem holds for the explicitly listed weight factors with polylogarithmic scopes and linear total input degree stated above.

The compressed equations define the same closed subscheme inside the Boolean scheme, because its product coordinate algebra has ideals determined by their vanishing components. The sparsification leaves form a disjoint union decomposition of $Y_F$, and multiplication trace adds over that decomposition. The nonzero two-rail coefficient records an actual component; trace at all weights equal to one is only its rank modulo two. Every point of $Y_F$ is $\mathbb F_2$-rational, so Frobenius on $H^0_{\mathrm{et}}(Y_F\times\overline{\mathbb F}_2,\mathbb Q_\ell)$, for $\ell\ne2$, is the identity. This weight-zero purity does not compute the symbolic trace. The exact classical multiplication-trace and base-change locator is EGA IV4,18.2.1--18.2.2, printed p111, equation18.2.1.2. No reductive-group multiplicity or canonical GCT object is asserted.

\subsection{NN-C036 --- Slower sparsification gives integral symbolic weights and ideal certificates for arbitrary-density CNFs}
\begin{proposition}[NN-C036]
Fix $k\ge2$. For every $k$-CNF $F$ on $n$ declared variables with $m$ clauses, the integral satisfying-assignment polynomial
\[
       A_F(z)=\sum_{x\models F}\prod_{i=1}^n z_{i,x_i}
                      \in\mathbb Z[z_{i,b}:1\le i\le n,\ b\in\{0,1\}]
\]
has a deterministically constructible division-free integral circuit of size and bit construction time
\[
 (n+m)^{O_k(1)}
  2^{O_k\!\left(n(\log\log n)^{1+1/k}/(\log n)^{1/(2k)}\right)}
                                                    = (n+m)^{O_k(1)}2^{o(n)}.
\]
The circuit can be made constant-free with polynomial formal degree. It computes the formal polynomial over $\mathbb Z$, so the same circuit works after every commutative-ring base change. The number of original Boolean variables remains $n$.

The same assertion holds for a product of $O(n)$ explicitly listed weight factors with independent coefficient parameters, polylogarithmic scopes, and linear total input degree; include a polynomial factor in their description length and allow the constants to depend on the fixed scope and degree bounds. Thus arbitrary-density fixed-width CNFs admit symbolic integral weighted counting at a mildly subexponential original-variable bound. Arbitrary rational specializations of the weights can also be evaluated within the displayed time times a polynomial in their full bit length.

In particular the full multivariate independence polynomial of every $n$-vertex graph, with no edge-density restriction, has an integral constant-free circuit of size
\[
 \operatorname{poly}(n+|E|)\,
       2^{O\!\left(n(\log\log n)^{3/2}/(\log n)^{1/4}\right)}
\]
and polynomial formal degree. The characteristic-two construction of NN-C035 retains its sharper exponent. This integral construction does not lift an arbitrary characteristic-two circuit; it uses a different decomposition of the original Boolean problem.

For every unsatisfiable normalized width-$k$ CNF, the same bound constructs ordinary integral multipliers $h_j$ with $\sum_jv_jh_j=1$ for its clause-violation polynomials $v_j$, without Boolean axioms or extra variables. One has $\deg h_j\le n-|S_j|$ for clause scope $S_j$, and the resulting Hilbert-like IPS circuit has actual degree at most $n+1$ and polynomial formal degree. The associated hypersurface $t^2-t=\sum_jy_jv_j(x)$ admits explicit integral affine-space coordinate maps of the same circuit bound, of forward degree at most $n+2$ and inverse degree at most $n+2k+2$.
\end{proposition}
\paragraph{Exact inputs.}
E: The Exponential Time Hypothesis is False, supplied LaTeX. Lemma6.3 disjoint sparsification, exact occurrence bound B=O\_k((sigma\textasciicircum{}-1 log(2/sigma))\textasciicircum{}(k-1)); Section3.1 total input degree and Proposition3.4 with the actual q=floor(dR/Q); complete Lemmas4.7--4.8.

NN-C034. integral swap-matrix and cyclotomic Hasse routine, explicit polynomial coefficient-height bounds, formal coordinate descent and constant-free block compilation; the construction is reused below with a larger actual degree parameter.

NN-C035. faster characteristic-two formal weighted construction, retained as a distinct bound.

NN-C015. constant-free coefficient-norm induction at polynomial formal degree; used for exact rational specialization bit bounds.

Grochow--Pitassi, Circuit Complexity, Proof Complexity, and Polynomial Identity Testing: The Ideal Proof System. Author JACM manuscript Definition1.1 printed draftp1:4; Lemma4.2 full proof pp1:26--1:28, equations(4)--(6); classical first-violation identity.

NN-C029. classical first-violation coefficient factorization and the integral unit-ideal affine coordinate formulas.

EGA IV, fourth part, Publications Mathematiques IHES32 (1967). Printed p.111,18.2.1--18.2.2; base-change equation18.2.1.2; primary PDF page110 (zero-based109).

\begin{proof}
The construction uses the integral clause-satisfaction polynomial instead of compressing the constraints in characteristic two. For a clause $C$, let $\ell(x)$ be the affine integer polynomial for each literal and put
\[
                   c_C(x)=1-\prod_{\ell\in C}(1-\ell(x))\in\mathbb Z[x].
\]
For every Boolean assignment this factor is exactly one when $C$ is satisfied and zero otherwise. Its degree is at most the clause width and its expanded list has at most $2^k+1$ terms. These are literal integer equalities, valid after every coefficient-ring map. An empty clause gives output zero. A formula without clauses has the explicit two-rail circuit $\prod_i(z_{i,0}+z_{i,1})$. Bounded exceptional input sizes are expanded directly.

We choose a slower sparsification parameter than the one used for the characteristic-two theorem. Put
\[
 \lambda=\lceil\log_2 n\rceil,\qquad
 \nu=\lceil\log_2\lambda\rceil,
 \qquad s=\max\{t\in\mathbb Z_{\ge0}:t^{2k}\nu^{2k+2}\le\lambda\}.
\]
For sufficiently large $n$, one has $s\ge2$ and
\[
 \frac1s=\Theta_k\!\left(
             \frac{\nu^{1+1/k}}{\lambda^{1/(2k)}}\right)=o(1).
\]
The integer $s$ is computed exactly by integer comparisons in polynomial time in $\log n$ for fixed $k$. Sizes with $s<2$ form a finite range. Apply E Lemma6.3 with $\sigma=1/s$. Its leaf satisfying sets partition that of $F$, its tree has $\operatorname{poly}_k(n,s)2^{n/s}$ nodes, and its explicit leaf occurrence bound satisfies
\[
 B=O_k\bigl((s\log(2s))^{k-1}\bigr)
       =O_k(s^{k-1}\nu^{k-1}).
\]
We may increase $B$ to at least one. The choice of $s$ implies
\begin{equation}\label{eq:C036-balance}
     \frac{B\nu^2}{\sqrt\lambda}
       =O_k\!\left(\frac{s^{k-1}\nu^{k+1}}{\sqrt\lambda}\right)
       =O_k(1/s),
\end{equation}
because $s^k\nu^{k+1}\le\sqrt\lambda$. This estimate balances the cost per leaf against the number of leaves.

Fix a nontrivial leaf $F_t$. It has at most $nB$ nonempty clauses and total clause-incidence length at most $nB$. Its product of satisfaction factors has independent-incidence input degree at most $nB$, since each factor's degree is bounded by its number of literals. Multiply by the $n$ unary factors
\[
                       W_i(x_i)=(1-x_i)z_{i,0}+x_i z_{i,1}.
\]
The desired leaf polynomial is exactly
\[
 Z_t(z)=\sum_{x\in\{0,1\}^n}\left(\prod_{C\in F_t}c_C(x)\right)
                                      \prod_iW_i(x_i)
                          \quad\hbox{in }\mathbb Z[z].
\]
For the more general weight factors, their specified linear total-degree budget adds $O(n)$ to the same bound. In either case the number of factors is $O(n(B+1))$, the total input degree is $d=O(n(B+1))$, and the total list length is polynomial in $n$ and the supplied weight-list length. The crucial difference from the literal theorem statement of NN-C034 is that $B$ now grows. We therefore reuse its construction with the actual $d$, rather than invoking its linear-degree corollary outside its hypotheses.

The clause scopes are constant, and their maximum variable incidence is at most $B$. General weight scopes may be polylogarithmic and have uneven incidences. If needed, append the degree-zero dummy factors and branch on high-incidence variables exactly as in NN-C034, using the full total scope length $M$ and threshold
\[
                   \max(1,\lceil M/n\rceil)\lceil\sqrt\lambda\rceil.
\]
There are at most $n/\sqrt\lambda$ branched bits. Every remaining instance has $n'=n-O(n/\sqrt\lambda)=\Theta(n)$ variables, polylogarithmic scopes and maximum incidence, and input degree $d=O(n'B)$. The unbranched unary-weight case already has this layout regime. A fixed exponent $\beta$, depending only on $k$ and the fixed weight-scope exponent, bounds the maximum scope and incidence by $(\lceil\log_2n'\rceil)^\beta$ for all sufficiently large $n$. This is possible because $B$ and every added weight-incidence threshold are fixed powers of $\log n$.

Apply the exact layout parameters of E Lemma4.8 to $n'$, writing $\lambda'=\lceil\log_2n'\rceil$ and $\nu'=\lceil\log_2\lambda'\rceil$. They satisfy
\[
 \frac RQ=O_\beta(\nu'/\sqrt{\lambda'}),\qquad
 \log(2Q/R)=O(\nu'),\qquad
 \tau\log(2D+1)+\frac1r=O_\beta((\nu')^2/\sqrt{\lambda'}).
\]
The source constructs and checks the numerical hypotheses exactly; its finite exceptional range is handled by direct expansion. The interpolation polynomial has degree at most $dR$. Set
\[
                         q=\lfloor dR/Q\rfloor,\qquad J=q+1.
\]
E Proposition3.4 supplies the enumeration factor
$2QrC_q(2D+1)^{2\tau n'}2^{2n'/r}$, where
$C_q=\sum_{j\le\min(q,n')}2^j\binom{n'}j$.
Its finite-difference estimate with $\eta=R/Q$ gives
\[
 \log_2 C_q
   \le q\log_2(Q/R)+n'\log_2(1+2R/Q)
   =O\!\left(n'B\frac RQ\log\frac{2Q}{R}\right).
\]
Consequently all three exponential costs are
\[
           O_{k,\beta}\!\left(n'B\frac{(\nu')^2}{\sqrt{\lambda'}}\right)
                       =O_{k,\beta}(n/s)
\]
by \eqref{eq:C036-balance} and $n'=\Theta(n)$. The high-incidence branching contributes at most $n/\sqrt\lambda=O(n/s)$ to the exponent. In particular $q,J,Q,R,D$ remain polynomially bounded in $n$; no numerical parameter with exponentially long encoding has entered the construction.

We check the additional hypotheses of the integral implementation. Take the same finite free cyclotomic algebra, unit-separated nodes, integral swap matrix and exact integer measurement fibers as NN-C034. Their construction depends on the layout, not on having only $O(n)$ factors. A clause factor evaluates by its at most $2^k+1$ monomials. A supplied weight factor evaluates directly from its list. Every substituted monomial has degree at most $d=O(nB)$ and is a product of that many row Taylor inputs at worst. The coefficient-norm estimates in NN-C034 bound the bit height of each row and each unit inverse polynomially in $n$. Repeated multiplication of $O(nB)$ such scalar inputs, in a polynomial-length truncated Taylor array, adds only a polynomial to the logarithmic norm. The factor count $O(nB)$ and $B=\operatorname{polylog}n$ are themselves polynomial in $n$. Hence all local scalar preprocessing has polynomial bit cost. The explicit Hasse cardinal recurrence has polynomial bit cost in $Q,J$, exactly as before.

All identities in the integer-fiber separation, Boolean-difference formula, Taylor readout and cyclotomic interpolation hold formally in the weight variables. Store symbolic ring operations as shared gate references, and descend by taking the scalar coordinate in the finite free algebra. This constructs an integral circuit for $Z_t$, with size and bit construction time a polynomial factor times $2^{O_{k,\beta}(n/s)}$. There is no symbolic identity test and no reduction modulo a prime. For constant-free compilation one may first put independent coefficient symbols on every monomial of every factor, use the same factor-block construction, and then specialize each clause coefficient to its fixed integer value and each weight coefficient to its specified parameter. Its formal degree is $O(nB+n)$, since the factor count is $O(nB+n)$ and every such coefficient specialization is constant or linear. This is polynomial in $n$. The scalar-height estimates justify replacing all integer scalar multiplications by binary doubling and addition/subtraction at polynomial overhead. This is the same explicit integral compilation as NN-C034, now with a polylogarithmic factor-count multiplier.

Sum the resulting integral circuits over the disjoint sparsification leaves. Every satisfying assignment contributes its weight exactly once, as an integer polynomial, and no other assignment contributes. The output is therefore $A_F$ over $\mathbb Z$, or the stipulated general weighted sum. The tree count, leaf work and high-incidence branches all have exponent $O_{k,\beta}(n/s)$. Polynomial preprocessing and input normalization give the asserted full time bound. An additive sum of leaf circuits preserves their polynomial formal-degree bound. Formal equality over $\mathbb Z$ proves all coefficient-ring base changes simultaneously.

For completeness, rational specialization of the two-rail weights also has the asserted bit-time bound, without an unproved bound on intermediate arithmetic values. Let the circuit have $S$ gates and formal degree $D=\operatorname{poly}(n)$, and let $L$ be the total bit length of the rational inputs. Choose a common input denominator $Q_0$ by multiplying their denominators, so $\log Q_0\le L$. The constant-free coefficient-norm induction of NN-C015 gives coefficient norm at most $2^{SD}$ at every gate. After evaluation, a denominator divides $Q_0^D$ and the numerator has bit length $O(D(S+L+1))$. Exact rational arithmetic therefore evaluates all $S$ gates in $\operatorname{poly}(S,n,L)$ bit operations. Taking a fixed polynomial power of $S$ preserves the displayed $2^{O_k(n/s)}$ bound. The same argument applies to the general weight coefficients. We make no polynomial-working-space claim from this circuit evaluation argument.

For the independent-set example use the width-two CNF $\bigwedge_{uv\in E}(\neg x_u\vee\neg x_v)$ on the original vertex bits and then substitute $z_{i,0}=1$, $z_{i,1}=y_i$. Setting $k=2$ gives the stated exponent. Its coefficients are the integer zero-one indicators of independent subsets. This proves an integral circuit for their chosen lift by direct integer clause indicators; it is not a generic procedure for lifting circuits from $\mathbb F_2$ to $\mathbb Z$. The sharper characteristic-two bound of NN-C035 is therefore retained alongside this result.

\paragraph{Ordinary integral ideal certificates without a density bound.}
Let $F=C_1\wedge\cdots\wedge C_m$ be an unsatisfiable fixed-width
$k$-CNF on $n$ declared variables.  Use the standard clause encoding:
delete repeated literals, discard tautological clauses, and let
$v_j(x)$ be the product of the distinct falsifying literal factors
of $C_j$.  Empty clauses have $v_j=1$ and give an immediate
certificate.  Then one can deterministically construct integral
constant-free circuits for polynomials $h_j$ such that
\[
                 \sum_{j=1}^m v_j(x)h_j(x)=1
                              \quad\hbox{in }\mathbb Z[x_1,\ldots,x_n],
\]
without adding Boolean equations or verifier variables.  Their
total size and bit construction time are
\[
 (n+m)^{O_k(1)}
  2^{O_k(n(\log\log n)^{1+1/k}/(\log n)^{1/(2k)})}.
\]
If $S_j$ is the scope of $C_j$, the multipliers may be chosen
with $\deg h_j\le n-|S_j|$, so each product $v_jh_j$ has actual
degree at most $n$.  The circuit formal degree is polynomial
in $n$; it is not identified with the sharper actual degree.

The ideal identity is the classical first-violated-clause
certificate of Grochow--Pitassi, \emph{Circuit Complexity, Proof
Complexity, and Polynomial Identity Testing: The Ideal Proof
System}, author JACM manuscript, Lemma~4.2 and equations
(4)--(6), printed draft pp.~1:26--1:28; the primary file is
\url{https://www.cs.toronto.edu/~toni/Papers/jacm-gp.pdf}.
Its Definition~1.1 is on printed draft p.~1:4.  These are
the author manuscript's page markers, not an asserted match
with the final journal pagination.  The new consequence is
the constructive original-$n$ circuit bound for all clause
densities, obtained by the following exact factorization of
those classical multipliers.

Suppose there is no empty clause.  For each $j$, let
$b^{(j)}\in\{0,1\}^{S_j}$ be the unique assignment falsifying
$C_j$.  Restrict the prefix $C_1\wedge\cdots\wedge C_{j-1}$
by these fixed bits, keeping all the remaining $n_j=n-|S_j|$
declared variables.  Call the resulting width-at-most-$k$
CNF $R_j$.  Simplified false or empty prefixes are handled
by the same direct cases as in the main theorem.  Apply
NN-C036 to its full two-rail polynomial
\[
 H_j(z)=\sum_{a\models R_j}\prod_{i\notin S_j}z_{i,a_i},
\]
and specialize $z_{i,0}=1-x_i$, $z_{i,1}=x_i$.  The output is
\[
 h_j(x_{[n]\setminus S_j})=
 \sum_{a\models R_j}\prod_{i\notin S_j}
                       x_i^{a_i}(1-x_i)^{1-a_i}.
\]
Thus $h_j$ is exactly the coefficient of the $j$th placeholder
in the cited first-violation certificate: fixing the falsifying
tuple leaves precisely the assignments satisfying all earlier
clauses.  Multiplication by $v_j$ restores the missing factors
on $S_j$.  In the notation
$\delta_b(x)=\prod_i x_i^{b_i}(1-x_i)^{1-b_i}$, this identifies
$v_jh_j$ with the sum of $\delta_b$ over the first-violation
class of clause $j$.  Equations (4)--(6) of the cited classical
proof therefore give the displayed ordinary integral identity.
No identity merely modulo $x_i^2-x_i$ is substituted for it.

For $n>2k$, each nontrivial call uses $n-k\le n_j\le n$,
so $n_j=\Theta(n)$ and its exponent is bounded by the same
displayed expression in $n$.  There are at most $m$ calls,
and forming all restricted prefixes has polynomial cost
in $n+m$.  The fixed small range of $n$ is handled by direct
expansion.  Each coefficient substitution is a constant-free
linear expression, hence preserves the polynomial formal-degree
bound.  This proves the total circuit and construction bounds.
The displayed expression for $h_j$ proves its actual degree
bound without expanding it.  A uniform algorithm may first
use E Theorem~1.1 to decide satisfiability, reporting that
case or constructing the certificate when the answer is
negative; E's decision cost is absorbed by this weaker exponent.

With placeholders $u_1,\ldots,u_m$, the circuit
\[
                    C(x,u)=\sum_j u_jh_j(x)
\]
is a Hilbert-like IPS certificate: $C(x,0)=0$ and
$C(x,v(x))=1$.  It has actual degree at most $n+1$,
polynomial formal degree, and the same circuit-size bound.
The certificate works over every commutative ring by base
change from its integral identity.  These superpolynomial
size bounds do not make the system polynomially bounded,
give polynomial-time deterministic PIT verification, or
imply a corresponding small monomial-PC derivation.
\end{proof}
\paragraph{Status and complexity impact.}
Integral dense-CNF construction, optimized parameter bound and ordinary ideal-certificate consequence proved. Independent complete deduction-proof review by root: passed; imported input theorems are not claimed independently proved. This closes the specific integral formal-weight gate for arbitrary-density fixed-width CNFs and arbitrary-graph independence polynomials. The cost is a weaker exponent than the characteristic-two compressed construction: it balances disjoint sparsification with the actual larger interpolation degree and does not lift a finite-field circuit. It also constructs ordinary integral Hilbert-like IPS certificates without Boolean axioms and explicit affine coordinate maps at the original-variable bound, for every clause density. The construction is uniform and preserves the original Boolean variables, but remains superpolynomial. It supplies no polynomial-size IPS or arithmetic circuits, no general deterministic polynomial-time PIT, and no VP/VNP or P/NP collapse. Generic verifier reductions still charge their gate variables before this theorem is applied.

\paragraph{Geometric and representation dictionary.}
Let $F$ be a fixed-width CNF on $n$ declared Boolean coordinates. Its integral Boolean scheme is
\[
 Y_F=\operatorname{Spec}\!\left(
 \mathbb Z[x_1,\ldots,x_n]/(x_i^2-x_i,\ 1-c_C(x):C\in F)\right).
\]
Here $c_C=1-\prod_{\ell\in C}(1-\ell(x))$ is the integral clause-satisfaction indicator. The scheme is the disjoint union of one $\operatorname{Spec}\mathbb Z$ for each satisfying assignment. Put $B=\mathbb Z[z_{i,b}]$, $\mathscr A_F=\mathcal O(Y_F)\otimes_{\mathbb Z}B$, and
\[
                         w=\prod_i((1-x_i)z_{i,0}+x_i z_{i,1}).
\]
Then the finite locally free multiplication trace $\operatorname{Tr}_{\mathscr A_F/B}(w)$ equals the integral assignment polynomial $A_F$. For fixed width $k$, this regular trace function has a deterministically constructible integral constant-free circuit of size $(n+m)^{O_k(1)}2^{O_k(n(\log\log n)^{1+1/k}/(\log n)^{1/(2k)})}$ and polynomial formal degree. The same circuit computes the trace after every base change, including all characteristics and nonreduced coefficient rings. Explicit weight factors with polylogarithmic scopes and linear total input degree have the same bound with their description length and fixed parameters included.

Disjoint sparsification is a decomposition of $Y_F$ into clopen finite etale subschemes over $\mathbb Z$, and the trace adds over it. The construction retains the actual integer component multiplicities, rather than only their residues. Equivalently, the torus with one character $z_{i,b}$ per literal rail has an effective module with one basis vector of character $\prod_i z_{i,x_i}$ for each satisfying assignment; its character is computed by the integral circuit using virtual differences. For a graph, this yields the effective independent-subset character after setting the zero rails to one. The exact trace/base-change source is EGA IV4,18.2.1--18.2.2, printed p111, equation18.2.1.2; the split-torus character and Grothendieck-ring locators are those in NN-S019. This is an effective statement for these presented finite schemes and characters, not a theorem about general GCT multiplicities.

For an unsatisfiable normalized width-$k$ CNF with clause
polynomials $v_1,\ldots,v_m$, the closed subscheme
$V(v_1,\ldots,v_m)\subseteq\mathbb A^n_{\mathbb Z}$ is empty,
and its unit ideal has the displayed explicit integral
circuit certificate of size
$(n+m)^{O_k(1)}2^{O_k(n(\log\log n)^{1+1/k}/(\log n)^{1/(2k)})}$.
No Boolean equations occur in that ideal.  The multipliers
have degree at most $n-|S_j|$ and integral constant-free
circuits of polynomial formal degree.

For the affine hypersurface
\[
 X_F=\{t^2-t=\sum_{j=1}^m y_jv_j(x)\}
                          \subseteq\mathbb A^{n+m+1}_{\mathbb Z},
\]
the same certificate gives the explicit isomorphism
$X_F\simeq\mathbb A^{n+m}_{\mathbb Z}$ by the unit-ideal
coordinate formulas already proved in NN-C029:
\[
 z_j=y_j+h_j(x)(2t-t^2),\qquad
 t=\sum_jv_j(x)z_j,\qquad
 y_j=z_j+h_j(x)(t^2-2t).
\]
The first formula is the map from $X_F$ to coordinates
$(x,z)$, and the last two define its inverse.  The equality
$\sum_jv_jh_j=1$ makes $\sum_jv_jz_j=t$ on $X_F$ and
$\sum_jv_jy_j=t^2-t$ on the inverse image; substitution
then gives both identity composites.  These calculations
use no division and hold over every base ring.  The forward
map has degree at most $n+2$ and the inverse has degree at
most $n+2k+2$, since $\deg(\sum_jv_jz_j)\le k+1$.
Sharing the common sum and multiplier circuits gives both
maps the same mildly subexponential original-$n$ circuit
bound and polynomial formal degree.  This is an effective
integral trivialization for the stated empty clause systems,
not a general recognition theorem for affine space.

